\documentclass[10pt]{article}
\usepackage[utf8]{inputenc}
\usepackage[T1]{fontenc}
\usepackage{graphicx}
\usepackage[export]{adjustbox}
\graphicspath{ {./images/} }
\usepackage{hyperref}
\hypersetup{colorlinks=true, linkcolor=blue, filecolor=magenta, urlcolor=cyan,}
\urlstyle{same}
\usepackage{amsmath}
\usepackage{amsfonts}
\usepackage{amssymb}
\usepackage[version=4]{mhchem}
\usepackage{stmaryrd}
\usepackage{bbold}
\usepackage{mathrsfs}
\usepackage{caption}
\usepackage{multirow}

%New command to display footnote whose markers will always be hidden
\let\svthefootnote\thefootnote
\newcommand\blfootnotetext[1]{%
  \let\thefootnote\relax\footnote{#1}%
  \addtocounter{footnote}{-1}%
  \let\thefootnote\svthefootnote%
}
%Overriding the \footnotetext command to hide the marker if its value is `0`
\let\svfootnotetext\footnotetext
\renewcommand\footnotetext[2][?]{%
  \if\relax#1\relax%
    \ifnum\value{footnote}=0\blfootnotetext{#2}\else\svfootnotetext{#2}\fi%
  \else%
    \if?#1\ifnum\value{footnote}=0\blfootnotetext{#2}\else\svfootnotetext{#2}\fi%
    \else\svfootnotetext[#1]{#2}\fi%
  \fi
}
\DeclareUnicodeCharacter{22C5}{\ifmmode\cdot\else{$\cdot$}\fi}
\DeclareUnicodeCharacter{0394}{\ifmmode\Delta\else{$\Delta$}\fi}
\DeclareUnicodeCharacter{22EE}{\ifmmode\vdots\else{$\vdots$}\fi}
\title{Chapter 12: M-Estimators and Nonlinear Estimation}

\author{}
\date{}

\begin{document}
\maketitle

\includegraphics[max width=\textwidth, alt={}]{2e17339c-7635-4ef8-aded-2191d8d853ff-0426_102_1224_116_283}
\end{center}}
In this part we begin our study of nonlinear econometric methods. What we mean by nonlinear needs some explanation because it does not necessarily mean that the underlying model is what we would think of as nonlinear. For example, suppose the population model of interest can be written as $y=\mathbf{x} \boldsymbol{\beta}+u$, but, rather than assuming $\mathrm{E}(u \mid \mathbf{x})=0$, we assume that the median of $u$ given $\mathbf{x}$ is zero for all $\mathbf{x}$. This assumption implies $\operatorname{Med}(y \mid \mathbf{x})=\mathbf{x} \boldsymbol{\beta}$, which is a linear model for the conditional median of $y$ given $\mathbf{x}$. (The conditional mean, $\mathrm{E}(y \mid \mathbf{x})$, may or may not be linear in $\mathbf{x}$.) The standard estimator for a conditional median turns out to be least absolute deviations (LAD), not ordinary least squares. Like OLS, the LAD estimator solves a minimization problem: it minimizes the sum of absolute residuals. However, there is a key difference between LAD and OLS: the LAD estimator cannot be obtained in closed form. The lack of a closed-form expression for LAD has implications not only for obtaining the LAD estimates from a sample of data, but also for the asymptotic theory of LAD.

All the estimators we studied in Part II were obtained in closed form, a feature that greatly facilitates asymptotic analysis: we needed nothing more than the weak law of large numbers, the central limit theorem, and the basic algebra of probability limits. When an estimation method does not deliver closed-form solutions, we need to use more advanced asymptotic theory. In what follows, "nonlinear" describes any problem in which the estimators cannot be obtained in closed form.

The three chapters in this part provide the foundation for asymptotic analysis of most nonlinear models encountered in applications with cross section or panel data. We will make certain assumptions concerning continuity and differentiability, and so problems violating these conditions will not be covered. In the general development of M-estimators in Chapter 12, we will mention some of the applications that are ruled out and provide references.

This part of the book is by far the most technical. We will not dwell on the sometimes intricate arguments used to establish consistency and asymptotic normality in nonlinear contexts. For completeness, we do provide some general results on consistency and asymptotic normality for general classes of estimators. However, for specific estimation methods, such as nonlinear least squares, we will only state assumptions that have real impact for performing inference. Unless the underlying regularity conditions-which involve assuming that certain moments of the population random variables are finite, as well as assuming continuity and differentiability of the regression function or log-likelihood function-are obviously false, they are usually just assumed. Where possible, the assumptions will correspond closely with those given previously for linear models.

The analysis of maximum likelihood methods in Chapter 13 is greatly simplified once we have given a general treatment of M-estimators. Chapter 14 contains results for generalized method of moments estimators for models nonlinear in parameters. We also briefly discuss the related topic of minimum distance estimation in Chapter 14.

Readers who are not interested in general approaches to nonlinear estimation might use the more technical material in these chapters only when needed for reference in Part IV.

\section*{\begin{center}
\includegraphics[max width=\textwidth, alt={}]{2e17339c-7635-4ef8-aded-2191d8d853ff-0428_111_1129_116_283}
\end{center}}
\subsection*{12.1 Introduction}
We begin our study of nonlinear estimation with a general class of estimators known as M-estimators, a term introduced by Huber (1967). (You might think of the "M" as standing for minimization or maximization.) M-estimation methods include maximum likelihood, nonlinear least squares, least absolute deviations, quasi-maximum likelihood, and many other procedures used by econometricians.

Much of this chapter is somewhat abstract and technical, but it is useful to develop a unified theory early on so that it can be applied in a variety of situations. We will carry along the example of nonlinear least squares for cross section data to motivate the general approach. In Sections 12.9 and 12.10, we study multivariate nonlinear regression and quantile regression, two practically important estimation methods.

In a nonlinear regression model, we have a random variable, $y$, and we would like to model $\mathrm{E}(y \mid \mathbf{x})$ as a function of the explanatory variables $\mathbf{x}$, a $K$-vector. We already know how to estimate models of $\mathrm{E}(y \mid \mathbf{x})$ when the model is linear in its parameters: OLS produces consistent, asymptotically normal estimators. What happens if the regression function is nonlinear in its parameters?

Generally, let $m(\mathbf{x}, \boldsymbol{\theta})$ be a parametric model for $\mathrm{E}(y \mid \mathbf{x})$, where $m$ is a known function of $\mathbf{x}$ and $\boldsymbol{\theta}$, and $\boldsymbol{\theta}$ is a $P \times 1$ parameter vector. (This is a parametric model because $m(\cdot, \boldsymbol{\theta})$ is assumed to be known up to a finite number of parameters.) The dimension of the parameters, $P$, can be less than or greater than $K$. The parameter space, $\boldsymbol{\Theta}$, is a subset of $\mathbb{R}^{P}$. This is the set of values of $\boldsymbol{\theta}$ that we are willing to consider in the regression function. Unlike in linear models, for nonlinear models the asymptotic analysis requires explicit assumptions on the parameter space.

An example of a nonlinear regression function is the exponential regression function, $m(\mathbf{x}, \boldsymbol{\theta})=\exp (\mathbf{x} \boldsymbol{\theta})$, where $\mathbf{x}$ is a row vector and contains unity as its first element. This is a useful functional form whenever $y \geq 0$. A regression model suitable when the response $y$ is restricted to the unit interval is the logistic function, $m(\mathbf{x}, \boldsymbol{\theta})= \exp (\mathbf{x} \boldsymbol{\theta}) /[1+\exp (\mathbf{x} \boldsymbol{\theta})]$. Both the exponential and logistic functions are nonlinear in $\boldsymbol{\theta}$.

In any application, there is no guarantee that our chosen model is adequate for $\mathrm{E}(y \mid \mathbf{x})$. We say that we have a correctly specified model for the conditional mean, $\mathrm{E}(y \mid \mathbf{x})$, if, for some $\boldsymbol{\theta}_{\mathrm{o}} \in \boldsymbol{\Theta}$,


\begin{equation*}
\mathrm{E}(y \mid \mathbf{x})=m\left(\mathbf{x}, \boldsymbol{\theta}_{\mathrm{o}}\right) . \tag{12.1}
\end{equation*}


We introduce the subscript "o" on theta to distinguish the parameter vector appearing in $\mathrm{E}(y \mid \mathbf{x})$ from other candidates for that vector. (Often, the value $\boldsymbol{\theta}_{\mathrm{o}}$ is called "the true value of theta," a phrase that is somewhat loose but still useful as\\
shorthand.) As an example, for $y>0$ and a single explanatory variable $x>0$, consider the model $m(x, \boldsymbol{\theta})=\theta_{1} x^{\theta_{2}}$. If the population regression function is $\mathrm{E}(y \mid x)=4 x^{1.5}$, then $\theta_{\mathrm{o} 1}=4$ and $\theta_{\mathrm{o} 2}=1.5$. We will never know the actual $\theta_{\mathrm{o} 1}$ and $\theta_{\mathrm{o} 2}$ (unless we somehow control the way the data have been generated), but, if the model is correctly specified, then these values exist, and we would like to estimate them. Generic candidates for $\theta_{\mathrm{o} 1}$ and $\theta_{\mathrm{o} 2}$ are labeled $\theta_{1}$ and $\theta_{2}$, and, without further information, $\theta_{1}$ is any positive number and $\theta_{2}$ is any real number: the parameter space is $\boldsymbol{\Theta} \equiv\left\{\left(\theta_{1}, \theta_{2}\right): \theta_{1}>0, \theta_{2} \in \mathbb{R}\right\}$. For an exponential regression model, $m(\mathbf{x}, \boldsymbol{\theta})=\exp (\mathbf{x} \boldsymbol{\theta})$ is a correctly specified model for $\mathrm{E}(y \mid \mathbf{x})$ if and only if there is some $K$-vector $\boldsymbol{\theta}_{\mathrm{o}}$ such that $\mathrm{E}(y \mid \mathbf{x})=\exp \left(\mathbf{x} \boldsymbol{\theta}_{\mathrm{o}}\right)$.

In our analysis of linear models, there was no need to make the distinction between the parameter vector in the population regression function and other candidates for this vector because the estimators in linear contexts are obtained in closed form, and so their asymptotic properties can be studied directly. As we will see, in our theoretical development we need to distinguish the vector appearing in $\mathrm{E}(y \mid \mathbf{x})$ from a generic element of $\boldsymbol{\Theta}$. We will often drop the subscripting by "o" when studying particular applications because the notation can be cumbersome when there are many parameters.

Equation (12.1) is the most general way of thinking about what nonlinear least squares is intended to do: estimate models of conditional expectations. But, as a statistical matter, equation (12.1) is equivalent to a model with an additive, unobservable error with a zero conditional mean:


\begin{equation*}
y=m\left(\mathbf{x}, \boldsymbol{\theta}_{\mathrm{o}}\right)+u, \quad \mathrm{E}(u \mid \mathbf{x})=0 \tag{12.2}
\end{equation*}


Given equation (12.2), equation (12.1) clearly holds. Conversely, given equation (12.1), we obtain equation (12.2) by defining the error to be $u \equiv y-m\left(\mathbf{x}, \boldsymbol{\theta}_{\mathrm{o}}\right)$. In interpreting the model and deciding on appropriate estimation methods, we should not focus on the error form in equation (12.2) because, evidently, the additivity of $u$ has some unintended connotations. In particular, we must remember that, in writing the model in additive error form, the only thing implied by equation (12.1) is $\mathrm{E}(u \mid \mathbf{x})=0$. Depending on the nature of $y$, the error $u$ may have some unusual properties. For example, if $y \geq 0$ then $u \geq-m\left(\mathbf{x}, \boldsymbol{\theta}_{\mathrm{o}}\right)$, in which case $u$ and $\mathbf{x}$ cannot be independent. Heteroskedasticity in the error-that is, $\operatorname{Var}(u \mid \mathbf{x}) \neq \operatorname{Var}(u)$-is present whenever $\operatorname{Var}(y \mid \mathbf{x})$ depends on $\mathbf{x}$, as is very common when $y$ takes on a restricted range of values. Plus, when we introduce randomly sampled observations $\left\{\left(\mathbf{x}_{i}, y_{i}\right): i=1,2, \ldots, N\right\}$, it is too tempting to write the model and its assumptions as " $y_{i}=m\left(\mathbf{x}_{i}, \boldsymbol{\theta}_{\mathrm{o}}\right)+u_{i}$ where the $u_{i}$ are i.i.d. errors." As we discussed in Section 1.4 for\\
the linear model, under random sampling the $\left\{u_{i}\right\}$ are always i.i.d. What is usually meant is that $u_{i}$ and $\mathbf{x}_{i}$ are independent, but, for the reasons we just gave, this assumption is often much too strong. The error form of the model does turn out to be useful for defining estimators of asymptotic variances and for obtaining test statistics.

For later reference, we formalize the first nonlinear least squares (NLS) assumption as follows:\\
assumption NLS.1: For some $\boldsymbol{\theta}_{\mathrm{o}} \in \boldsymbol{\Theta}, \mathrm{E}(y \mid \mathbf{x})=m\left(\mathbf{x}, \boldsymbol{\theta}_{\mathrm{o}}\right)$.\\
This form of presentation represents the level at which we will state assumptions for particular econometric methods. In our general development of M-estimators that follows, we will need to add conditions involving moments of $m(\mathbf{x}, \boldsymbol{\theta})$ and $y$, as well as continuity assumptions on $m(\mathbf{x}, \cdot)$.

If we let $\mathbf{w} \equiv(\mathbf{x}, y)$, then $\boldsymbol{\theta}_{\mathrm{o}}$ indexes a feature of the population distribution of $\mathbf{w}$, namely, the conditional mean of $y$ given $\mathbf{x}$. More generally, let $\mathbf{w}$ be an $M$-vector of random variables with some distribution in the population. We let $\mathscr{W}$ denote the subset of $\mathbb{R}^{M}$ representing the possible values of $\mathbf{w}$. Let $\boldsymbol{\theta}_{\mathrm{o}}$ denote a parameter vector describing some feature of the distribution of $\mathbf{w}$. This could be a conditional mean, a conditional mean and conditional variance, a conditional median, or a conditional distribution. As shorthand, we call $\boldsymbol{\theta}_{\mathrm{o}}$ "the true parameter" or "the true value of theta." These phrases simply mean that $\boldsymbol{\theta}_{\mathrm{o}}$ is the parameter vector describing the underlying population, something we will make precise later. We assume that $\boldsymbol{\theta}_{\mathrm{o}}$ belongs to a known parameter space $\boldsymbol{\Theta} \subset \mathbb{R}^{P}$.

We assume that our data come as a random sample of size $N$ from the population; we label this random sample $\left\{\mathbf{w}_{i}: i=1,2, \ldots\right\}$, where each $\mathbf{w}_{i}$ is an $M$-vector. This assumption is much more general than it may initially seem. It covers cross section models with many equations, and it also covers panel data settings with small time series dimension. The extension to independently pooled cross sections is almost immediate. In the NLS example, $\mathbf{w}_{i}$ consists of $\mathbf{x}_{i}$ and $y_{i}$, the $i$ th draw from the population on $\mathbf{x}$ and $y$.

What allows us to estimate $\boldsymbol{\theta}_{\mathrm{o}}$ when it indexes $\mathrm{E}(y \mid \mathbf{x})$ ? It is the fact that $\boldsymbol{\theta}_{\mathrm{o}}$ is the value of $\boldsymbol{\theta}$ that minimizes the expected squared error between $y$ and $m(\mathbf{x}, \boldsymbol{\theta})$. That is, $\boldsymbol{\theta}_{\mathrm{o}}$ solves the population problem


\begin{equation*}
\min _{\boldsymbol{\theta} \in \boldsymbol{\Theta}} \mathrm{E}\left\{[y-m(\mathbf{x}, \boldsymbol{\theta})]^{2}\right\}, \tag{12.3}
\end{equation*}


where the expectation is over the joint distribution of $(\mathbf{x}, y)$. This conclusion follows immediately from basic properties of conditional expectations (in particular, condition CE. 8 in Chapter 2). We will give a slightly different argument here. Write


\begin{align*}
{[y-m(\mathbf{x}, \boldsymbol{\theta})]^{2}=} & {\left[y-m\left(\mathbf{x}, \boldsymbol{\theta}_{\mathrm{o}}\right)\right]^{2}+2\left[m\left(\mathbf{x}, \boldsymbol{\theta}_{\mathrm{o}}\right)-m(\mathbf{x}, \boldsymbol{\theta})\right] u } \\
& +\left[m\left(\mathbf{x}, \boldsymbol{\theta}_{\mathrm{o}}\right)-m(\mathbf{x}, \boldsymbol{\theta})\right]^{2} \tag{12.4}
\end{align*}


where $u$ is defined in equation (12.2). Now, since $\mathrm{E}(u \mid \mathbf{x})=0, u$ is uncorrelated with any function of $\mathbf{x}$, including $m\left(\mathbf{x}, \boldsymbol{\theta}_{\mathrm{o}}\right)-m(\mathbf{x}, \boldsymbol{\theta})$. Thus, taking the expected value of equation (12.4) gives


\begin{equation*}
\mathrm{E}\left\{[y-m(\mathbf{x}, \boldsymbol{\theta})]^{2}\right\}=\mathrm{E}\left\{\left[y-m\left(\mathbf{x}, \boldsymbol{\theta}_{\mathrm{o}}\right)\right]^{2}\right\}+\mathrm{E}\left\{\left[m\left(\mathbf{x}, \boldsymbol{\theta}_{\mathrm{o}}\right)-m(\mathbf{x}, \boldsymbol{\theta})\right]^{2}\right\} . \tag{12.5}
\end{equation*}


Since the last term in equation (12.5) is nonnegative, it follows that


\begin{equation*}
\mathrm{E}\left\{[y-m(\mathbf{x}, \boldsymbol{\theta})]^{2}\right\} \geq \mathrm{E}\left\{\left[y-m\left(\mathbf{x}, \boldsymbol{\theta}_{\mathrm{o}}\right)\right]^{2}\right\}, \quad \text { all } \boldsymbol{\theta} \in \boldsymbol{\Theta} . \tag{12.6}
\end{equation*}


The inequality is strict when $\boldsymbol{\theta} \neq \boldsymbol{\theta}_{\mathrm{o}}$ unless $\mathrm{E}\left\{\left[m\left(\mathbf{x}, \boldsymbol{\theta}_{\mathrm{o}}\right)-m(\mathbf{x}, \boldsymbol{\theta})\right]^{2}\right\}=0$; for $\boldsymbol{\theta}_{\mathrm{o}}$ to be identified, we will have to rule this possibility out.

Because $\boldsymbol{\theta}_{\mathrm{o}}$ solves the population problem in expression (12.3), the analogy principle-which we introduced in Chapter 4-suggests estimating $\boldsymbol{\theta}_{\mathrm{o}}$ by solving the sample analogue. In other words, we replace the population moment $\mathrm{E}\left\{\left[(y-m(\mathbf{x}, \boldsymbol{\theta})]^{2}\right\}\right.$ with the sample average. The NLS estimator of $\boldsymbol{\theta}_{\mathrm{o}}, \hat{\boldsymbol{\theta}}$, solves


\begin{equation*}
\min _{\boldsymbol{\theta} \in \boldsymbol{\Theta}} N^{-1} \sum_{i=1}^{N}\left[y_{i}-m\left(\mathbf{x}_{i}, \boldsymbol{\theta}\right)\right]^{2} . \tag{12.7}
\end{equation*}


For now, we assume that a solution to this problem exists.\\
The NLS objective function in expression (12.7) is a special case of a more general class of estimators. Let $q(\mathbf{w}, \boldsymbol{\theta})$ be a function of the random vector $\mathbf{w}$ and the parameter vector $\boldsymbol{\theta}$. An M-estimator of $\boldsymbol{\theta}_{\mathrm{o}}$ solves the problem


\begin{equation*}
\min _{\boldsymbol{\theta} \in \boldsymbol{\Theta}} N^{-1} \sum_{i=1}^{N} q\left(\mathbf{w}_{i}, \boldsymbol{\theta}\right), \tag{12.8}
\end{equation*}


assuming that a solution, call it $\hat{\boldsymbol{\theta}}$, exists. The estimator clearly depends on the sample $\left\{\mathbf{w}_{i}: i=1,2, \ldots, N\right\}$, but we suppress that fact in the notation.

The objective function for an M-estimator is a sample average of a function of $\mathbf{w}_{i}$ and $\boldsymbol{\theta}$. The division by $N$, while needed for the theoretical development, does not affect the minimization problem. Also, the focus on minimization rather than maximization is without loss of generality because maximization can be trivially turned into minimization.

The parameter vector $\boldsymbol{\theta}_{\mathrm{o}}$ is assumed to uniquely solve the population problem


\begin{equation*}
\min _{\boldsymbol{\theta} \in \boldsymbol{\Theta}} \mathrm{E}[q(\mathbf{w}, \boldsymbol{\theta})] . \tag{12.9}
\end{equation*}


Comparing equations (12.8) and (12.9), we see that M -estimators are based on the analogy principle. Once $\boldsymbol{\theta}_{\mathrm{O}}$ has been defined, finding an appropriate function $q$ that delivers $\boldsymbol{\theta}_{\mathrm{o}}$ as the solution to problem (12.9) requires basic results from probability theory. Often there is more than one choice of $q$ such that $\boldsymbol{\theta}_{\mathrm{o}}$ solves problem (12.9), in which case the choice depends on efficiency or computational issues. For the next several sections, we carry along the NLS example; we treat maximum likelihood estimation in Chapter 13.

How do we translate the fact that $\boldsymbol{\theta}_{\mathrm{o}}$ solves the population problem (12.9) into consistency of the M-estimator $\hat{\boldsymbol{\theta}}$ that solves problem (12.8)? Heuristically, the argument is as follows. Since for each $\boldsymbol{\theta} \in \boldsymbol{\Theta},\left\{q\left(\mathbf{w}_{i}, \boldsymbol{\theta}\right): i=1,2, \ldots\right\}$ is just an i.i.d. sequence, the law of large numbers implies that


\begin{equation*}
N^{-1} \sum_{i=1}^{N} q\left(\mathbf{w}_{i}, \boldsymbol{\theta}\right) \xrightarrow{p} \mathrm{E}[q(\mathbf{w}, \boldsymbol{\theta})] \tag{12.10}
\end{equation*}


under very weak finite moment assumptions. Since $\hat{\boldsymbol{\theta}}$ minimizes the function on the left side of equation (12.10) and $\boldsymbol{\theta}_{\mathrm{o}}$ minimizes the function on the right, it seems plausible that $\hat{\boldsymbol{\theta}} \xrightarrow{p} \boldsymbol{\theta}_{\mathrm{o}}$. This informal argument turns out to be correct, except in pathological cases. There are essentially two issues to address. The first is identifiability of $\boldsymbol{\theta}_{\mathrm{o}}$, which is purely a population issue. The second is the sense in which the convergence in equation (12.10) happens across different values of $\boldsymbol{\theta}$ in $\boldsymbol{\Theta}$.

\subsection*{12.2 Identification, Uniform Convergence, and Consistency}
We now present a formal consistency result for M-estimators under fairly weak assumptions. As mentioned previously, the conditions can be broken down into two parts. The first part is the identification or identifiability of $\boldsymbol{\theta}_{\mathrm{O}}$. For nonlinear regression, we showed how $\boldsymbol{\theta}_{\mathrm{o}}$ solves the population problem (12.3). However, we did not argue that $\boldsymbol{\theta}_{\mathrm{o}}$ is always the unique solution to problem (12.3). Whether or not this is the case depends on the distribution of $\mathbf{x}$ and the nature of the regression function:

ASSUMPTION NLS.2: $\quad \mathrm{E}\left\{\left[m\left(\mathbf{x}, \boldsymbol{\theta}_{\mathrm{o}}\right)-m(\mathbf{x}, \boldsymbol{\theta})\right]^{2}\right\}>0$, all $\boldsymbol{\theta} \in \boldsymbol{\Theta}, \boldsymbol{\theta} \neq \boldsymbol{\theta}_{\mathrm{o}}$.\\
Assumption NLS. 2 plays the same role as Assumption OLS. 2 in Chapter 4. It can fail if the explanatory variables $\mathbf{x}$ do not have sufficient variation in the population. In fact, in the linear case $m(\mathbf{x}, \boldsymbol{\theta})=\mathbf{x} \boldsymbol{\theta}$, Assumption NLS. 2 holds if and only if rank $\mathrm{E}\left(\mathbf{x}^{\prime} \mathbf{x}\right)=K$, which is just Assumption OLS. 2 from Chapter 4. In nonlinear models, Assumption NLS. 2 can fail if $m\left(\mathbf{x}, \boldsymbol{\theta}_{\mathrm{O}}\right)$ depends on fewer parameters than are actually in $\boldsymbol{\theta}$. For example, suppose that we choose as our model $m(\mathbf{x}, \boldsymbol{\theta})=\theta_{1}+\theta_{2} x_{2}+\theta_{3} x_{3}^{\theta_{4}}$,\\
but the true model is linear in $x_{2}: \theta_{\mathrm{o} 3}=0$. Then $\mathrm{E}[(y-m(\mathbf{x}, \boldsymbol{\theta}))]^{2}$ is minimized for any $\boldsymbol{\theta}$ with $\theta_{1}=\theta_{\mathrm{o} 1}, \theta_{2}=\theta_{\mathrm{o} 2}, \theta_{3}=0$, and $\theta_{4}$ any value. If $\theta_{\mathrm{o} 3} \neq 0$, Assumption NLS. 2 would typically hold provided there is sufficient variation in $x_{2}$ and $x_{3}$. Because identification fails for certain values of $\boldsymbol{\theta}_{\mathrm{o}}$, this is an example of a poorly identified model. (See Section 9.5 for other examples of poorly identified models.)

Identification in commonly used nonlinear regression models, such as exponential and logistic regression functions, holds under weak conditions, provided perfect collinearity in $\mathbf{x}$ can be ruled out. For the most part, we will just assume that, when the model is correctly specified, $\boldsymbol{\theta}_{\mathrm{o}}$ is the unique solution to problem (12.3). For the general M-estimation case, we assume that $q(\mathbf{w}, \boldsymbol{\theta})$ has been chosen so that $\boldsymbol{\theta}_{\mathrm{o}}$ is a solution to problem (12.9). Identification requires that $\boldsymbol{\theta}_{\mathrm{O}}$ be the unique solution:


\begin{equation*}
\mathrm{E}\left[q\left(\mathbf{w}, \boldsymbol{\theta}_{\mathrm{o}}\right)\right]<\mathrm{E}[q(\mathbf{w}, \boldsymbol{\theta})], \quad \text { all } \boldsymbol{\theta} \in \boldsymbol{\Theta}, \quad \boldsymbol{\theta} \neq \boldsymbol{\theta}_{\mathrm{o}} . \tag{12.11}
\end{equation*}


The second component for consistency of the M-estimator is convergence of the sample average $N^{-1} \sum_{i=1}^{N} q\left(\mathbf{w}_{i}, \boldsymbol{\theta}\right)$ to its expected value. It turns out that pointwise convergence in probability, as stated in equation (12.10), is not sufficient for consistency. That is, it is not enough to simply invoke the usual weak law of large numbers at each $\boldsymbol{\theta} \in \boldsymbol{\Theta}$. Instead, uniform convergence in probability is sufficient. Mathematically,


\begin{equation*}
\max _{\boldsymbol{\theta} \in \boldsymbol{\Theta}}\left|N^{-1} \sum_{i=1}^{N} q\left(\mathbf{w}_{i}, \boldsymbol{\theta}\right)-\mathrm{E}[q(\mathbf{w}, \boldsymbol{\theta})]\right| \xrightarrow{p} 0 . \tag{12.12}
\end{equation*}


Uniform convergence clearly implies pointwise convergence, but the converse is not true: it is possible for equation (12.10) to hold but equation (12.12) to fail. Nevertheless, under certain regularity conditions, the pointwise convergence in equation (12.10) translates into the uniform convergence in equation (12.12).

To state a formal result concerning uniform convergence, we need to be more careful in stating assumptions about the function $q(\cdot, \cdot)$ and the parameter space $\boldsymbol{\Theta}$. Since we are taking expected values of $q(\mathbf{w}, \boldsymbol{\theta})$ with respect to the distribution of $\mathbf{w}$, $q(\mathbf{w}, \boldsymbol{\theta})$ must be a random variable for each $\boldsymbol{\theta} \in \boldsymbol{\Theta}$. Technically, we should assume that $q(\cdot, \boldsymbol{\theta})$ is a Borel measurable function on $\mathscr{W}$ for each $\boldsymbol{\theta} \in \boldsymbol{\Theta}$. Since it is very difficult to write down a function that is not Borel measurable, we spend no further time on it. Rest assured that any objective function that arises in econometrics is Borel measurable. You are referred to Billingsley (1979) and Davidson (1994, Chap. 3).

The next assumption concerning $q$ is practically more important. We assume that, for each $\mathbf{w} \in \mathscr{W}, q(\mathbf{w}, \cdot)$ is a continuous function over the parameter space $\boldsymbol{\Theta}$. All of the problems we treat in detail have objective functions that are continuous in the\\
parameters, but these do not cover all cases of interest. For example, Manski's (1975) maximum score estimator for binary response models has an objective function that is not continuous in $\boldsymbol{\theta}$. (We cover binary response models in Chapter 15.) It is possible to somewhat relax the continuity assumption in order to handle such cases, but we will not need that generality. See Manski (1988, Sect. 7.3) and Newey and McFadden (1994).

Obtaining uniform convergence is generally difficult for unbounded parameter sets, such as $\boldsymbol{\Theta}=\mathbb{R}^{P}$. It is easiest to assume that $\boldsymbol{\Theta}$ is a compact subset of $\mathbb{R}^{P}$, which means that $\boldsymbol{\Theta}$ is closed and bounded (see Rudin, 1976, Theorem 2.41). Because the natural parameter spaces in most applications are not bounded (and sometimes not closed), the compactness assumption is untidy for developing a general theory of estimation. However, for most applications it is not an assumption to worry about: $\boldsymbol{\Theta}$ can be defined to be such a large closed and bounded set as to always contain $\boldsymbol{\theta}_{\mathrm{o}}$. Some consistency results for nonlinear estimation without compact parameter spaces are available; see the discussion and references in Newey and McFadden (1994).

We can now state a theorem concerning uniform convergence appropriate for the random sampling environment. This result, known as the uniform weak law of large numbers (UWLLN), dates back to LeCam (1953). See also Newey and McFadden (1994, Lemma 2.4).\\
theorem 12.1 (Uniform Weak Law of Large Numbers): Let $\mathbf{w}$ be a random vector taking values in $\mathscr{W} \subset \mathbb{R}^{M}$, let $\boldsymbol{\Theta}$ be a subset of $\mathbb{R}^{P}$, and let $q: \mathscr{W} \times \boldsymbol{\Theta} \rightarrow \mathbb{R}$ be a realvalued function. Assume that (a) $\boldsymbol{\Theta}$ is compact; (b) for each $\boldsymbol{\theta} \in \boldsymbol{\Theta}, q(\cdot, \boldsymbol{\theta})$ is Borel measurable on $\mathscr{W}$; (c) for each $\mathbf{w} \in \mathscr{W}, q(\mathbf{w}, \cdot)$ is continuous on $\boldsymbol{\Theta}$; and (d) $|q(\mathbf{w}, \boldsymbol{\theta})| \leq b(\mathbf{w})$ for all $\boldsymbol{\theta} \in \boldsymbol{\Theta}$, where $b$ is a nonnegative function on $\mathscr{W}$ such that $\mathrm{E}[b(\mathbf{w})]<\infty$. Then equation (12.12) holds.

The only assumption we have not discussed is assumption d , which requires the expected absolute value of $q(\mathbf{w}, \boldsymbol{\theta})$ to be bounded across $\boldsymbol{\theta}$. This kind of moment condition is rarely verified in practice, although it can be. For example, for NLS with $0 \leq y \leq 1$ and the mean function such that $0 \leq m(\mathbf{x}, \boldsymbol{\theta}) \leq 1$ for all $\mathbf{x}$ and $\boldsymbol{\theta}$, we can take $b(\mathbf{w}) \equiv 1$. See Newey and McFadden (1994) for more complicated examples.

The continuity and compactness assumptions are important for establishing uniform convergence, and they also ensure that both the sample minimization problem (12.8) and the population minimization problem (12.9) actually have solutions. Consider problem (12.8) first. Under the assumptions of Theorem 12.1, the sample average is a continuous function of $\boldsymbol{\theta}$, since $q\left(\mathbf{w}_{i}, \boldsymbol{\theta}\right)$ is continuous for each $\mathbf{w}_{i}$. Since a continuous function on a compact space always achieves its minimum, the M-estimation\\
problem is well defined (there could be more than one solution). As a technical matter, it can be shown that $\hat{\boldsymbol{\theta}}$ is actually a random variable under the measurability assumption on $q(\cdot, \boldsymbol{\theta})$. See, for example, Gallant and White (1988).

It can also be shown that, under the assumptions of Theorem 12.1, the function $\mathrm{E}[q(\mathbf{w}, \boldsymbol{\theta})]$ is continuous as a function of $\boldsymbol{\theta}$. Therefore, problem (12.9) also has at least one solution; identifiability ensures that it has only one solution, which in turn implies consistency of the M-estimator.\\
theorem 12.2 (Consistency of M-Estimators): Under the assumptions of Theorem 12.1, assume that the identification assumption (12.11) holds. Then a random vector, $\hat{\boldsymbol{\theta}}$, solves problem (12.8), and $\hat{\boldsymbol{\theta}} \xrightarrow{p} \boldsymbol{\theta}_{\mathrm{o}}$.

A proof of Theorem 12.2 is given in Newey and McFadden (1994). For nonlinear least squares, once Assumptions NLS. 1 and NLS. 2 are maintained, the practical requirement is that $m(\mathbf{x}, \cdot)$ be a continuous function over $\boldsymbol{\Theta}$. Since this assumption is almost always true in applications of NLS, we do not list it as a separate assumption. Noncompactness of $\boldsymbol{\Theta}$ is not much of a concern for most applications.

Theorem 12.2 also applies to median regression. Suppose that the conditional median of $y$ given $\mathbf{x}$ is $\operatorname{Med}(y \mid \mathbf{x})=m\left(\mathbf{x}, \boldsymbol{\theta}_{\mathrm{o}}\right)$, where $m(\mathbf{x}, \boldsymbol{\theta})$ is a known function of $\mathbf{x}$ and $\boldsymbol{\theta}$. The leading case is a linear model, $m(\mathbf{x}, \boldsymbol{\theta})=\mathbf{x} \boldsymbol{\theta}$, where $\mathbf{x}$ contains unity. The least absolute deviations (LAD) estimator of $\boldsymbol{\theta}_{\mathrm{o}}$ solves\\
$\min _{\boldsymbol{\theta} \in \boldsymbol{\Theta}} N^{-1} \sum_{i=1}^{N}\left|y_{i}-m\left(\mathbf{x}_{i}, \boldsymbol{\theta}\right)\right|$.\\
If $\boldsymbol{\Theta}$ is compact and $m(\mathbf{x}, \cdot)$ is continuous over $\boldsymbol{\Theta}$ for each $\mathbf{x}$, a solution always exists. The LAD estimator is motivated by the fact that $\boldsymbol{\theta}_{\mathrm{o}}$ minimizes $\mathrm{E}[|y-m(\mathbf{x}, \boldsymbol{\theta})|]$ over the parameter space $\boldsymbol{\Theta}$; this follows by the fact that for each $\mathbf{x}$, the conditional median is the minimum absolute loss predictor conditional on $\mathbf{x}$. (See, for example, Bassett and Koenker, 1978, and Manski, 1988, Sect. 4.2.2.) If we assume that $\boldsymbol{\theta}_{\mathrm{o}}$ is the unique solution-a standard identification assumption-then the LAD estimator is consistent very generally. In addition to the continuity, compactness, and identification assumptions, it suffices that $\mathrm{E}[|y|]<\infty$ and $|m(\mathbf{x}, \boldsymbol{\theta})| \leq a(\mathbf{x})$ for some function $a(\cdot)$ such that $\mathrm{E}[a(\mathbf{x})]<\infty$. (To see this point, take $b(\mathbf{w}) \equiv|y|+a(\mathbf{x})$ in Theorem 12.2.)

Median regression is a special case of quantile regression, where we model quantiles in the distribution of $y$ given $\mathbf{x}$. For example, in addition to the median, we can estimate how the first and third quartiles in the distribution of $y$ given $\mathbf{x}$ change with $\mathbf{x}$. Except for the median (which leads to LAD), the objective function that identifies a\\
conditional quantile is asymmetric about zero. We study quantile regression in Section 12.10.

We end this section with a lemma that we use repeatedly in the rest of this chapter. It follows from Lemma 4.3 in Newey and McFadden (1994).\\
lemma 12.1: Suppose that $\hat{\boldsymbol{\theta}} \xrightarrow{p} \boldsymbol{\theta}_{\mathrm{o}}$, and assume that $r(\mathbf{w}, \boldsymbol{\theta})$ satisfies the same assumptions on $q(\mathbf{w}, \boldsymbol{\theta})$ in Theorem 12.2. Then


\begin{equation*}
N^{-1} \sum_{i=1}^{N} r\left(\mathbf{w}_{i}, \hat{\boldsymbol{\theta}}\right) \xrightarrow{p} \mathrm{E}\left[r\left(\mathbf{w}, \boldsymbol{\theta}_{\mathrm{o}}\right)\right] \tag{12.13}
\end{equation*}


That is, $N^{-1} \sum_{i=1}^{N} r\left(\mathbf{w}_{i}, \hat{\boldsymbol{\theta}}\right)$ is a consistent estimator of $\mathrm{E}\left[r\left(\mathbf{w}, \boldsymbol{\theta}_{\mathrm{o}}\right)\right]$.\\
Intuitively, Lemma 12.1 is quite reasonable. We know that $N^{-1} \sum_{i=1}^{N} r\left(\mathbf{w}_{i}, \boldsymbol{\theta}_{\mathrm{o}}\right)$ generally converges in probability to $\mathrm{E}\left[r\left(\mathbf{w}, \boldsymbol{\theta}_{\mathrm{o}}\right)\right]$ by the law of large numbers. Lemma 12.1 shows that, if we replace $\boldsymbol{\theta}_{\mathrm{o}}$ in the sample average with a consistent estimator, the convergence still holds, at least under standard regularity conditions. In fact, as shown by Newey and McFadden (1994), we need only assume $r(\mathbf{w}, \boldsymbol{\theta})$ is continuous at $\boldsymbol{\theta}_{\mathrm{o}}$ with probability one.

\subsection*{12.3 Asymptotic Normality}
Under additional assumptions on the objective function, we can also show that Mestimators are asymptotically normally distributed (and converge at the rate $\sqrt{N}$ ). It turns out that continuity over the parameter space does not ensure asymptotic normality.

The simplest asymptotic normality proof proceeds as follows. Assume that $\boldsymbol{\theta}_{\mathrm{o}}$ is in the interior of $\boldsymbol{\Theta}$, which means that $\boldsymbol{\Theta}$ must have nonempty interior; this assumption is true in most applications. Then, since $\hat{\boldsymbol{\theta}} \xrightarrow{p} \boldsymbol{\theta}_{\mathrm{o}}, \hat{\boldsymbol{\theta}}$ is in the interior of $\boldsymbol{\Theta}$ with probability approaching one. If $q(\mathbf{w}, \cdot)$ is continuously differentiable on the interior of $\boldsymbol{\Theta}$, then (with probability approaching one) $\hat{\boldsymbol{\theta}}$ solves the first-order condition\\
$\sum_{i=1}^{N} \mathbf{s}\left(\mathbf{w}_{i}, \hat{\boldsymbol{\theta}}\right)=\mathbf{0}$,\\
where $\mathbf{s}(\mathbf{w}, \boldsymbol{\theta})$ is the $P \times 1$ vector of partial derivatives of $q(\mathbf{w}, \boldsymbol{\theta}): \mathbf{s}(\mathbf{w}, \boldsymbol{\theta})^{\prime}= \nabla_{\theta} q(\mathbf{w}, \boldsymbol{\theta}) \equiv\left[\partial q(\mathbf{w}, \boldsymbol{\theta}) / \partial \theta_{1}, \partial q(\mathbf{w}, \boldsymbol{\theta}) / \partial \theta_{2}, \ldots, \partial q(\mathbf{w}, \boldsymbol{\theta}) / \partial \theta_{P}\right]$. (That is, $\mathbf{s}(\mathbf{w}, \boldsymbol{\theta})$ is the transpose of the gradient of $q(\mathbf{w}, \boldsymbol{\theta})$.) We call $\mathbf{s}(\mathbf{w}, \boldsymbol{\theta})$ the score of the objective function $q(\mathbf{w}, \boldsymbol{\theta})$. While condition (12.14) can only be guaranteed to hold with probability\\
approaching one, usually it holds exactly; at any rate, we will drop the qualifier, as it does not affect the derivation of the limiting distribution.

If $q(\mathbf{w}, \cdot)$ is twice continuously differentiable, then each row of the left-hand side of equation (12.14) can be expanded about $\boldsymbol{\theta}_{\mathrm{O}}$ in a mean-value expansion:


\begin{equation*}
\sum_{i=1}^{N} \mathbf{s}\left(\mathbf{w}_{i}, \hat{\boldsymbol{\theta}}\right)=\sum_{i=1}^{N} \mathbf{s}\left(\mathbf{w}_{i}, \boldsymbol{\theta}_{\mathrm{o}}\right)+\left(\sum_{i=1}^{N} \ddot{\mathbf{H}}_{i}\right)\left(\hat{\boldsymbol{\theta}}-\boldsymbol{\theta}_{\mathrm{o}}\right) . \tag{12.15}
\end{equation*}


The notation $\ddot{\mathbf{H}}_{i}$ denotes the $P \times P$ Hessian of the objective function, $q\left(\mathbf{w}_{i}, \boldsymbol{\theta}\right)$, with respect to $\boldsymbol{\theta}$, but with each row of $\mathbf{H}\left(\mathbf{w}_{i}, \boldsymbol{\theta}\right) \equiv \partial^{2} q\left(\mathbf{w}_{i}, \boldsymbol{\theta}\right) / \partial \boldsymbol{\theta} \partial \boldsymbol{\theta}^{\prime} \equiv \nabla_{\theta}^{2} q\left(\mathbf{w}_{i}, \boldsymbol{\theta}\right)$ evaluated at a different mean value. Each of the $P$ mean values is on the line segment between $\boldsymbol{\theta}_{\mathrm{o}}$ and $\hat{\boldsymbol{\theta}}$. We cannot know what these mean values are, but we do know that each must converge in probability to $\boldsymbol{\theta}_{\mathrm{o}}$ (since each is "trapped" between $\hat{\boldsymbol{\theta}}$ and $\boldsymbol{\theta}_{\mathrm{o}}$ ).

Combining equations (12.14) and (12.15) and multiplying through by $1 / \sqrt{N}$ gives

$$
\mathbf{0}=N^{-1 / 2} \sum_{i=1}^{N} \mathbf{s}\left(\mathbf{w}_{i}, \boldsymbol{\theta}_{\mathrm{o}}\right)+\left(N^{-1} \sum_{i=1}^{N} \ddot{\mathbf{H}}_{i}\right) \sqrt{N}\left(\hat{\boldsymbol{\theta}}-\boldsymbol{\theta}_{\mathrm{o}}\right) .
$$

Now, we can apply Lemma 12.1 to get $N^{-1} \sum_{i=1}^{N} \ddot{\mathbf{H}}_{i} \xrightarrow{p} \mathrm{E}\left[\mathbf{H}\left(\mathbf{w}, \boldsymbol{\theta}_{\mathrm{o}}\right)\right]$ (under some moment conditions). If $\mathbf{A}_{\mathrm{o}} \equiv \mathrm{E}\left[\mathbf{H}\left(\mathbf{w}, \boldsymbol{\theta}_{\mathrm{o}}\right)\right]$ is nonsingular, then $N^{-1} \sum_{i=1}^{N} \ddot{\mathbf{H}}_{i}$ is nonsingular w.p.a. 1 and $\left(N^{-1} \sum_{i=1}^{N} \ddot{\mathbf{H}}_{i}\right)^{-1} \xrightarrow{p} \mathbf{A}_{\mathrm{o}}^{-1}$. Therefore, we can write

$$
\sqrt{N}\left(\hat{\boldsymbol{\theta}}-\boldsymbol{\theta}_{\mathrm{o}}\right)=\left(N^{-1} \sum_{i=1}^{N} \ddot{\mathbf{H}}_{i}\right)^{-1}\left[-N^{-1 / 2} \sum_{i=1}^{N} \mathbf{s}_{i}\left(\boldsymbol{\theta}_{\mathrm{o}}\right)\right]
$$

where $\mathbf{s}_{i}\left(\boldsymbol{\theta}_{\mathrm{o}}\right) \equiv \mathbf{s}\left(\mathbf{w}_{i}, \boldsymbol{\theta}_{\mathrm{o}}\right)$. As we will show, $\mathrm{E}\left[\mathbf{s}_{i}\left(\boldsymbol{\theta}_{\mathrm{o}}\right)\right]=\boldsymbol{0}$. Therefore, $N^{-1 / 2} \sum_{i=1}^{N} \mathbf{s}_{i}\left(\boldsymbol{\theta}_{\mathrm{o}}\right)$ generally satisfies the central limit theorem because it is the average of i.i.d. random vectors with zero mean, multiplied by the usual $\sqrt{N}$. Since $\mathrm{o}_{p}(1) \cdot \mathrm{O}_{p}(1)=\mathrm{o}_{p}(1)$, we have


\begin{equation*}
\sqrt{N}\left(\hat{\boldsymbol{\theta}}-\boldsymbol{\theta}_{\mathrm{o}}\right)=\mathbf{A}_{\mathrm{o}}^{-1}\left[-N^{-1 / 2} \sum_{i=1}^{N} \mathbf{s}_{i}\left(\boldsymbol{\theta}_{\mathrm{o}}\right)\right]+\mathrm{o}_{p}(1) \tag{12.16}
\end{equation*}


This is an important equation. It shows that $\sqrt{N}\left(\hat{\boldsymbol{\theta}}-\boldsymbol{\theta}_{\mathrm{o}}\right)$ inherits its limiting distribution from the average of the scores, evaluated at $\boldsymbol{\theta}_{\mathrm{O}}$. The matrix $\mathbf{A}_{\mathrm{O}}^{-1}$ simply acts as a linear transformation. If we absorb this linear transformation into $\mathbf{s}_{i}\left(\boldsymbol{\theta}_{\mathrm{o}}\right)$, we can write


\begin{equation*}
\sqrt{N}\left(\hat{\boldsymbol{\theta}}-\boldsymbol{\theta}_{\mathrm{o}}\right)=N^{-1 / 2} \sum_{i=1}^{N} \mathbf{e}_{i}\left(\boldsymbol{\theta}_{\mathrm{o}}\right)+\mathrm{o}_{p}(1) \tag{12.17}
\end{equation*}


where $\mathbf{e}_{i}\left(\boldsymbol{\theta}_{\mathrm{o}}\right) \equiv-\mathbf{A}_{\mathrm{o}}^{-1} \mathbf{s}_{i}\left(\boldsymbol{\theta}_{\mathrm{o}}\right)$; this is sometimes called the influence function representation of $\hat{\boldsymbol{\theta}}$, where $\mathbf{e}(\mathbf{w}, \boldsymbol{\theta})$ is the influence function.

Equation (12.16) (or (12.17)) allows us to derive the first-order asymptotic distribution of $\hat{\boldsymbol{\theta}}$. Higher order representations attempt to reduce the error in the $\mathrm{o}_{p}(1)$ term in equation (12.16); such derivations are much more complicated than equation (12.16) and are beyond the scope of this book.

We have essentially proved the following result:\\
theorem 12.3 (Asymptotic Normality of M-Estimators): In addition to the assumptions in Theorem 12.2, assume (a) $\boldsymbol{\theta}_{\mathrm{o}}$ is in the interior of $\boldsymbol{\Theta}$;(b) $\mathbf{s}(\mathbf{w}, \cdot)$ is continuously differentiable on the interior of $\boldsymbol{\Theta}$ for all $\mathbf{w} \in \mathscr{W}$; (c) Each element of $\mathbf{H}(\mathbf{w}, \boldsymbol{\theta})$ is bounded in absolute value by a function $b(\mathbf{w})$, where $\mathrm{E}[b(\mathbf{w})]<\infty$; (d) $\mathbf{A}_{\mathrm{o}} \equiv \mathrm{E}\left[\mathbf{H}\left(\mathbf{w}, \boldsymbol{\theta}_{\mathrm{o}}\right)\right]$ is positive definite; (e) $\mathrm{E}\left[\mathbf{s}\left(\mathbf{w}, \boldsymbol{\theta}_{\mathrm{o}}\right)\right]=\mathbf{0}$; and (f) each element of $\mathbf{s}\left(\mathbf{w}, \boldsymbol{\theta}_{\mathrm{o}}\right)$ has finite second moment.

Then\\
$\sqrt{N}\left(\hat{\boldsymbol{\theta}}-\boldsymbol{\theta}_{\mathrm{o}}\right) \xrightarrow{d} \operatorname{Normal}\left(0, \mathbf{A}_{\mathrm{o}}^{-1} \mathbf{B}_{\mathrm{o}} \mathbf{A}_{\mathrm{o}}^{-1}\right)$,\\
where


\begin{equation*}
\mathbf{A}_{\mathrm{o}} \equiv \mathrm{E}\left[\mathbf{H}\left(\mathbf{w}, \boldsymbol{\theta}_{\mathrm{o}}\right)\right] \tag{12.19}
\end{equation*}


and\\
$\mathbf{B}_{\mathrm{o}} \equiv \mathrm{E}\left[\mathbf{s}\left(\mathbf{w}, \boldsymbol{\theta}_{\mathrm{o}}\right) \mathbf{s}\left(\mathbf{w}, \boldsymbol{\theta}_{\mathrm{o}}\right)^{\prime}\right]=\operatorname{Var}\left[\mathbf{s}\left(\mathbf{w}, \boldsymbol{\theta}_{\mathrm{o}}\right)\right]$.

Thus,\\
$\operatorname{Avar}(\hat{\boldsymbol{\theta}})=\mathbf{A}_{\mathrm{o}}^{-1} \mathbf{B}_{\mathrm{o}} \mathbf{A}_{\mathrm{o}}^{-1} / N$.

Theorem 12.3 implies asymptotic normality of most of the estimators we study in the remainder of the book. A leading example that is not covered by Theorem 12.3 is the LAD estimator. Even if $m(\mathbf{x}, \boldsymbol{\theta})$ is twice continuously differentiable in $\boldsymbol{\theta}$, the objective function for each $i, q\left(\mathbf{w}_{i}, \boldsymbol{\theta}\right) \equiv\left|y_{i}-m\left(\mathbf{x}_{i}, \boldsymbol{\theta}\right)\right|$, is not twice continuously differentiable because the absolute value function is nondifferentiable at zero. By itself, nondifferentiability at zero is a minor nuisance. More important, by any reasonable definition, the Hessian of the LAD objective function is the zero matrix in the leading case of a linear conditional median function, and this feature violates assumption d of Theorem 12.3. It turns out that the LAD estimator is generally $\sqrt{N}$-asymptotically normal, but Theorem 12.3 cannot be applied. We discuss the asymptotic distribution of LAD and other quantile estimators in Section 12.10.

A key component of Theorem 12.3 is that the score evaluated at $\boldsymbol{\theta}_{\mathrm{o}}$ has expected value zero. In many applications, including NLS, we can show this result directly. But it is also useful to know that it holds in the abstract M-estimation framework, at least if we can interchange the expectation and the derivative. To see this point, note that, if $\boldsymbol{\theta}_{\mathrm{o}}$ is in the interior of $\boldsymbol{\Theta}$, and $\mathrm{E}[q(\mathbf{w}, \boldsymbol{\theta})]$ is differentiable for $\boldsymbol{\theta} \in \operatorname{int} \boldsymbol{\Theta}$, then\\
$\left.\nabla_{\theta} \mathrm{E}[q(\mathbf{w}, \boldsymbol{\theta})]\right|_{\boldsymbol{\theta}=\boldsymbol{\theta}_{\mathrm{o}}}=\mathbf{0}$.

Now, if the derivative and expectations operator can be interchanged (which is the case quite generally), then equation (12.22) implies\\
$\mathrm{E}\left[\nabla_{\theta} q\left(\mathbf{w}, \boldsymbol{\theta}_{\mathrm{o}}\right)\right]=\mathrm{E}\left[\mathbf{s}\left(\mathbf{w}, \boldsymbol{\theta}_{\mathrm{o}}\right)\right]=\mathbf{0}$.

A similar argument shows that, in general, $\mathrm{E}\left[\mathbf{H}\left(\mathbf{w}, \boldsymbol{\theta}_{\mathrm{o}}\right)\right]$ is positive semidefinite. If $\boldsymbol{\theta}_{\mathrm{o}}$ is identified, $\mathrm{E}\left[\mathbf{H}\left(\mathbf{w}, \boldsymbol{\theta}_{\mathrm{o}}\right)\right]$ is positive definite.

For the remainder of this chapter, it is convenient to divide the original NLS objective function by two:\\
$q(\mathbf{w}, \boldsymbol{\theta})=[y-m(\mathbf{x}, \boldsymbol{\theta})]^{2} / 2$.

The score of equation (12.24) can be written as


\begin{equation*}
\mathbf{s}(\mathbf{w}, \boldsymbol{\theta})=-\nabla_{\theta} m(\mathbf{x}, \boldsymbol{\theta})^{\prime}[y-m(\mathbf{x}, \boldsymbol{\theta})] \tag{12.25}
\end{equation*}


where $\nabla_{\theta} m(\mathbf{x}, \boldsymbol{\theta})$ is the $1 \times P$ gradient of $m(\mathbf{x}, \boldsymbol{\theta})$, and therefore $\nabla_{\theta} m(\mathbf{x}, \boldsymbol{\theta})^{\prime}$ is $P \times 1$. We can show directly that this expression has an expected value of zero at $\boldsymbol{\theta}=\boldsymbol{\theta}_{\mathrm{o}}$ by showing that expected value of $\mathbf{s}\left(\mathbf{w}, \boldsymbol{\theta}_{\mathrm{o}}\right)$ conditional on $\mathbf{x}$ is zero:\\
$\mathrm{E}\left[\mathbf{s}\left(\mathbf{w}, \boldsymbol{\theta}_{\mathrm{o}}\right) \mid \mathbf{x}\right]=-\nabla_{\theta} m\left(\mathbf{x}, \boldsymbol{\theta}_{\mathrm{o}}\right)^{\prime}\left[\mathrm{E}(y \mid \mathbf{x})-m\left(\mathbf{x}, \boldsymbol{\theta}_{\mathrm{o}}\right)\right]=\mathbf{0}$.

The variance of $\mathbf{s}\left(\mathbf{w}, \boldsymbol{\theta}_{\mathrm{o}}\right)$ is\\
$\mathbf{B}_{\mathrm{o}} \equiv \mathrm{E}\left[\mathbf{s}\left(\mathbf{w}, \boldsymbol{\theta}_{\mathrm{o}}\right) \mathbf{s}\left(\mathbf{w}, \boldsymbol{\theta}_{\mathrm{o}}\right)^{\prime}\right]=\mathrm{E}\left[u^{2} \nabla_{\theta} m\left(\mathbf{x}, \boldsymbol{\theta}_{\mathrm{o}}\right)^{\prime} \nabla_{\theta} m\left(\mathbf{x}, \boldsymbol{\theta}_{\mathrm{o}}\right)\right]$,\\
where the error $u \equiv y-m\left(\mathbf{x}, \boldsymbol{\theta}_{\mathrm{o}}\right)$ is the difference between $y$ and $\mathrm{E}(y \mid \mathbf{x})$.\\
The Hessian of $q(\mathbf{w}, \boldsymbol{\theta})$ is\\
$\mathbf{H}(\mathbf{w}, \boldsymbol{\theta})=\nabla_{\theta} m(\mathbf{x}, \boldsymbol{\theta})^{\prime} \nabla_{\theta} m(\mathbf{x}, \boldsymbol{\theta})-\nabla_{\theta}^{2} m(\mathbf{x}, \boldsymbol{\theta})[y-m(\mathbf{x}, \boldsymbol{\theta})]$,\\
where $\nabla_{\theta}^{2} m(\mathbf{x}, \boldsymbol{\theta})$ is the $P \times P$ Hessian of $m(\mathbf{x}, \boldsymbol{\theta})$ with respect to $\boldsymbol{\theta}$. To find the expected value of $\mathbf{H}(\mathbf{w}, \boldsymbol{\theta})$ at $\boldsymbol{\theta}=\boldsymbol{\theta}_{\mathrm{o}}$, we first find the expectation conditional on $\mathbf{x}$. When evaluated at $\boldsymbol{\theta}_{\mathrm{o}}$, the second term in equation (12.28) is $\nabla_{\theta}^{2} m\left(\mathbf{x}, \boldsymbol{\theta}_{\mathrm{o}}\right) u$, and it therefore has a zero mean conditional on $\mathbf{x}$ (since $\mathrm{E}(u \mid \mathbf{x})=0$ ). Therefore,\\
$\mathrm{E}\left[\mathbf{H}\left(\mathbf{w}, \boldsymbol{\theta}_{\mathrm{o}}\right) \mid \mathbf{x}\right]=\nabla_{\theta} m\left(\mathbf{x}, \boldsymbol{\theta}_{\mathrm{o}}\right)^{\prime} \nabla_{\theta} m\left(\mathbf{x}, \boldsymbol{\theta}_{\mathrm{o}}\right)$.

Taking the expected value of equation (12.29) over the distribution of $\mathbf{x}$ gives\\
$\mathbf{A}_{\mathrm{o}}=\mathrm{E}\left[\nabla_{\theta} m\left(\mathbf{x}, \boldsymbol{\theta}_{\mathrm{o}}\right)^{\prime} \nabla_{\theta} m\left(\mathbf{x}, \boldsymbol{\theta}_{\mathrm{o}}\right)\right]$.

This matrix plays a fundamental role in nonlinear regression. When $\boldsymbol{\theta}_{\mathrm{o}}$ is identified, $\mathbf{A}_{\mathrm{o}}$ is generally positive definite. In the linear case $m(\mathbf{x}, \boldsymbol{\theta})=\mathbf{x} \boldsymbol{\theta}, \mathbf{A}_{\mathrm{o}}=\mathrm{E}\left(\mathbf{x}^{\prime} \mathbf{x}\right)$. In the exponential case $m(\mathbf{x}, \boldsymbol{\theta})=\exp (\mathbf{x} \boldsymbol{\theta}), \mathbf{A}_{\mathrm{o}}=\mathrm{E}\left[\exp \left(2 \mathbf{x} \boldsymbol{\theta}_{\mathrm{o}}\right) \mathbf{x}^{\prime} \mathbf{x}\right]$, which is generally positive definite whenever $\mathrm{E}\left(\mathbf{x}^{\prime} \mathbf{x}\right)$ is. In the example $m(\mathbf{x}, \boldsymbol{\theta})=\theta_{1}+\theta_{2} x_{2}+\theta_{3} x_{3}^{\theta_{4}}$ with $\theta_{\mathrm{o} 3}=0$, it is easy to show that matrix (12.30) has rank less than four.

For nonlinear regression, $\mathbf{A}_{\mathrm{o}}$ and $\mathbf{B}_{\mathrm{o}}$ appear to be similar in that they both depend on $\nabla_{\theta} m\left(\mathbf{x}, \boldsymbol{\theta}_{\mathrm{o}}\right)^{\prime} \nabla_{\theta} m\left(\mathbf{x}, \boldsymbol{\theta}_{\mathrm{o}}\right)$. Generally, though, there is no simple relationship between $\mathbf{A}_{\mathrm{o}}$ and $\mathbf{B}_{\mathrm{o}}$ because the latter depends on the distribution of $u^{2}$, the squared population error. In Section 12.5 we will show that a homoskedasticity assumption implies that $\mathbf{B}_{\mathrm{o}}$ is proportional to $\mathbf{A}_{\mathrm{o}}$.

The previous analysis of nonlinear regression assumes that the conditional mean function is correctly specified. If we drop this assumption, then we cannot simplify the expected Hessian as in equation $(12.30)$. White $(1981,1994)$ studies the properties of NLS when the mean function is misspecified. Under weak conditions, the NLS estimator $\hat{\boldsymbol{\theta}}$ converges in probability to a vector $\boldsymbol{\theta}^{*}$, where $\boldsymbol{\theta}^{*}$ provides the best mean square approximation to the actual regression function, $\mathrm{E}\left(y_{i} \mid \mathbf{x}_{i}=\mathbf{x}\right)$. More precisely, $\mathrm{E}\left\{\left[\mathrm{E}\left(y_{i} \mid \mathbf{x}_{i}\right)-m\left(\mathbf{x}_{i}, \boldsymbol{\theta}^{*}\right)\right]^{2}\right\} \leq \mathrm{E}\left\{\left[\mathrm{E}\left(y_{i} \mid \mathbf{x}_{i}\right)-m\left(\mathbf{x}_{i}, \boldsymbol{\theta}\right)\right]^{2}\right\}$ for $\boldsymbol{\theta} \in \boldsymbol{\Theta}$. The asymptotic normality result in equation (12.18) still holds (with the obvious notational changes). But the second term in equation (12.28), evaluated at $\boldsymbol{\theta}^{*}$, is no longer guaranteed to be zero: $\nabla_{\theta}^{2} m\left(\mathbf{x}_{i}, \boldsymbol{\theta}^{*}\right)$ is generally correlated with $y_{i}-m\left(\mathbf{x}_{i}, \boldsymbol{\theta}^{*}\right)$, because $\mathrm{E}\left(y_{i} \mid \mathbf{x}_{i}\right) \neq m\left(\mathbf{x}_{i}, \boldsymbol{\theta}^{*}\right)$. We have focused on the case of correctly specified conditional mean functions, but one should be aware that the assumption has implications for estimating the asymptotic variance, a topic we return to in Section 12.5.

\subsection*{12.4 Two-Step M-Estimators}
Sometimes applications of M-estimators involve a first-stage estimation (an example is OLS with generated regressors, as in Chapter 6). Let $\hat{\gamma}$ be a preliminary estimator, usually based on the random sample $\left\{\mathbf{w}_{i}: i=1,2, \ldots, N\right\}$. Where this estimator comes from must be vague for the moment.

A two-step M-estimator $\hat{\boldsymbol{\theta}}$ of $\boldsymbol{\theta}_{\mathrm{o}}$ solves the problem


\begin{equation*}
\min _{\boldsymbol{\theta} \in \boldsymbol{\Theta}} \sum_{i=1}^{N} q\left(\mathbf{w}_{i}, \boldsymbol{\theta} ; \hat{\boldsymbol{\gamma}}\right) \tag{12.31}
\end{equation*}


where $q$ is now defined on $\mathscr{W} \times \boldsymbol{\Theta} \times \boldsymbol{\Gamma}$, and $\boldsymbol{\Gamma}$ is a subset of $\mathbb{R}^{J}$. We will see several examples of two-step M-estimators in Chapter 13 and the applications in Part IV. An example of a two-step M-estimator is the weighted nonlinear least squares (WNLS) estimator, where the weights are estimated in a first stage. The WNLS estimator solves


\begin{equation*}
\min _{\boldsymbol{\theta} \in \boldsymbol{\Theta}}(1 / 2) \sum_{i=1}^{N}\left[y_{i}-m\left(\mathbf{x}_{i}, \boldsymbol{\theta}\right)\right]^{2} / h\left(\mathbf{x}_{i}, \hat{\boldsymbol{\gamma}}\right), \tag{12.32}
\end{equation*}


where the weighting function, $h(\mathbf{x}, \boldsymbol{\gamma})$, depends on the explanatory variables and a parameter vector. As with NLS, $m(\mathbf{x}, \boldsymbol{\theta})$ is a model of $\mathrm{E}(y \mid \mathbf{x})$. The function $h(\mathbf{x}, \boldsymbol{\gamma})$ is chosen to be a model of $\operatorname{Var}(y \mid \mathbf{x})$. The estimator $\hat{\gamma}$ comes from a problem used to estimate the conditional variance. We list the key assumptions needed for WNLS to have desirable properties here, but several of the derivations are left for the problems.\\
assumption WNLS.1: Same as Assumption NLS.1.

\subsection*{12.4.1 Consistency}
For the general two-step M-estimator, when will $\hat{\boldsymbol{\theta}}$ be consistent for $\boldsymbol{\theta}_{\mathrm{o}}$ ? In practice, the important condition is the identification assumption. To state the identification condition, we need to know about the asymptotic behavior of $\hat{\gamma}$. A general assumption is that $\hat{\gamma} \xrightarrow{p} \boldsymbol{\gamma}^{*}$, where $\boldsymbol{\gamma}^{*}$ is some element in $\boldsymbol{\Gamma}$. We label this value $\boldsymbol{\gamma}^{*}$ to allow for the possibility that $\hat{\gamma}$ does not converge to a parameter indexing some interesting feature of the distribution of $\mathbf{w}$. In some cases, the plim of $\hat{\gamma}$ will be of direct interest. In the weighted regression case, if we assume that $h(\mathbf{x}, \gamma)$ is a correctly specified model for $\operatorname{Var}(y \mid \mathbf{x})$, then it is possible to choose an estimator such that $\hat{\gamma} \xrightarrow{p} \gamma_{\mathrm{o}}$, where $\operatorname{Var}(y \mid \mathbf{x})=h\left(\mathbf{x}, \boldsymbol{\gamma}_{\mathrm{o}}\right)$. (For an example, see Problem 12.2.) If the variance model is misspecified, plim $\hat{\gamma}$ is generally well defined, but $\operatorname{Var}(y \mid \mathbf{x}) \neq h\left(\mathbf{x}, \gamma^{*}\right)$; it is for this reason that we use the notation $\gamma^{*}$.

The identification condition for the two-step M-estimator is

$$
\mathrm{E}\left[q\left(\mathbf{w}, \boldsymbol{\theta}_{\mathrm{o}} ; \boldsymbol{\gamma}^{*}\right)\right]<\mathrm{E}\left[q\left(\mathbf{w}, \boldsymbol{\theta} ; \boldsymbol{\gamma}^{*}\right)\right], \quad \text { all } \boldsymbol{\theta} \in \boldsymbol{\Theta}, \quad \boldsymbol{\theta} \neq \boldsymbol{\theta}_{\mathrm{o}} .
$$

The consistency argument is essentially the same as that underlying Theorem 12.2. If $q\left(\mathbf{w}_{i}, \boldsymbol{\theta} ; \gamma\right)$ satisfies the UWLLN over $\boldsymbol{\Theta} \times \boldsymbol{\Gamma}$ then expression (12.31) can be shown to converge to $\mathrm{E}\left[q\left(\mathbf{w}, \boldsymbol{\theta} ; \boldsymbol{\gamma}^{*}\right)\right]$ uniformly over $\boldsymbol{\Theta}$. Along with identification, this result can be shown to imply consistency of $\hat{\boldsymbol{\theta}}$ for $\boldsymbol{\theta}_{\mathrm{o}}$.

In some applications of two-step M-estimation, identification of $\boldsymbol{\theta}_{\mathrm{o}}$ holds for any $\boldsymbol{\gamma} \in \boldsymbol{\Gamma}$. This result can be shown for the WNLS estimator (see Problem 12.4). It is for this reason that WNLS is still consistent even if the function $h(\mathbf{x}, \gamma)$ is not correctly\\
specified for $\operatorname{Var}(y \mid \mathbf{x})$. The weakest version of the identification assumption for WNLS is the following:

ASSUMPTION WNLS.2: $\quad \mathrm{E}\left\{\left[m\left(\mathbf{x}, \boldsymbol{\theta}_{\mathrm{o}}\right)-m(\mathbf{x}, \boldsymbol{\theta})\right]^{2} / h\left(\mathbf{x}, \boldsymbol{\gamma}^{*}\right)\right\}>0$, all $\boldsymbol{\theta} \in \boldsymbol{\Theta}, \quad \boldsymbol{\theta} \neq \boldsymbol{\theta}_{\mathrm{o}}$, where $\boldsymbol{\gamma}^{*}=\operatorname{plim} \hat{\boldsymbol{\gamma}}$.

As with the case of NLS, we know that weak inequality holds in Assumption WNLS. 2 under Assumption WNLS.1. The strict inequality in Assumption WNLS. 2 puts restrictions on the distribution of $\mathbf{x}$ and the functional forms of $m$ and $h$.

In other cases, including several two-step maximum likelihood estimators we encounter in Part IV, the identification condition for $\boldsymbol{\theta}_{\mathrm{o}}$ holds only for $\boldsymbol{\gamma}=\boldsymbol{\gamma}^{*}=\boldsymbol{\gamma}_{\mathrm{o}}$, where $\gamma_{\mathrm{o}}$ also indexes some feature of the distribution of $\mathbf{w}$.

\subsection*{12.4.2 Asymptotic Normality}
With the two-step M-estimator, there are two cases worth distinguishing. The first occurs when the asymptotic variance of $\sqrt{N}\left(\hat{\boldsymbol{\theta}}-\boldsymbol{\theta}_{\mathrm{o}}\right)$ does not depend on the asymptotic variance of $\sqrt{N}\left(\hat{\gamma}-\gamma^{*}\right)$, and the second occurs when the asymptotic variance of $\sqrt{N}\left(\hat{\boldsymbol{\theta}}-\boldsymbol{\theta}_{\mathrm{o}}\right)$ should be adjusted to account for the first-stage estimation of $\boldsymbol{\gamma}^{*}$. We first derive conditions under which we can ignore the first-stage estimation error.

Using arguments similar to those in Section 12.3, it can be shown that, under standard regularity conditions,


\begin{equation*}
\sqrt{N}\left(\hat{\boldsymbol{\theta}}-\boldsymbol{\theta}_{\mathrm{o}}\right)=\mathbf{A}_{\mathrm{o}}^{-1}\left(-N^{-1 / 2} \sum_{i=1}^{N} \mathbf{s}_{i}\left(\boldsymbol{\theta}_{\mathrm{o}} ; \hat{\boldsymbol{\gamma}}\right)\right)+\mathrm{o}_{p}(1), \tag{12.33}
\end{equation*}


where now $\mathbf{A}_{\mathrm{o}}=\mathrm{E}\left[\mathbf{H}\left(\mathbf{w}, \boldsymbol{\theta}_{\mathrm{o}} ; \boldsymbol{\gamma}^{*}\right)\right]$. In obtaining the score and the Hessian, we take derivatives only with respect to $\boldsymbol{\theta} ; \boldsymbol{\gamma}^{*}$ simply appears as an extra argument. Now, if


\begin{equation*}
N^{-1 / 2} \sum_{i=1}^{N} \mathbf{s}_{i}\left(\boldsymbol{\theta}_{\mathrm{o}} ; \hat{\boldsymbol{\gamma}}\right)=N^{-1 / 2} \sum_{i=1}^{N} \mathbf{s}_{i}\left(\boldsymbol{\theta}_{\mathrm{o}} ; \boldsymbol{\gamma}^{*}\right)+\mathrm{o}_{p}(1), \tag{12.34}
\end{equation*}


then $\sqrt{N}\left(\hat{\boldsymbol{\theta}}-\boldsymbol{\theta}_{\mathrm{o}}\right)$ behaves the same asymptotically whether we used $\hat{\boldsymbol{\gamma}}$ or its plim in defining the M-estimator.

When does equation (12.34) hold? Assuming that $\sqrt{N}\left(\hat{\gamma}-\gamma^{*}\right)=\mathrm{O}_{p}(1)$, which is standard, a mean value expansion similar to the one in Section 12.3 gives


\begin{equation*}
N^{-1 / 2} \sum_{i=1}^{N} \mathbf{s}_{i}\left(\boldsymbol{\theta}_{\mathrm{o}} ; \hat{\boldsymbol{\gamma}}\right)=N^{-1 / 2} \sum_{i=1}^{N} \mathbf{s}_{i}\left(\boldsymbol{\theta}_{\mathrm{o}} ; \boldsymbol{\gamma}^{*}\right)+\mathbf{F}_{\mathrm{o}} \sqrt{N}\left(\hat{\boldsymbol{\gamma}}-\boldsymbol{\gamma}^{*}\right)+\mathrm{o}_{p}(1), \tag{12.35}
\end{equation*}


where $\mathbf{F}_{\mathrm{o}}$ is the $P \times J$ matrix\\
$\mathbf{F}_{\mathrm{o}} \equiv \mathrm{E}\left[\nabla_{\gamma} \mathbf{s}\left(\mathbf{w}, \boldsymbol{\theta}_{\mathrm{o}} ; \gamma^{*}\right)\right]$.\\
(Remember, $J$ is the dimension of $\gamma$.) Therefore, if\\
$\mathrm{E}\left[\nabla_{\gamma} \mathbf{s}\left(\mathbf{w}, \boldsymbol{\theta}_{\mathrm{o}} ; \boldsymbol{\gamma}^{*}\right)\right]=\mathbf{0}$,\\
then equation (12.34) holds, and the asymptotic variance of the two-step M-estimator is the same as if $\boldsymbol{\gamma}^{*}$ were plugged in. In other words, under assumption (12.37), we conclude that equation (12.18) holds, where $\mathbf{A}_{\mathrm{o}}$ and $\mathbf{B}_{\mathrm{o}}$ are given in expressions (12.19) and (12.20), respectively, except that $\gamma^{*}$ appears as an argument in the score and Hessian. For deriving the asymptotic distribution of $\sqrt{N}\left(\hat{\boldsymbol{\theta}}-\boldsymbol{\theta}_{\mathrm{o}}\right)$, we can ignore the fact that $\hat{\gamma}$ was obtained in a first-stage estimation.

One case where assumption (12.37) holds is weighted nonlinear least squares, something you are asked to show in Problem 12.4. Naturally, we must assume that the conditional mean is correctly specified, but, interestingly, assumption (12.37) holds whether or not the conditional variance is correctly specified.

There are many problems for which assumption (12.37) does not hold, including some of the methods for correcting for endogeneity in probit and Tobit models in Part IV. In Chapter 21 we will see that two-step methods for correcting sample selection bias are two-step M-estimators, but assumption (12.37) fails. In such cases we need to make an adjustment to the asymptotic variance of $\sqrt{N}\left(\hat{\boldsymbol{\theta}}-\boldsymbol{\theta}_{\mathrm{o}}\right)$. The adjustment is easily obtained from equation (12.35), once we have a first-order representation for $\sqrt{N}\left(\hat{\gamma}-\gamma^{*}\right)$. We assume that


\begin{equation*}
\sqrt{N}\left(\hat{\gamma}-\gamma^{*}\right)=N^{-1 / 2} \sum_{i=1}^{N} \mathbf{r}_{i}\left(\gamma^{*}\right)+\mathrm{o}_{p}(1) \tag{12.38}
\end{equation*}


where $\mathbf{r}_{i}\left(\gamma^{*}\right)$ is a $J \times 1$ vector with $\mathrm{E}\left[\mathbf{r}_{i}\left(\gamma^{*}\right)\right]=\mathbf{0}$ (in practice, $\mathbf{r}_{i}$ depends on parameters other than $\gamma^{*}$, but we suppress those here for simplicity). Therefore, $\hat{\gamma}$ could itself be an M-estimator or, as we will see in Chapter 14, a generalized method of moments estimator. In fact, every estimator considered in this book has a representation as in equation (12.38).

Now we can write


\begin{equation*}
\sqrt{N}\left(\hat{\boldsymbol{\theta}}-\boldsymbol{\theta}_{\mathrm{o}}\right)=\mathbf{A}_{\mathrm{o}}^{-1} N^{-1 / 2} \sum_{i=1}^{N}\left[-\mathbf{g}_{i}\left(\boldsymbol{\theta}_{\mathrm{o}} ; \boldsymbol{\gamma}^{*}\right)\right]+\mathrm{o}_{p}(1), \tag{12.39}
\end{equation*}


where $\mathbf{g}_{i}\left(\boldsymbol{\theta}_{\mathrm{o}} ; \boldsymbol{\gamma}^{*}\right) \equiv \mathbf{s}_{i}\left(\boldsymbol{\theta}_{\mathrm{o}} ; \boldsymbol{\gamma}^{*}\right)+\mathbf{F}_{\mathrm{o}} \mathbf{r}_{i}\left(\boldsymbol{\gamma}^{*}\right)$. Since $\mathbf{g}_{i}\left(\boldsymbol{\theta}_{\mathrm{o}} ; \boldsymbol{\gamma}^{*}\right)$ has zero mean, the standardized partial sum in equation (12.39) can be assumed to satisfy the central limit theorem. Define the $P \times P$ matrix\\
$\mathbf{D}_{\mathrm{o}} \equiv \mathrm{E}\left[\mathbf{g}_{i}\left(\boldsymbol{\theta}_{\mathrm{o}} ; \boldsymbol{\gamma}^{*}\right) \mathbf{g}_{i}\left(\boldsymbol{\theta}_{\mathrm{o}} ; \boldsymbol{\gamma}^{*}\right)^{\prime}\right]=\operatorname{Var}\left[\mathbf{g}_{i}\left(\boldsymbol{\theta}_{\mathrm{o}} ; \boldsymbol{\gamma}^{*}\right)\right]$.

Then\\
Avar $\sqrt{N}\left(\hat{\boldsymbol{\theta}}-\boldsymbol{\theta}_{\mathrm{o}}\right)=\mathbf{A}_{\mathrm{o}}^{-1} \mathbf{D}_{\mathrm{o}} \mathbf{A}_{\mathrm{o}}^{-1}$.

We will discuss estimation of this matrix in the next section.\\
Sometimes it is informative to compare the correct asymptotic variance in (12.41) with the asymptotic variance one would obtain by ignoring the sampling error in $\hat{\gamma}$; that is, by using the incorrect formula in equation (12.18). (If $\mathbf{F}_{\mathrm{o}}=\mathbf{0}$, both formulas are the same.) Most often, the concern is that ignoring estimation of $\gamma^{*}$ leads one to be too optimistic about the precision in $\hat{\boldsymbol{\theta}}$, which happens when $\mathbf{D}_{\mathrm{o}}-\mathbf{B}_{\mathrm{o}}$ is positive semidefinite (p.s.d.). One case where $\mathbf{D}_{\mathrm{o}}-\mathbf{B}_{\mathrm{o}}$ is unambiguously p.s.d. (but not zero) is when $\mathbf{F}_{\mathrm{o}} \neq \mathbf{0}$ and the scores from the first- and second-step estimation are uncorrelated: $\mathrm{E}\left[\mathbf{s}_{i}\left(\boldsymbol{\theta}_{\mathrm{o}} ; \boldsymbol{\gamma}^{*}\right) \mathbf{r}_{i}\left(\boldsymbol{\gamma}^{*}\right)^{\prime}\right]=\mathbf{0}$. Then,\\
$\mathbf{D}_{\mathrm{o}}=\mathrm{E}\left[\mathbf{s}_{i}\left(\boldsymbol{\theta}_{\mathrm{o}} ; \boldsymbol{\gamma}^{*}\right) \mathrm{E}\left[\mathbf{s}_{i}\left(\boldsymbol{\theta}_{\mathrm{o}} ; \boldsymbol{\gamma}^{*}\right)^{\prime}\right]+\mathbf{F}_{\mathrm{o}} \mathrm{E}\left[\mathbf{r}_{i}\left(\boldsymbol{\gamma}^{*}\right) \mathbf{r}_{i}\left(\boldsymbol{\gamma}^{*}\right)^{\prime}\right] \mathbf{F}_{\mathrm{o}}^{\prime}=\mathbf{B}_{\mathrm{o}}+\mathbf{F}_{\mathrm{o}} \mathrm{E}\left[\mathbf{r}_{i}\left(\boldsymbol{\gamma}^{*}\right) \mathbf{r}_{i}\left(\boldsymbol{\gamma}^{*}\right)^{\prime}\right] \mathbf{F}_{\mathrm{o}}^{\prime}\right.$,\\
and the last term is p.s.d. In other cases, $\mathbf{D}_{\mathrm{o}}-\mathbf{B}_{\mathrm{o}}$ is indefinite, in which case it is not true that the correct asymptotic variances are all larger than the incorrect ones. Perhaps surprisingly, it is also possible that $\mathbf{D}_{\mathrm{o}}-\mathbf{B}_{\mathrm{o}}$ is negative semidefinite. An immediate implication is that it is possible that estimating $\gamma^{*}$, rather than knowing $\gamma^{*}$, can actually lead to a more precise estimator of $\boldsymbol{\theta}_{\mathrm{o}}$. Situations where estimation in a first stage reduces the asymptotic variance in a second stage are special, and when it happens, the estimator of $\gamma^{*}$ is usually from a correctly specified maximum likelihood procedure (in which case we would use the notation $\gamma_{\mathrm{o}}$ rather than $\gamma^{*}$ ). Therefore, we postpone further discussion until the next chapter, where we explicitly cover maximum likelihood estimation.

\subsection*{12.5 Estimating the Asymptotic Variance}
\subsection*{12.5.1 Estimation without Nuisance Parameters}
We first consider estimating the asymptotic variance of $\hat{\boldsymbol{\theta}}$ in the case where there are no nuisance parameters. This task requires consistently estimating the matrices $\mathbf{A}_{\mathrm{o}}$ and $\mathbf{B}_{\mathrm{o}}$. One thought is to solve for the expected values of $\mathbf{H}\left(\mathbf{w}, \boldsymbol{\theta}_{\mathrm{o}}\right)$ and $\mathbf{s}\left(\mathbf{w}, \boldsymbol{\theta}_{\mathrm{o}}\right)$. $\mathbf{s}\left(\mathbf{w}, \boldsymbol{\theta}_{\mathrm{o}}\right)^{\prime}$ over the distribution of $\mathbf{w}$, and then to plug in $\hat{\boldsymbol{\theta}}$ for $\boldsymbol{\theta}_{\mathrm{o}}$. When we have completely specified the distribution of $\mathbf{w}$, obtaining closed-form expressions for $\mathbf{A}_{\mathrm{o}}$ and $\mathbf{B}_{\mathrm{o}}$ is, in principle, possible. However, except in simple cases, it would be difficult. More important, we rarely specify the entire distribution of $\mathbf{w}$. Even in a maximum\\
likelihood setting, $\mathbf{w}$ is almost always partitioned into two parts: a set of endogenous variables, $\mathbf{y}$, and conditioning variables, $\mathbf{x}$. Rarely do we wish to specify the distribution of $\mathbf{x}$, and so the expected values needed to obtain $\mathbf{A}_{\mathrm{o}}$ and $\mathbf{B}_{\mathrm{o}}$ are not available.

We can always estimate $\mathbf{A}_{\mathrm{o}}$ consistently by taking away the expectation and replacing $\boldsymbol{\theta}_{\mathrm{o}}$ with $\hat{\boldsymbol{\theta}}$. Under regularity conditions that ensure uniform converge of the Hessian, the estimator


\begin{equation*}
N^{-1} \sum_{i=1}^{N} \mathbf{H}\left(\mathbf{w}_{i}, \hat{\boldsymbol{\theta}}\right) \equiv N^{-1} \sum_{i=1}^{N} \hat{\mathbf{H}}_{i} \tag{12.42}
\end{equation*}


is consistent for $\mathbf{A}_{\mathrm{o}}$, by Lemma 12.1. The advantage of the estimator (12.42) is that it is always available in problems with a twice continuously differentiable objective function. This includes cases where $\boldsymbol{\theta}_{\mathrm{o}}$ does not index a feature of an unconditional or conditional distribution (such as nonlinear regression with the conditional mean function misspecified). Its main drawback is computational: it requires calculation of second derivatives, which is a nontrivial task for certain estimation problems.

Because we assume $\hat{\boldsymbol{\theta}}$ solves the minimization problem in equation (12.8), and $\hat{\boldsymbol{\theta}}$ is in the interior of $\boldsymbol{\Theta}$, we know that the Hessian evaluated at $\hat{\boldsymbol{\theta}}$ is at least p.s.d. If the model is well specified in the sense that $\boldsymbol{\theta}_{\mathrm{o}}$ is identified, then the estimator in (12.42) will actually be positive definite with high probability because $\mathrm{E}\left[\mathbf{H}\left(\mathbf{w}_{i}, \boldsymbol{\theta}_{\mathrm{o}}\right)\right]$ is positive definite. In some cases, the objective function is strictly convex on $\boldsymbol{\Theta}$, in which case $\sum_{i=1}^{N} \mathbf{H}\left(\mathbf{w}_{i}, \boldsymbol{\theta}\right)$ is positive definite for all $\boldsymbol{\theta} \in \boldsymbol{\Theta}$.

In many econometric applications, more structure is available that allows a different estimator. Suppose we can partition $\mathbf{w}$ into $\mathbf{x}$ and $\mathbf{y}$, and that $\boldsymbol{\theta}_{\mathrm{o}}$ indexes some feature of the distribution of $\mathbf{y}$ given $\mathbf{x}$ (such as the conditional mean or, in the case of maximum likelihood, the conditional distribution). Define


\begin{equation*}
\mathbf{A}\left(\mathbf{x}, \boldsymbol{\theta}_{\mathrm{o}}\right) \equiv \mathrm{E}\left[\mathbf{H}\left(\mathbf{w}, \boldsymbol{\theta}_{\mathrm{o}}\right) \mid \mathbf{x}\right] . \tag{12.43}
\end{equation*}


While $\mathbf{H}\left(\mathbf{w}, \boldsymbol{\theta}_{\mathrm{o}}\right)$ is generally a function of $\mathbf{x}$ and $\mathbf{y}, \mathbf{A}\left(\mathbf{x}, \boldsymbol{\theta}_{\mathrm{o}}\right)$ is a function only of $\mathbf{x}$. By the law of iterated expectations, $\mathrm{E}\left[\mathbf{A}\left(\mathbf{x}, \boldsymbol{\theta}_{\mathrm{o}}\right)\right]=\mathrm{E}\left[\mathbf{H}\left(\mathbf{w}, \boldsymbol{\theta}_{\mathrm{o}}\right)\right]=\mathbf{A}_{\mathrm{o}}$. From Lemma 12.1 and standard regularity conditions it follows that


\begin{equation*}
N^{-1} \sum_{i=1}^{N} \mathbf{A}\left(\mathbf{x}_{i}, \hat{\boldsymbol{\theta}}\right) \equiv N^{-1} \sum_{i=1}^{N} \hat{\mathbf{A}}_{i} \xrightarrow{p} \mathbf{A}_{\circ} . \tag{12.44}
\end{equation*}


The estimator (12.44) of $\mathbf{A}_{\mathrm{o}}$ is useful in cases where $\mathrm{E}\left[\mathbf{H}\left(\mathbf{w}, \boldsymbol{\theta}_{\mathrm{o}}\right) \mid \mathbf{x}\right]$ can be obtained in closed form or is easily approximated. In some leading cases, including NLS and certain maximum likelihood problems, $\mathbf{A}\left(\mathbf{x}, \boldsymbol{\theta}_{\mathrm{o}}\right)$ depends only on the first derivatives of the conditional mean function.

When the estimator (12.44) is available, it is usually the case that $\boldsymbol{\theta}_{\mathrm{O}}$ actually minimizes $\mathrm{E}[q(\mathbf{w}, \boldsymbol{\theta}) \mid \mathbf{x}]$ for any value of $\mathbf{x}$; from equation (12.4) this is easily seen to be the case for NLS with a correctly specified conditional mean function. Under assumptions that allow the interchange of derivative and expectation, this result implies that $\mathbf{A}\left(\mathbf{x}, \boldsymbol{\theta}_{\mathrm{o}}\right)$ is p.s.d. The expected value of $\mathbf{A}\left(\mathbf{x}, \boldsymbol{\theta}_{\mathrm{o}}\right)$ over the distribution of $\mathbf{x}$ is positive definite provided $\boldsymbol{\theta}_{\mathrm{o}}$ is identified. Therefore, the estimator (12.44) is usually positive definite in the sample.

Obtaining a p.s.d. estimator of $\mathbf{B}_{\mathrm{o}}$ is straightforward. By Lemma 12.1, under standard regularity conditions we have


\begin{equation*}
N^{-1} \sum_{i=1}^{N} \mathbf{s}\left(\mathbf{w}_{i}, \hat{\boldsymbol{\theta}}\right) \mathbf{s}\left(\mathbf{w}_{i}, \hat{\boldsymbol{\theta}}\right)^{\prime} \equiv N^{-1} \sum_{i=1}^{N} \hat{\mathbf{s}}_{i} \hat{\mathbf{s}}_{i}^{\prime} \xrightarrow{p} \mathbf{B}_{\mathrm{o}} . \tag{12.45}
\end{equation*}


Combining the estimator (12.45) with the consistent estimators for $\mathbf{A}_{\mathrm{o}}$, we can consistently estimate Avar $\sqrt{N}\left(\hat{\boldsymbol{\theta}}-\boldsymbol{\theta}_{\mathrm{o}}\right)$ by


\begin{equation*}
\widehat{\operatorname{Avar} \sqrt{N}}\left(\hat{\boldsymbol{\theta}}-\boldsymbol{\theta}_{\mathrm{o}}\right)=\hat{\mathbf{A}}^{-1} \hat{\mathbf{B}} \hat{\mathbf{A}}^{-1} \tag{12.46}
\end{equation*}


where $\hat{\mathbf{A}}$ is one of the estimators (12.42) or (12.44). The asymptotic standard errors are obtained from the matrix


\begin{equation*}
\hat{\mathbf{V}} \equiv \widehat{\operatorname{var}}(\hat{\boldsymbol{\theta}})=\hat{\mathbf{A}}^{-1} \hat{\mathbf{B}} \hat{\mathbf{A}}^{-1} / N, \tag{12.47}
\end{equation*}


which can be expressed as


\begin{equation*}
\left(\sum_{i=1}^{N} \hat{\mathbf{H}}_{i}\right)^{-1}\left(\sum_{i=1}^{N} \hat{\mathbf{s}}_{i} \hat{\mathbf{s}}_{i}^{\prime}\right)\left(\sum_{i=1}^{N} \hat{\mathbf{H}}_{i}\right)^{-1} \tag{12.48}
\end{equation*}


or


\begin{equation*}
\left(\sum_{i=1}^{N} \hat{\mathbf{A}}_{i}\right)^{-1}\left(\sum_{i=1}^{N} \hat{\mathbf{s}}_{i} \hat{\mathbf{s}}_{i}^{\prime}\right)\left(\sum_{i=1}^{N} \hat{\mathbf{A}}_{i}\right)^{-1} \tag{12.49}
\end{equation*}


depending on the estimator used for $\mathbf{A}_{\mathrm{o}}$. Expressions (12.48) and (12.49) are both at least p.s.d. when they are well defined.

The variance matrix estimators in (12.48) and (12.49) are examples of a HuberWhite sandwich estimator, after Huber (1967) and White (1982a). When they differ (due to $\mathbf{H}\left(\mathbf{w}, \boldsymbol{\theta}_{\mathrm{o}}\right) \neq \mathbf{A}\left(\mathbf{x}, \boldsymbol{\theta}_{\mathrm{o}}\right)$ ), the estimator in equation (12.49) is usually valid only when some feature of the conditional distribution $\mathrm{D}\left(\mathbf{y}_{i} \mid \mathbf{x}_{i}\right)$ is correctly specified (because otherwise the calculation of $\mathrm{E}\left[\mathbf{H}\left(\mathbf{w}, \boldsymbol{\theta}_{\mathrm{o}}\right) \mid \mathbf{x}\right]$ would be incorrect). Sometimes it is\\
useful to make a distinction by calling the estimator in (12.49) a semirobust variance matrix estimator. The variance matrix estimator that requires the fewest assumptions to produce valid inference is in (12.48) (and so, for emphasis, it might be called a fully robust variance matrix estimator).

In the case of NLS, the estimator of $\mathbf{A}_{\mathrm{o}}$ in equation (12.44) is always available when $\mathrm{E}\left(y_{i} \mid \mathbf{x}_{i}\right)=m\left(\mathbf{x}_{i}, \boldsymbol{\theta}_{\mathrm{o}}\right)$, and it is usually used if we think the mean function is correctly specified:\\
$\sum_{i=1}^{N} \hat{\mathbf{A}}_{i}=\sum_{i=1}^{N} \nabla_{\theta} \hat{m}_{i}^{\prime} \nabla_{\theta} \hat{m}_{i}$,\\
where $\nabla_{\theta} \hat{m}_{i} \equiv \nabla_{\theta} m\left(\mathbf{x}_{i}, \hat{\boldsymbol{\theta}}\right)$ for every observation $i$. Also, the estimated score for NLS can be written as\\
$\hat{\mathbf{s}}_{i}=-\nabla_{\theta} \hat{m}_{i}^{\prime}\left[y_{i}-m\left(\mathbf{x}_{i}, \hat{\boldsymbol{\theta}}\right)\right]=-\nabla_{\theta} \hat{m}_{i}^{\prime} \hat{u}_{i}$,\\
where the nonlinear least squares residuals, $\hat{u}_{i}$, are defined as\\
$\hat{u}_{i} \equiv y_{i}-m\left(\mathbf{x}_{i}, \hat{\boldsymbol{\theta}}\right)$.

The estimated asymptotic variance of the NLS estimator is\\
$\widehat{\operatorname{Avar}(\hat{\boldsymbol{\theta}})}=\left(\sum_{i=1}^{N} \nabla_{\theta} \hat{m}_{i}^{\prime} \nabla_{\theta} \hat{m}_{i}\right)^{-1}\left(\sum_{i=1}^{N} \hat{u}_{i}^{2} \nabla_{\theta} \hat{m}_{i}^{\prime} \nabla_{\theta} \hat{m}_{i}\right)\left(\sum_{i=1}^{N} \nabla_{\theta} \hat{m}_{i}^{\prime} \nabla_{\theta} \hat{m}_{i}\right)^{-1}$.

This is called the heteroskedasticity-robust variance matrix estimator for NLS because it places no restrictions on $\operatorname{Var}(y \mid \mathbf{x})$. It was first proposed by White (1980a). (Sometimes the expression is multiplied by $N /(N-P)$ as a degrees-of-freedom adjustment, where $P$ is the dimension of $\boldsymbol{\theta}$.) As always, the asymptotic standard error of each element of $\hat{\boldsymbol{\theta}}$ is the square root of the appropriate diagonal element of matrix (12.52).

As a specific example, suppose that $m(\mathbf{x}, \boldsymbol{\theta})=\exp (\mathbf{x} \boldsymbol{\theta})$. Then $\nabla_{\theta} \hat{m}_{i}^{\prime} \nabla_{\theta} \hat{m}_{i}= \exp \left(2 \mathbf{x}_{i} \hat{\boldsymbol{\theta}}\right) \mathbf{x}_{i}^{\prime} \mathbf{x}_{i}$, which has dimension $K \times K$. We can plug this equation into expression (12.52) along with $\hat{u}_{i}=y_{i}-\exp \left(\mathbf{x}_{i} \hat{\boldsymbol{\theta}}\right)$.

Using the previous terminology, the estimator in (12.52) is a semirobust variance matrix estimator because it assumes that the conditional mean is correctly specified. The sense in which (12.52) is robust is that it is valid without any assumptions on $\operatorname{Var}\left(y_{i} \mid \mathbf{x}_{i}\right)$, but it is not valid if $\mathrm{E}\left(y_{i} \mid \mathbf{x}_{i}\right)$ is misspecified. If we allow misspecification of $m(\mathbf{x}, \boldsymbol{\theta})$, as in White (1981), the formula (12.52) is no longer valid. Instead, the summands in the outer terms of (12.52) should be replaced with $\nabla_{\theta} m\left(\mathbf{x}_{i}, \hat{\boldsymbol{\theta}}\right) \nabla_{\theta} m\left(\mathbf{x}_{i}, \hat{\boldsymbol{\theta}}\right)^{\prime}-\nabla_{\theta}^{2} m\left(\mathbf{x}_{i}, \hat{\boldsymbol{\theta}}\right)\left[y_{i}-m\left(\mathbf{x}_{i}, \hat{\boldsymbol{\theta}}\right)\right]$, which is simply the $P \times P$ Hessian\\
of $\left[y_{i}-m\left(\mathbf{x}_{i}, \boldsymbol{\theta}\right)\right]^{2} / 2$ evaluated at $\hat{\boldsymbol{\theta}}$. In most applications of NLS-and in standard software packages-the conditional mean is assumed to be correctly specified, and equation (12.52) is used as the robust variance matrix estimator.

In many contexts, including NLS and certain quasi-likelihood methods, the asymptotic variance estimator can be simplified under additional assumptions. For our purposes, we state the assumption as follows: For some $\sigma_{\mathrm{o}}^{2}>0$,\\
$\mathrm{E}\left[\mathbf{s}\left(\mathbf{w}, \boldsymbol{\theta}_{\mathrm{o}}\right) \mathbf{s}\left(\mathbf{w}, \boldsymbol{\theta}_{\mathrm{o}}\right)^{\prime}\right]=\sigma_{\mathrm{o}}^{2} \mathrm{E}\left[\mathbf{H}\left(\mathbf{w}, \boldsymbol{\theta}_{\mathrm{o}}\right)\right]$.

This assumption simply says that the expected outer product of the score, evaluated at $\boldsymbol{\theta}_{\mathrm{O}}$, is proportional to the expected value of the Hessian (evaluated at $\boldsymbol{\theta}_{\mathrm{O}}$ ): $\mathbf{B}_{\mathrm{o}}= \sigma_{\mathrm{o}}^{2} \mathbf{A}_{\mathrm{o}}$. Shortly we will provide an assumption under which assumption (12.53) holds for NLS. In the next chapter we will show that assumption (12.53) holds for $\sigma_{\mathrm{o}}^{2}=1$ in the context of maximum likelihood with a correctly specified conditional density. For reasons we will see in Chapter 13, we refer to assumption (12.53) as the generalized information matrix equality (GIME).

LEMMA 12.2: Under regularity conditions of the type contained in Theorem 12.3 and assumption (12.53), $\operatorname{Avar}(\hat{\boldsymbol{\theta}})=\sigma_{\mathrm{o}}^{2} \mathbf{A}_{\mathrm{o}}^{-1} / N$. Therefore, under assumption (12.53), the asymptotic variance of $\hat{\boldsymbol{\theta}}$ can be estimated as


\begin{equation*}
\hat{\mathbf{V}}=\hat{\sigma}^{2}\left(\sum_{i=1}^{N} \hat{\mathbf{H}}_{i}\right)^{-1} \tag{12.54}
\end{equation*}


or\\
$\hat{\mathbf{V}}=\hat{\sigma}^{2}\left(\sum_{i=1}^{N} \hat{\mathbf{A}}_{i}\right)^{-1}$,\\
where $\hat{\mathbf{H}}_{i}$ and $\hat{\mathbf{A}}_{i}$ are defined as before, and $\hat{\sigma}^{2} \xrightarrow{p} \sigma_{\mathrm{o}}^{2}$.\\
In the case of nonlinear regression, the parameter $\sigma_{\mathrm{o}}^{2}$ is the variance of $y$ given $\mathbf{x}$, or equivalently $\operatorname{Var}(u \mid \mathbf{x})$, under homoskedasticity:\\
assumption NLS.3: $\operatorname{Var}(y \mid \mathbf{x})=\operatorname{Var}(u \mid \mathbf{x})=\sigma_{\mathrm{o}}^{2}$.\\
Under Assumption NLS.3, we can show that assumption (12.53) holds with $\sigma_{\mathrm{o}}^{2}= \operatorname{Var}(y \mid \mathbf{x})$. First, since $\mathbf{s}\left(\mathbf{w}, \boldsymbol{\theta}_{\mathrm{o}}\right) \mathbf{s}\left(\mathbf{w}, \boldsymbol{\theta}_{\mathrm{o}}\right)^{\prime}=u^{2} \nabla_{\theta} m\left(\mathbf{x}, \boldsymbol{\theta}_{\mathrm{o}}\right)^{\prime} \nabla_{\theta} m\left(\mathbf{x}, \boldsymbol{\theta}_{\mathrm{o}}\right)$, it follows that


\begin{align*}
\mathrm{E}\left[\mathbf{s}\left(\mathbf{w}, \boldsymbol{\theta}_{\mathrm{o}}\right) \mathbf{s}\left(\mathbf{w}, \boldsymbol{\theta}_{\mathrm{o}}\right)^{\prime} \mid \mathbf{x}\right] & =\mathrm{E}\left(u^{2} \mid \mathbf{x}\right) \nabla_{\theta} m\left(\mathbf{x}, \boldsymbol{\theta}_{\mathrm{o}}\right)^{\prime} \nabla_{\theta} m\left(\mathbf{x}, \boldsymbol{\theta}_{\mathrm{o}}\right) \\
& =\sigma_{\mathrm{o}}^{2} \nabla_{\theta} m\left(\mathbf{x}, \boldsymbol{\theta}_{\mathrm{o}}\right)^{\prime} \nabla_{\theta} m\left(\mathbf{x}, \boldsymbol{\theta}_{\mathrm{o}}\right) \tag{12.56}
\end{align*}


under Assumptions NLS. 1 and NLS.3. Taking the expected value with respect to $\mathbf{x}$ gives equation (12.53).

Under Assumption NLS.3, a simplified estimator of the asymptotic variance of the NLS estimator exists from equation (12.55). Let


\begin{equation*}
\hat{\sigma}^{2}=\frac{1}{(N-P)} \sum_{i=1}^{N} \hat{u}_{i}^{2}=\mathrm{SSR} /(N-P), \tag{12.57}
\end{equation*}


where the $\hat{u}_{i}$ are the NLS residuals (12.51) and SSR is the sum of squared NLS residuals. Using Lemma 12.1, $\hat{\sigma}^{2}$ can be shown to be consistent very generally. The subtraction of $P$ in the denominator of equation (12.57) is an adjustment that is thought to improve the small sample properties of $\hat{\sigma}^{2}$.

Under Assumptions NLS.1-NLS.3, the asymptotic variance of the NLS estimator is estimated as


\begin{equation*}
\hat{\sigma}^{2}\left(\sum_{i=1}^{N} \nabla_{\theta} \hat{m}_{i}^{\prime} \nabla_{\theta} \hat{m}_{i}\right)^{-1} \tag{12.58}
\end{equation*}


This is the default asymptotic variance estimator for NLS, but it is valid only under homoskedasticity; the estimator (12.52) is valid with or without Assumption NLS.3. For an exponential regression function, expression (12.58) becomes $\hat{\sigma}^{2}\left(\sum_{i=1}^{N} \exp \left(2 \mathbf{x}_{i} \hat{\boldsymbol{\theta}}\right) \mathbf{x}_{i}^{\prime} \mathbf{x}_{i}\right)^{-1}$. (Remember, if we want to allow Assumption NLS. 1 to fail, we should use expression [12.48].)

\subsection*{12.5.2 Adjustments for Two-Step Estimation}
In the case of the two-step M-estimator, we may or may not need to adjust the asymptotic variance. If assumption (12.37) holds, estimation is very simple. The most general estimators are expressions (12.48) and (12.49), where $\hat{\mathbf{s}}_{i}, \hat{\mathbf{H}}_{i}$, and $\hat{\mathbf{A}}_{i}$ depend on $\hat{\gamma}$, but we only compute derivatives with respect to $\boldsymbol{\theta}$.

In some cases under assumption (12.37), the analogue of assumption (12.53) holds (with $\gamma_{\mathrm{o}}=\operatorname{plim} \hat{\gamma}$ appearing in $\mathbf{H}$ and $\mathbf{s}$ ). If so, the simpler estimators (12.54) and (12.55) are available. In Problem 12.4 you are asked to show this result for weighted NLS when $\operatorname{Var}(y \mid \mathbf{x})=\sigma_{\mathrm{o}}^{2} h\left(\mathbf{x}, \gamma_{\mathrm{o}}\right)$ and $\gamma_{\mathrm{o}}=\operatorname{plim} \hat{\gamma}$. The natural third assumption for WNLS is that the variance function is correctly specified:\\
assumption WNLS.3: For some $\gamma_{\mathrm{o}} \in \boldsymbol{\Gamma}$ and $\sigma_{\mathrm{o}}^{2}>0, \operatorname{Var}(y \mid \mathbf{x})=\sigma_{\mathrm{o}}^{2} h\left(\mathbf{x}, \gamma_{\mathrm{o}}\right)$. Further, $\sqrt{N}\left(\hat{\gamma}-\gamma_{\mathrm{o}}\right)=\mathrm{O}_{p}(1)$.

Under Assumptions WNLS.1-WNLS.3, the asymptotic variance of the WNLS estimator is estimated as\\
$\hat{\sigma}^{2}\left(\sum_{i=1}^{N}\left(\nabla_{\theta} \hat{m}_{i}^{\prime} \nabla_{\theta} \hat{m}_{i}\right) / \hat{h}_{i}\right)^{-1}$,\\
where $\hat{h}_{i}=h\left(\mathbf{x}_{i}, \hat{\gamma}\right)$ and $\hat{\sigma}^{2}$ is as in equation (12.57) except that the residual $\hat{u}_{i}$ is replaced with the standardized residual, $\hat{u}_{i} / \sqrt{\hat{h}_{i}}$. The sum in expression (12.59) is simply the outer product of the weighted gradients, $\nabla_{\theta} \hat{m}_{i} / \sqrt{\hat{h}_{i}}$. Thus the NLS formulas can be used but with all quantities weighted by $1 / \sqrt{\hat{h}_{i}}$. It is important to remember that expression (12.59) is not valid without Assumption WNLS.3. Without Assumption WNLS.3, but with correct specification of the conditional mean, the asymptotic variance estimator of the WNLS estimator is


\begin{equation*}
\left(\sum_{i=1}^{N} \nabla_{\theta} \hat{m}_{i}^{\prime} \nabla_{\theta} \hat{m}_{i} / \hat{h}_{i}\right)^{-1}\left(\sum_{i=1}^{N} \hat{u}_{i}^{2} \nabla_{\theta} \hat{m}_{i}^{\prime} \nabla_{\theta} \hat{m}_{i} / \hat{h}_{i}^{2}\right)\left(\sum_{i=1}^{N} \nabla_{\theta} \hat{m}_{i}^{\prime} \nabla_{\theta} \hat{m}_{i} / \hat{h}_{i}\right)^{-1} \tag{12.60}
\end{equation*}


where $\hat{u}_{i}=y_{i}-m\left(\mathbf{x}_{i}, \hat{\boldsymbol{\theta}}\right)$ are the WNLS residuals. Notice that this estimator has the same structure as equation (12.52), but where the residuals are replaced with the standardized residuals, $\hat{u}_{i} / \sqrt{\hat{h}_{i}}$, and the gradients are replaced with the weighted gradients, $\nabla_{\theta} \hat{m}_{i} / \sqrt{\hat{h}_{i}}$.

When assumption (12.37) is violated, the asymptotic variance estimator of $\hat{\boldsymbol{\theta}}$ must account for the asymptotic variance of $\hat{\gamma}$; we must estimate equation (12.41). We already know how to consistently estimate $\mathbf{A}_{0}$ : use expression (12.42) or (12.44) where $\hat{\gamma}$ is also plugged in. Estimation of $\mathbf{D}_{\mathrm{o}}$ is also straightforward. First, we need to estimate $\mathbf{F}_{\mathrm{o}}$. An estimator that is always available is the $P \times J$ matrix


\begin{equation*}
\hat{\mathbf{F}}=N^{-1} \sum_{i=1}^{N} \nabla_{\gamma} \mathbf{s}_{i}(\hat{\boldsymbol{\theta}} ; \hat{\gamma}) \tag{12.61}
\end{equation*}


In cases with conditioning variables, such as NLS, a simpler estimator can be obtained by computing $\mathrm{E}\left[\nabla_{\gamma} \mathbf{s}\left(\mathbf{w}_{i}, \boldsymbol{\theta}_{\mathrm{o}}, \boldsymbol{\gamma}^{*}\right) \mid \mathbf{x}_{i}\right]$, replacing $\left(\boldsymbol{\theta}_{\mathrm{o}}, \boldsymbol{\gamma}^{*}\right)$ with $(\hat{\boldsymbol{\theta}}, \hat{\boldsymbol{\gamma}})$, and using this in place of $\nabla_{\gamma} \mathbf{s}_{i}(\hat{\boldsymbol{\theta}} ; \hat{\gamma})$. Next, replace $\mathbf{r}_{i}\left(\gamma^{*}\right)$ with $\hat{\mathbf{r}}_{i} \equiv \mathbf{r}_{i}(\hat{\gamma})$. Then


\begin{equation*}
\hat{\mathbf{D}} \equiv N^{-1} \sum_{i=1}^{N} \hat{\mathbf{g}}_{i} \hat{\mathbf{g}}_{i}^{\prime} \tag{12.62}
\end{equation*}


is consistent for $\mathbf{D}_{\mathrm{o}}$, where $\hat{\mathbf{g}}_{i}=\hat{\mathbf{s}}_{i}+\hat{\mathbf{F}}_{i}$. The asymptotic variance of the two-step Mestimator can be obtained as in expression (12.48) or (12.49), but where $\hat{\mathbf{s}}_{i}$ is replaced with $\hat{\mathbf{g}}_{i}$.

In some cases, a subtle issue arises in deciding whether adjustment for the presence of $\hat{\gamma}$ is needed. For example, in WNLS with a correctly specified conditional mean (and without any assumption about the conditional variance), $\mathbf{F}_{\mathrm{o}}=\mathbf{0}$, and so no adjustment to the asymptotic variance of $\hat{\boldsymbol{\theta}}$ is needed. But if the mean is misspecified so that $\operatorname{plim}(\hat{\boldsymbol{\theta}})=\boldsymbol{\theta}^{*}$, where $\boldsymbol{\theta}^{*}$ minimizes the weighted mean squared error, then the corresponding matrix, say $\mathbf{F}^{*} \equiv \mathrm{E}\left[\nabla_{\gamma} \mathbf{s}_{i}\left(\boldsymbol{\theta}^{*} ; \boldsymbol{\gamma}^{*}\right)\right]$, is not zero. Then, neither (12.59) nor (12.60) is valid. Further, it is not enough to replace the expected Hessian (under correct specification of the conditional mean) with the Hessian in (12.60). A version that allows the mean to be misspecified must account for the asymptotic variance of $\sqrt{N}\left(\hat{\gamma}-\gamma^{*}\right)$ (which would itself require recognition that the conditional mean is misspecified in the initial NLS estimation). For general two-step M-estimation, we typically require some feature of $\mathrm{D}\left(\mathbf{y}_{i} \mid \mathbf{x}_{i}\right)$ to be correctly specified (such as the mean in the case of WNLS) in order to justify ignoring the sampling variation in $\hat{\gamma}$. In practice, calculating the variance for two-step methods that allows misspecification of the feature of $\mathrm{D}\left(\mathbf{y}_{i} \mid \mathbf{x}_{i}\right)$ of interest is rarely done.

\subsection*{12.6 Hypothesis Testing}
\subsection*{12.6.1 Wald Tests}
Wald tests are easily obtained once we choose a form of the asymptotic variance. To test the $Q$ restrictions $\mathrm{H}_{0}: \mathbf{c}\left(\boldsymbol{\theta}_{\mathrm{o}}\right)=\mathbf{0}$, we can form the Wald statistic


\begin{equation*}
W \equiv \mathbf{c}(\hat{\boldsymbol{\theta}})^{\prime}\left(\hat{\mathbf{C}} \hat{\mathbf{V}} \hat{\mathbf{C}}^{\prime}\right)^{-1} \mathbf{c}(\hat{\boldsymbol{\theta}}), \tag{12.63}
\end{equation*}


where $\hat{\mathbf{V}}$ is an asymptotic variance matrix estimator of $\hat{\boldsymbol{\theta}}, \hat{\mathbf{C}} \equiv \mathbf{C}(\hat{\boldsymbol{\theta}})$, and $\mathbf{C}(\boldsymbol{\theta})$ is the $Q \times P$ Jacobian of $\mathbf{c}(\boldsymbol{\theta})$. The estimator $\hat{\mathbf{V}}$ can be chosen to be fully robust, as in expression (12.48) or (12.49); under assumption (12.53), the simpler forms in Lemma 12.2 are available. Also, $\hat{\mathbf{V}}$ can be chosen to account for two-step estimation, when necessary. Provided $\hat{\mathbf{V}}$ has been chosen appropriately, $W \stackrel{a}{\sim} \chi_{Q}^{2}$ under $\mathrm{H}_{0}$.

A couple of practical restrictions are needed for $W$ to have a limiting $\chi_{Q}^{2}$ distribution. First, $\boldsymbol{\theta}_{\mathrm{o}}$ must be in the interior of $\boldsymbol{\Theta}$; that is, $\boldsymbol{\theta}_{\mathrm{o}}$ cannot be on the boundary. If, for example, the first element of $\boldsymbol{\theta}$ must be nonnegative-and we impose this restriction in the estimation-then expression (12.63) does not have a limiting chi-square distribution under $\mathrm{H}_{0}: \theta_{\mathrm{o} 1}=0$. The second condition is that $\mathbf{C}\left(\boldsymbol{\theta}_{\mathrm{o}}\right)=\nabla_{\theta} \mathbf{c}\left(\boldsymbol{\theta}_{\mathrm{o}}\right)$ must have rank $Q$. This rules out cases where $\boldsymbol{\theta}_{\mathrm{o}}$ is unidentified under the null hypothesis, such as the NLS example where $m(\mathbf{x}, \boldsymbol{\theta})=\theta_{1}+\theta_{2} x_{2}+\theta_{3} x_{3}^{\theta_{4}}$ and $\theta_{\mathrm{o} 3}=0$ under $\mathrm{H}_{0}$.

One drawback to the Wald statistic is that it is not invariant to how the nonlinear restrictions are imposed. We can change the outcome of a hypothesis test by rede-\\
fining the constraint function, $\mathbf{c}(\cdot)$. The easiest way to illustrate the lack of invariance of the Wald statistic is to use an aymptotic $t$ statistic. Just as in the classical linear model where the $F$ statistic for a single restriction is the square of the $t$ statistic for that same restriction, the Wald statistic is the square of the corresponding asymptotic $t$ statistic. Suppose that for a parameter $\theta_{1}>0$, the null hypothesis is $\mathrm{H}_{0}: \theta_{\mathrm{o} 1}=1$. The asymptotic $t$ statistic is $\left(\hat{\theta}_{1}-1\right) / \operatorname{se}\left(\hat{\theta}_{1}\right)$, where $\operatorname{se}\left(\hat{\theta}_{1}\right)$ is the asymptotic standard error of $\hat{\theta}_{1}$. Now define $\phi_{1}=\log \left(\theta_{1}\right)$, so that $\phi_{\mathrm{o} 1}=\log \left(\theta_{\mathrm{o} 1}\right)$ and $\hat{\phi}_{1}=\log \left(\hat{\theta}_{1}\right)$. The null hypothesis can be stated as $\mathrm{H}_{0}: \phi_{\mathrm{ol}}=0$. Using the delta method (see Chapter 3), $\operatorname{se}\left(\hat{\phi}_{1}\right)=\hat{\theta}_{1}^{-1} \operatorname{se}\left(\hat{\theta}_{1}\right)$, and so the $t$ statistic based on $\hat{\phi}_{1}$ is $\hat{\phi}_{1} / \operatorname{se}\left(\hat{\phi}_{1}\right)=\log \left(\hat{\theta}_{1}\right) \hat{\theta}_{1} / \operatorname{se}\left(\hat{\theta}_{1}\right) \neq\left(\hat{\theta}_{1}-1\right) / \operatorname{se}\left(\hat{\theta}_{1}\right)$.

The lack of invariance of the Wald statistic is discussed in more detail by Gregory and Veall (1985), Phillips and Park (1988), and Davidson and MacKinnon (1993, Sect. 13.6). The lack of invariance is a cause for concern because it suggests that the Wald statistic can have poor finite sample properties for testing nonlinear hypotheses. What is much less clear is that the lack of invariance has led empirical researchers to search over different statements of the null hypothesis in order to obtain a desired result.

\subsection*{12.6.2 Score (or Lagrange Multiplier) Tests}
In cases where the unrestricted model is difficult to estimate but the restricted model is relatively simple to estimate, it is convenient to have a statistic that only requires estimation under the null. Such a statistic is Rao's (1948) score statistic, also called the Lagrange multiplier statistic in econometrics, based on the work of Aitchison and Silvey (1958). We will focus on Rao's original motivation for the statistic because it leads more directly to test statistics that are used in econometrics. An important point is that, even though Rao, Aitchison and Silvey, Engle (1984), and many others focused on the maximum likelihood setup, the score principle is applicable to any problem where the estimators solve a first-order condition, including the general class of M estimators.

The score approach is ideally suited for specification testing. Typically, the first step in specification testing is to begin with a popular model-one that is relatively easy to estimate and interpret-and nest it within a more complicated model. Then the popular model is tested against the more general alternative to determine if the original model is misspecified. We do not want to estimate the more complicated model unless there is significant evidence against the restricted form of the model. In stating the null and alternative hypotheses, there is no difference between specification testing and classical tests of parameter restrictions. However, in practice, specification testing gives primary importance to the restricted model, and we may\\
have no intention of actually estimating the general model even if the null model is rejected.

We will derive the score test only in the case where no correction is needed for preliminary estimation of nuisance parameters: either there are no such parameters present, or assumption (12.37) holds under $\mathrm{H}_{0}$. If nuisance parameters are present, we do not explicitly show the score and Hessian depending on $\hat{\gamma}$.

We again assume that there are $Q$ continuously differentiable restrictions imposed on $\boldsymbol{\theta}_{\mathrm{o}}$ under $\mathrm{H}_{0}, \mathbf{c}\left(\boldsymbol{\theta}_{\mathrm{o}}\right)=\mathbf{O}$. However, we must also assume that the restrictions define a mapping from $\mathbb{R}^{P-Q}$ to $\mathbb{R}^{P}$, say, $\mathbf{d}: \mathbb{R}^{P-Q} \rightarrow \mathbb{R}^{P}$. In particular, under the null hypothesis, we can write $\boldsymbol{\theta}_{\mathrm{o}}=\mathbf{d}\left(\lambda_{\mathrm{o}}\right)$, where $\lambda_{\mathrm{o}}$ is a $(P-Q) \times 1$ vector. We must assume that $\lambda_{\mathrm{o}}$ is in the interior of its parameter space, $\boldsymbol{\Lambda}$, under $\mathrm{H}_{0}$. We also assume that $\mathbf{d}$ is twice continuously differentiable on the interior of $\boldsymbol{\Lambda}$.

Let $\tilde{\lambda}$ be the solution to the constrained minimization problem


\begin{equation*}
\min _{\lambda \in \Lambda} \sum_{i=1}^{N} q\left[\mathbf{w}_{i}, \mathbf{d}(\lambda)\right] \tag{12.64}
\end{equation*}


The constrained estimator of $\boldsymbol{\theta}_{\mathrm{o}}$ is simply $\tilde{\boldsymbol{\theta}} \equiv \mathbf{d}(\tilde{\boldsymbol{\lambda}})$. In practice, we do not have to explicitly find the function d; solving problem (12.64) is easily done just by directly imposing the restrictions, especially when the restrictions set certain parameters to hypothesized values (such as zero). Then, we just minimize the resulting objective function over the free parameters.

As an example, consider the nonlinear regression model

$$
m(\mathbf{x}, \boldsymbol{\theta})=\exp \left[\mathbf{x} \boldsymbol{\beta}+\delta_{1}(\mathbf{x} \boldsymbol{\beta})^{2}+\delta_{2}(\mathbf{x} \boldsymbol{\beta})^{3}\right]
$$

where $\mathbf{x}$ is $1 \times K$ and contains unity as its first element. The null hypothesis is $\mathrm{H}_{0}: \delta_{1}=\delta_{2}=0$, so that the model with the restrictions imposed is just an exponential regression function, $m(\mathbf{x}, \boldsymbol{\beta})=\exp (\mathbf{x} \boldsymbol{\beta})$.

The simplest method for deriving the LM test is to use Rao's score principle extended to the M-estimator case. The LM statistic is based on the limiting distribution of


\begin{equation*}
N^{-1 / 2} \sum_{i=1}^{N} \mathbf{s}_{i}(\tilde{\boldsymbol{\theta}}) \tag{12.65}
\end{equation*}


under $\mathrm{H}_{0}$. This is the score with respect to the entire vector $\boldsymbol{\theta}$, but we are evaluating it at the restricted estimates. If $\tilde{\boldsymbol{\theta}}$ were replaced by $\hat{\boldsymbol{\theta}}$, then expression (12.65) would be identically zero, which would make it useless as a test statistic. If the restrictions\\
imposed by the null hypothesis are true, then expression (12.65) will not be statistically different from zero.

Assume initially that $\boldsymbol{\theta}_{\mathrm{o}}$ is in the interior of $\boldsymbol{\Theta}$ under $\mathrm{H}_{0}$; we will discuss how to relax this assumption later. Now $\sqrt{N}\left(\tilde{\boldsymbol{\theta}}-\boldsymbol{\theta}_{\mathrm{o}}\right)=\mathrm{O}_{p}(1)$ by the delta method because $\sqrt{N}\left(\tilde{\boldsymbol{\lambda}}-\boldsymbol{\lambda}_{\mathrm{o}}\right)=\mathrm{O}_{p}(1)$ under the given assumptions. A standard mean value expansion yields


\begin{equation*}
N^{-1 / 2} \sum_{i=1}^{N} \mathbf{s}_{i}(\tilde{\boldsymbol{\theta}})=N^{-1 / 2} \sum_{i=1}^{N} \mathbf{s}_{i}\left(\boldsymbol{\theta}_{\mathrm{o}}\right)+\mathbf{A}_{\mathrm{o}} \sqrt{N}\left(\tilde{\boldsymbol{\theta}}-\boldsymbol{\theta}_{\mathrm{o}}\right)+\mathrm{o}_{p}(1) \tag{12.66}
\end{equation*}


under $\mathrm{H}_{0}$, where $\mathbf{A}_{\mathrm{o}}$ is given in expression (12.19). But $\mathbf{0}=\sqrt{N} \mathbf{c}(\tilde{\boldsymbol{\theta}})=\sqrt{N} \mathbf{c}\left(\boldsymbol{\theta}_{\mathrm{o}}\right)+ \ddot{\mathbf{C}} \sqrt{N}\left(\tilde{\boldsymbol{\theta}}-\boldsymbol{\theta}_{\mathrm{o}}\right)$, where $\ddot{\mathbf{C}}$ is the $Q \times P$ Jacobian matrix $\mathbf{C}(\boldsymbol{\theta})$ with rows evaluated at mean values between $\tilde{\boldsymbol{\theta}}$ and $\boldsymbol{\theta}_{\mathrm{o}}$. Under $\mathrm{H}_{0}, \mathbf{c}\left(\boldsymbol{\theta}_{\mathrm{o}}\right)=\mathbf{0}$, and plim $\ddot{\mathbf{C}}=\mathbf{C}\left(\boldsymbol{\theta}_{\mathrm{o}}\right) \equiv \mathbf{C}_{\mathrm{o}}$. Therefore, under $\mathrm{H}_{0}, \mathbf{C}_{\mathrm{o}} \sqrt{N}\left(\tilde{\boldsymbol{\theta}}-\boldsymbol{\theta}_{\mathrm{o}}\right)=\mathrm{o}_{p}(1)$, and so multiplying equation (12.66) through by $\mathbf{C}_{\mathrm{o}} \mathbf{A}_{\mathrm{o}}^{-1}$ gives


\begin{equation*}
\mathbf{C}_{\mathrm{o}} \mathbf{A}_{\mathrm{o}}^{-1} N^{-1 / 2} \sum_{i=1}^{N} \mathbf{s}_{i}(\tilde{\boldsymbol{\theta}})=\mathbf{C}_{\mathrm{o}} \mathbf{A}_{\mathrm{o}}^{-1} N^{-1 / 2} \sum_{i=1}^{N} \mathbf{s}_{i}\left(\boldsymbol{\theta}_{\mathrm{o}}\right)+\mathrm{o}_{p}(1) \tag{12.67}
\end{equation*}


By the CLT, $\mathbf{C}_{\mathrm{o}} \mathbf{A}_{\mathrm{o}}^{-1} N^{-1 / 2} \sum_{i=1}^{N} \mathbf{s}_{i}\left(\boldsymbol{\theta}_{\mathrm{o}}\right) \xrightarrow{d} \operatorname{Normal}\left(\mathbf{0}, \mathbf{C}_{\mathrm{o}} \mathbf{A}_{\mathrm{o}}^{-1} \mathbf{B}_{\mathrm{o}} \mathbf{A}_{\mathrm{o}}^{-1} \mathbf{C}_{\mathrm{o}}^{\prime}\right)$, where $\mathbf{B}_{\mathrm{o}}$ is defined in expression (12.20). Under our assumptions, $\mathbf{C}_{\mathrm{o}} \mathbf{A}_{\mathrm{o}}^{-1} \mathbf{B}_{\mathrm{o}} \mathbf{A}_{\mathrm{o}}^{-1} \mathbf{C}_{\mathrm{o}}^{\prime}$ has full rank $Q$, and so

$$
\left[N^{-1 / 2} \sum_{i=1}^{N} \mathbf{s}_{i}(\tilde{\boldsymbol{\theta}})\right]^{\prime} \mathbf{A}_{\mathrm{o}}^{-1} \mathbf{C}_{\mathrm{o}}^{\prime}\left(\mathbf{C}_{\mathrm{o}} \mathbf{A}_{\mathrm{o}}^{-1} \mathbf{B}_{\mathrm{o}} \mathbf{A}_{\mathrm{o}}^{-1} \mathbf{C}_{\mathrm{o}}^{\prime}\right)^{-1} \mathbf{C}_{\mathrm{o}} \mathbf{A}_{\mathrm{o}}^{-1}\left[N^{-1 / 2} \sum_{i=1}^{N} \mathbf{s}_{i}(\tilde{\boldsymbol{\theta}})\right] \xrightarrow{d} \chi_{Q}^{2}
$$

The score or LM statistic is given by\\
$L M \equiv\left(\sum_{i=1}^{N} \tilde{\mathbf{s}}_{i}\right)^{\prime} \tilde{\mathbf{A}}^{-1} \tilde{\mathbf{C}}^{\prime}\left(\tilde{\mathbf{C}} \tilde{\mathbf{A}}^{-1} \tilde{\mathbf{B}} \tilde{\mathbf{A}}^{-1} \tilde{\mathbf{C}}^{\prime}\right)^{-1} \tilde{\mathbf{C}} \tilde{\mathbf{A}}^{-1}\left(\sum_{i=1}^{N} \tilde{\mathbf{s}}_{i}\right) / N$,\\
where all quantities are evaluated at $\tilde{\boldsymbol{\theta}}$. For example, $\tilde{\mathbf{C}} \equiv \mathbf{C}(\tilde{\boldsymbol{\theta}}), \tilde{\mathbf{B}}$ is given in expression (12.45) but with $\tilde{\boldsymbol{\theta}}$ in place of $\hat{\boldsymbol{\theta}}$, and $\tilde{\mathbf{A}}$ is one of the estimators in expression (12.42) or (12.44), again evaluated at $\tilde{\boldsymbol{\theta}}$. Under $\mathrm{H}_{0}, L M \xrightarrow{d} \chi_{Q}^{2}$. Because $\tilde{\mathbf{B}}$ is at least p.s.d., $L M \geq 0$.

For the Wald statistic we assumed that $\boldsymbol{\theta}_{\mathrm{o}} \in \operatorname{int}(\boldsymbol{\Theta})$ under $\mathrm{H}_{0}$; this assumption is crucial for the statistic to have a limiting chi-square distribution. We will not consider the Wald statistic when $\boldsymbol{\theta}_{\mathrm{o}}$ is on the boundary of $\boldsymbol{\Theta}$ under $\mathrm{H}_{0}$; see Wolak (1991) for some results. The general derivation of the LM statistic also assumed that $\boldsymbol{\theta}_{\mathrm{o}} \in \operatorname{int}(\boldsymbol{\Theta})$\\
under $\mathrm{H}_{0}$. Nevertheless, for certain applications of the LM test we can drop the requirement that $\boldsymbol{\theta}_{\mathrm{o}}$ is in the interior of $\boldsymbol{\Theta}$ under $\mathrm{H}_{0}$. A leading case occurs when $\boldsymbol{\theta}$ can be partitioned as $\boldsymbol{\theta} \equiv\left(\boldsymbol{\theta}_{1}^{\prime}, \boldsymbol{\theta}_{2}^{\prime}\right)^{\prime}$, where $\boldsymbol{\theta}_{1}$ is $(P-Q) \times 1$ and $\boldsymbol{\theta}_{2}$ is $Q \times 1$. The null hypothesis is $\mathrm{H}_{0}: \boldsymbol{\theta}_{\mathrm{o} 2}=\mathbf{0}$, so that $\mathbf{c}(\boldsymbol{\theta}) \equiv \boldsymbol{\theta}_{2}$. It is easy to see that the mean value expansion used to derive the LM statistic is valid provided $\lambda_{\mathrm{o}} \equiv \boldsymbol{\theta}_{\mathrm{o} 1}$ is in the interior of its parameter space under $\mathrm{H}_{0} ; \boldsymbol{\theta}_{\mathrm{o}} \equiv\left(\boldsymbol{\theta}_{\mathrm{ol}}^{\prime}, \mathbf{0}\right)^{\prime}$ can be on the boundary of $\boldsymbol{\Theta}$. This observation is useful especially when testing hypotheses about parameters that must be either nonnegative or nonpositive.

If we assume the generalized information matrix equality (12.53) with $\sigma_{\mathrm{o}}^{2}=1$, the LM statistic simplifies. The simplification results from the following reasoning: (1) $\tilde{\mathbf{C}} \tilde{\mathbf{D}}=\mathbf{0}$ by the chain rule, where $\tilde{\mathbf{D}} \equiv \nabla_{\lambda} \mathbf{d}(\tilde{\boldsymbol{\lambda}})$, since $\mathbf{c}[\mathbf{d}(\lambda)] \equiv \mathbf{0}$ for $\lambda$ in $\boldsymbol{\Lambda}$. (2) If $\mathbf{E}$ is a $P \times Q$ matrix $\mathbf{E}$ with rank $Q, \mathbf{F}$ is a $P \times(P-Q)$ matrix with rank $P-Q$, and $\mathbf{E}^{\prime} \mathbf{F}=\mathbf{0}$, then $\mathbf{E}\left(\mathbf{E}^{\prime} \mathbf{E}\right)^{-1} \mathbf{E}^{\prime}=\mathbf{I}_{\mathrm{P}}-\mathbf{F}\left(\mathbf{F}^{\prime} \mathbf{F}\right)^{-1} \mathbf{F}^{\prime}$. (This is simply a statement about projections onto orthogonal subspaces.) Choosing $\mathbf{E} \equiv \tilde{\mathbf{A}}^{-1 / 2} \tilde{\mathbf{C}}^{\prime}$ and $\mathbf{F} \equiv \tilde{\mathbf{A}}^{1 / 2} \tilde{\mathbf{D}}$ gives $\tilde{\mathbf{A}}^{-1 / 2} \tilde{\mathbf{C}}^{\prime}\left(\tilde{\mathbf{C}} \tilde{\mathbf{A}}^{-1} \tilde{\mathbf{C}}^{\prime}\right)^{-1} \tilde{\mathbf{C}} \tilde{\mathbf{A}}^{-1 / 2}=\mathbf{I}_{\mathrm{P}}-\tilde{\mathbf{A}}^{1 / 2} \tilde{\mathbf{D}}\left(\tilde{\mathbf{D}}{ }^{\prime} \tilde{\mathbf{A}} \tilde{\mathbf{D}}\right)^{-1} \tilde{\mathbf{D}}^{\prime} \tilde{\mathbf{A}}^{1 / 2}$. Now, pre- and postmultiply this equality by $\tilde{\mathbf{A}}^{-1 / 2}$ to get $\tilde{\mathbf{A}}^{-1} \tilde{\mathbf{C}}^{\prime}\left(\tilde{\mathbf{C}} \tilde{\mathbf{A}}^{-1} \tilde{\mathbf{C}}^{\prime}\right)^{-1} \tilde{\mathbf{C}} \tilde{\mathbf{A}}^{-1}=\tilde{\mathbf{A}}^{-1}-\tilde{\mathbf{D}}\left(\tilde{\mathbf{D}^{\prime}} \tilde{\mathbf{A}} \tilde{\mathbf{D}}\right)^{-1} \tilde{\mathbf{D}}^{\prime}$. (3) Plug $\tilde{\mathbf{B}}=\tilde{\mathbf{A}}$ into expression (12.68) and use step 2, along with the first-order condition $\tilde{\mathbf{D}}^{\prime}\left(\sum_{i=1}^{N} \tilde{\mathbf{s}}_{i}\right)=\mathbf{0}$, to get\\
$L M=\left(\sum_{i=1}^{N} \tilde{\mathbf{s}}_{i}\right)^{\prime} \tilde{\mathbf{M}}^{-1}\left(\sum_{i=1}^{N} \tilde{\mathbf{s}}_{i}\right)$,\\
where $\tilde{\mathbf{M}}$ can be chosen as $\sum_{i=1}^{N} \tilde{\mathbf{A}}_{i}, \sum_{i=1}^{N} \tilde{\mathbf{H}}_{i}$, or $\sum_{i=1}^{N} \tilde{\mathbf{s}}_{i} \tilde{\mathbf{s}}_{i}^{\prime}$. (Each of these expressions consistently estimates $\mathbf{A}_{\mathrm{o}}=\mathbf{B}_{\mathrm{o}}$ when divided by $N$.) The last choice of $\tilde{\mathbf{M}}$ results in a statistic that is $N$ times the uncentered $R$-squared, say $R_{0}^{2}$, from the regression


\begin{equation*}
1 \text { on } \tilde{\mathbf{s}}_{i}^{\prime}, \quad i=1,2, \ldots, N . \tag{12.70}
\end{equation*}


(Recall that $\tilde{\mathbf{s}}_{i}^{\prime}$ is a $1 \times P$ vector.) Because the dependent variable in regression (12.70) is unity, $N R_{0}^{2}$ is equivalent to $N-\mathrm{SSR}_{0}$, where $\mathrm{SSR}_{0}$ is the sum of squared residuals from regression (12.70). This is often called the outer product of the score LM statistic because of the estimator it uses for $\mathbf{A}_{\mathrm{o}}$. While this statistic is simple to compute, there is ample evidence that it can have severe size distortions (typically, the null hypothesis is rejected much more often than the nominal size of the test). See, for example, Davidson and MacKinnon (1993), Bera and McKenzie (1986), Orme (1990), and Chesher and Spady (1991).

The Hessian form of the LM statistic uses $\tilde{\mathbf{M}}=\sum_{i=1}^{N} \tilde{\mathbf{H}}_{i}$, and it has a few drawbacks: (1) because $\tilde{\mathbf{M}}$ is the Hessian for the full model evaluated at the restricted\\
estimates it may not be positive definite, in which case the LM statistic can be negative; (2) it requires computation of the second derivatives; and (3) it is not invariant to reparameterizations. We will discuss the last problem later.

A statistic that always avoids the first problem, and often the second and third problems, is based on $\mathrm{E}\left[\mathbf{H}\left(\mathbf{w}, \boldsymbol{\theta}_{\mathrm{o}}\right) \mid \mathbf{x}\right]$, assuming that $\mathbf{w}$ partitions into endogenous variables $\mathbf{y}$ and exogenous variables $\mathbf{x}$. We call the LM statistic that uses $\tilde{\mathbf{M}}= \sum_{i=1}^{N} \tilde{\mathbf{A}}_{i}$ the expected Hessian form of the LM statistic. This name comes from the fact that the statistic is based on the conditional expectation of $\mathbf{H}\left(\mathbf{w}, \boldsymbol{\theta}_{\mathrm{o}}\right)$ given $\mathbf{x}$. When it can be computed, the expected Hessian form is usually preferred because it tends to have the best small sample properties.

The LM statistic in equation (12.69) is valid only when $\mathbf{B}_{\mathrm{o}}=\mathbf{A}_{\mathrm{o}}$, and therefore it is not robust to failures of auxiliary assumptions in some important models. If $\mathbf{B}_{\mathrm{o}} \neq \mathbf{A}_{\mathrm{o}}$, the limiting distribution of equation (12.69) is not chi-square and is not suitable for testing.

In the context of NLS, the expected Hessian form of the LM statistic needs to be modified for the presence of $\sigma_{\mathrm{o}}^{2}$, assuming that Assumption NLS. 3 holds under $\mathrm{H}_{0}$. Let $\tilde{\sigma}^{2} \equiv N^{-1} \sum_{i=1}^{N} \tilde{u}_{i}^{2}$ be the estimate of $\sigma_{\mathrm{o}}^{2}$ using the restricted estimator of $\boldsymbol{\theta}_{\mathrm{o}}: \tilde{u}_{i} \equiv y_{i}-m\left(\mathbf{x}_{i}, \tilde{\boldsymbol{\theta}}\right), i=1,2, \ldots, N$. It is customary not to make a degrees-of-freedom adjustment when estimating the variance using the null estimates, partly because the sum of squared residuals for the restricted model is always larger than for the unrestricted model. The score evaluated at the restricted estimates can be written as $\tilde{\mathbf{s}}_{i}=\nabla_{\theta} \tilde{m}_{i}^{\prime} \tilde{u}_{i}$. Thus the LM statistic that imposes homoskedasticity is\\
$L M=\left(\sum_{i=1}^{N} \nabla_{\theta} \tilde{m}_{i}^{\prime} \tilde{u}_{i}\right)^{\prime}\left(\sum_{i=1}^{N} \nabla_{\theta} \tilde{m}_{i}^{\prime} \nabla_{\theta} \tilde{m}_{i}\right)^{-1}\left(\sum_{i=1}^{N} \nabla_{\theta} \tilde{m}_{i}^{\prime} \tilde{u}_{i}\right) / \tilde{\sigma}^{2}$.

A little algebra shows that this expression is identical to $N$ times the uncentered $R$ squared, $R_{u}^{2}$, from the auxiliary regression\\
$\tilde{u}_{i}$ on $\nabla_{\theta} \tilde{m}_{i}, \quad i=1,2, \ldots, N$.

In other words, just regress the residuals from the restricted model on the gradient with respect to the unrestricted mean function but evaluated at the restricted estimates. Under $\mathrm{H}_{0}$ and Assumption NLS.3, $L M=N R_{u}^{2} \stackrel{a}{\sim} \chi_{Q}^{2}$.

In the nonlinear regression example with $m(\mathbf{x}, \boldsymbol{\theta})=\exp \left[\mathbf{x} \boldsymbol{\beta}+\delta_{1}(\mathbf{x} \boldsymbol{\beta})^{2}+\delta_{2}(\mathbf{x} \boldsymbol{\beta})^{3}\right]$, let $\tilde{\boldsymbol{\beta}}$ be the restricted NLS estimator with $\delta_{1}=0$ and $\delta_{2}=0$; in other words, $\tilde{\boldsymbol{\beta}}$ is from a nonlinear regression with an exponential regression function. The restricted residuals are $\tilde{u}_{i}=y_{i}-\exp \left(\mathbf{x}_{i} \tilde{\boldsymbol{\beta}}\right)$, and the gradient of $m(\mathbf{x}, \boldsymbol{\theta})$ with respect to all parameters, evaluated at the null, is\\
$\nabla_{\theta} m\left(\mathbf{x}_{i}, \boldsymbol{\beta}_{\mathrm{o}}, \mathbf{0}\right)=\left\{\mathbf{x}_{i} \exp \left(\mathbf{x}_{i} \boldsymbol{\beta}_{\mathrm{o}}\right),\left(\mathbf{x}_{i} \boldsymbol{\beta}_{\mathrm{o}}\right)^{2} \exp \left(\mathbf{x}_{i} \boldsymbol{\beta}_{\mathrm{o}}\right),\left(\mathbf{x}_{i} \boldsymbol{\beta}_{\mathrm{o}}\right)^{3} \exp \left(\mathbf{x}_{i} \boldsymbol{\beta}_{\mathrm{o}}\right)\right\}$.\\
Plugging in $\tilde{\boldsymbol{\beta}}$ gives $\nabla_{\theta} \tilde{m}_{i}=\left[\mathbf{x}_{i} \tilde{m}_{i},\left(\mathbf{x}_{i} \tilde{\boldsymbol{\beta}}\right)^{2} \tilde{m}_{i},\left(\mathbf{x}_{i} \tilde{\boldsymbol{\beta}}\right)^{3} \tilde{m}_{i}\right]$, where $\tilde{m}_{i} \equiv \exp \left(\mathbf{x}_{i} \tilde{\boldsymbol{\beta}}\right)$. Regression (12.72) becomes\\
$\tilde{u}_{i}$ on $\mathbf{x}_{i} \tilde{m}_{i},\left(\mathbf{x}_{i} \tilde{\boldsymbol{\beta}}\right)^{2} \tilde{m}_{i},\left(\mathbf{x}_{i} \tilde{\boldsymbol{\beta}}\right)^{3} \tilde{m}_{i}, \quad i=1,2, \ldots, N$.

Under $\mathrm{H}_{0}$ and homoskedasticity, $N R_{u}^{2} \sim \chi_{2}^{2}$, since there are two restrictions being tested. This is a fairly simple way to test the exponential functional form without ever estimating the more complicated alternative model. Other models that nest the exponential model are discussed in Wooldridge (1992).

This example illustrates an important point: even though $\sum_{i=1}^{N}\left(\mathbf{x}_{i} \tilde{m}_{i}\right)^{\prime} \tilde{u}_{i}$ is identically zero by the first-order condition for NLS, the term $\mathbf{x}_{i} \tilde{m}_{i}$ must generally be included in regression (12.73). The $R$-squared from the regression without $\mathbf{x}_{i} \tilde{m}_{i}$ will be different because the remaining regressors in regression (12.73) are usually correlated with $\mathbf{x}_{i} \tilde{m}_{i}$ in the sample. (More important, for $h=2$ and $3,\left(\mathbf{x}_{i} \boldsymbol{\beta}_{\mathrm{o}}\right)^{h} \exp \left(\mathbf{x}_{i} \boldsymbol{\beta}_{\mathrm{o}}\right)$ is probably correlated with $\mathbf{x}_{i} \exp \left(\mathbf{x}_{i} \boldsymbol{\beta}_{\mathrm{o}}\right)$ in the population.) As a general rule, the entire gradient $\nabla_{\theta} \tilde{m}_{i}$ must appear in the auxiliary regression.

In order to be robust against failure of Assumption NLS.3, the more general form of the statistic in expression (12.68) should be used. Fortunately, this statistic also can be easily computed for most hypotheses. Partition $\boldsymbol{\theta}$ into the $(P-Q) \times 1$ vector $\boldsymbol{\beta}$ and the $Q$ vector $\boldsymbol{\delta}$. Assume that the null hypothesis is $\mathrm{H}_{0}: \boldsymbol{\delta}_{\mathrm{o}}=\overline{\boldsymbol{\delta}}$, where $\overline{\boldsymbol{\delta}}$ is a prespecified vector (often containing all zeros, but not always). Let $\nabla_{\beta} \tilde{m}_{i} [1 \times(P-Q)]$ and $\nabla_{\delta} \tilde{\boldsymbol{m}}_{i}(1 \times Q)$ denote the gradients with respect to $\boldsymbol{\beta}$ and $\boldsymbol{\delta}$, respectively, evaluated at $\tilde{\boldsymbol{\beta}}$ and $\overline{\boldsymbol{\delta}}$. After tedious algebra, and using the special structure $\mathbf{C}(\boldsymbol{\theta})=\left[\mathbf{0} \mid \mathbf{I}_{Q}\right]$, where $\mathbf{0}$ is a $Q \times(P-Q)$ matrix of zero, the following procedure can be shown to produce expression (12.68):

\begin{enumerate}
  \item Run a multivariate regression\\
$\nabla_{\delta} \tilde{m}_{i}$ on $\nabla_{\beta} \tilde{m}_{i}, \quad i=1,2, \ldots, N$\\
and save the $1 \times Q$ vector residuals, say $\tilde{\mathbf{r}}_{i}$. Then, for each $i$, form $\tilde{u}_{i} \tilde{\mathbf{r}}_{i}$. (That is, multiply $\tilde{u}_{i}$ by each element of $\tilde{\mathbf{r}}_{i}$.)
  \item $L M=N-\mathrm{SSR}_{0}=N R_{0}^{2}$ from the regression
\end{enumerate}

1 on $\tilde{u}_{i} \tilde{\mathbf{r}}_{i}, \quad i=1,2, \ldots, N$\\
where $\mathrm{SSR}_{0}$ is the usual sum of squared residuals and $R_{0}^{2}$ is the uncentered Rsquared. This step produces a statistic that has a limiting $\chi_{Q}^{2}$ distribution whether or not Assumption NLS. 3 holds. See Wooldridge (1991a) for more discussion.

We can illustrate the heteroskedasticity-robust test using the preceding exponential model. Regression (12.74) is the same as regressing each of $\left(\mathbf{x}_{i} \tilde{\boldsymbol{\beta}}\right)^{2} \tilde{m}_{i}$ and $\left(\mathbf{x}_{i} \tilde{\boldsymbol{\beta}}\right)^{3} \tilde{m}_{i}$ onto $\mathbf{x}_{i} \tilde{m}_{i}$, and saving the residuals $\tilde{r}_{i 1}$ and $\tilde{r}_{i 2}$, respectively ( $N$ each). Then, regression (12.75) is simply 1 on $\tilde{u}_{i} \tilde{r}_{i 1}, \tilde{u}_{i} \tilde{r}_{i 2}$. The number of regressors in the final regression of the robust test is always the same as the degrees of freedom of the test.

Finally, these procedures are easily modified for WNLS. Simply multiply both $\tilde{u}_{i}$ and $\nabla_{\theta} \tilde{m}_{i}$ by $1 / \sqrt{\tilde{h}_{i}}$, where the variance estimates $\tilde{h}_{i}$ are based on the null model (so we use a $\sim$ rather than a $\wedge$ ). The nonrobust LM statistic that maintains Assumption WNLS. 3 is obtained as in regression (12.72). The robust form, which allows $\operatorname{Var}(y \mid \mathbf{x}) \neq \sigma_{\mathrm{o}}^{2} h\left(\mathbf{x}, \gamma_{\mathrm{o}}\right)$, follows exactly as in regressions (12.74) and (12.75).

In examples like the previous one, there is an alternative, asymptotically equivalent method of obtaining a test statistic that, like the score test, only requires estimation under the null hypothesis. A variable addition test (VAT) is obtained by adding (estimated) terms to a standard model. Once the additional variables have been defined, calculation is typically straightforward with existing software, and robust tests-robust to heteroskedasticity in the case of NLS-are easy to obtain.

In the previous example, where we focused on NLS, we again obtain $\tilde{\boldsymbol{\beta}}$ from NLS using the exponential regression function. We then compute $\left(\mathbf{x}_{i} \tilde{\boldsymbol{\beta}}\right)^{2}$ and $\left(\mathbf{x}_{i} \tilde{\boldsymbol{\beta}}\right)^{3}$. Next, we estimate an expanded exponential mean function that includes the original regressors, $\mathbf{x}_{i}$, and the two additional regressors, say, $\tilde{z}_{i 1} \equiv\left(\mathbf{x}_{i} \tilde{\boldsymbol{\beta}}\right)^{2}$ and $\tilde{z}_{i 2} \equiv\left(\mathbf{x}_{i} \tilde{\boldsymbol{\beta}}\right)^{3}$. That is, the auxiliary model (used for testing purposes only) is $\exp \left(\mathbf{x}_{i} \boldsymbol{\beta}+\alpha_{1} \tilde{z}_{i 1}+\alpha_{2} \tilde{z}_{i 2}\right)$. The VAT is obtained as a standard Wald test for joint exclusion of $\tilde{z}_{i 1}$ and $\tilde{z}_{i 2}$ (and therefore has an asymptotic $\chi_{2}^{2}$ distribution). For the nonrobust test, one difference between the LM test and the VAT is that the latter uses a different variance estimate (because of the additional terms $\tilde{z}_{i 1}$ and $\tilde{z}_{i 2}$ that have coefficients estimated rather than set to zero). Under the null, both estimators converge to $\sigma_{\mathrm{o}}^{2}$. It is easy to make VAT tests robust to heteroskedasticity whenever one is using an econometrics package that computes heteroskedasticity-robust Wald tests of exclusion restrictions; typically, it simply requires adding a "robust" option to an estimation command such as NLS.

For NLS (and WNLS), the VAT approach can be applied more generally when the mean function has the form $m(\mathbf{x}, \boldsymbol{\beta}, \boldsymbol{\delta})=R(a(\mathbf{x} \boldsymbol{\beta}, \mathbf{x}, \boldsymbol{\delta}))$, where $a(k, \mathbf{x}, \mathbf{0})=k$ and $\partial a(k, \mathbf{x}, \mathbf{0}) / \partial k=1$. Then $m(\mathbf{x}, \boldsymbol{\beta}, \mathbf{0})=R(\mathbf{x} \boldsymbol{\beta})$ and $\nabla_{\boldsymbol{\beta}} m(\mathbf{x}, \boldsymbol{\beta}, 0)=r(\mathbf{x} \boldsymbol{\beta}) \cdot \mathbf{x}$, where $r(\cdot)$ is the derivative of $R(\cdot)$. The vector of variables to be added is simply $\tilde{\mathbf{z}}_{i} \equiv \nabla_{\delta} a\left(\mathbf{x}_{i} \tilde{\boldsymbol{\beta}}, \mathbf{x}_{i}, \mathbf{0}\right)$, a $1 \times Q$ vector, where $\tilde{\boldsymbol{\beta}}$ is the NLS estimator using the mean function $R\left(\mathbf{x}_{i} \boldsymbol{\beta}\right)$. In other words, we use the augmented regression function $R\left(\mathbf{x}_{i} \boldsymbol{\beta}+\tilde{\mathbf{z}}_{i} \boldsymbol{\alpha}\right)$ in NLS and test joint exclusion of $\tilde{\mathbf{z}}_{i}$. Note that this "model" is not correctly specified under the alternative. We estimate the augmented equation to\\
obtain a simple test. The VAT procedure is easy to implement when the mean function $R(\cdot)$ with linear functions inside is easily estimated. (Examples include the exponential and logistic functions, and some others we encounter in Part IV.) We will see how to obtain VAT tests with other estimation methods in future chapters.

The invariance issue for the score statistic is somewhat complicated, but several results are known. First, it is easy to see that the outer product form of the statistic is invariant to differentiable reparameterizations. Write $\boldsymbol{\phi}=\mathbf{g}(\boldsymbol{\theta})$ as a twice continuously differentiable, invertible reparameterization; thus the $P \times P$ Jacobian of $\mathbf{g}$, $\mathbf{G}(\boldsymbol{\theta})$, is nonsingular for all $\boldsymbol{\theta} \in \boldsymbol{\Theta}$. The objective function in terms of $\boldsymbol{\phi}$ is $q^{g}(\mathbf{w}, \boldsymbol{\phi})$, and we must have $q^{g}[\mathbf{w}, \mathbf{g}(\boldsymbol{\theta})]=q(\mathbf{w}, \boldsymbol{\theta})$ for all $\boldsymbol{\theta} \in \boldsymbol{\Theta}$. Differentiating and transposing gives $\mathbf{s}(\mathbf{w}, \boldsymbol{\theta})=\mathbf{G}(\boldsymbol{\theta})^{\prime} \mathbf{s}^{g}[\mathbf{w}, \mathbf{g}(\boldsymbol{\theta})]$, where $\mathbf{s}^{g}(\mathbf{w}, \boldsymbol{\phi})$ is the score of $q^{g}[\mathbf{w}, \boldsymbol{\phi}]$. If $\tilde{\boldsymbol{\phi}}$ is the restricted estimator of $\boldsymbol{\phi}$, then $\tilde{\boldsymbol{\phi}}=\mathbf{g}(\tilde{\boldsymbol{\theta}})$, and so, for each observation $i, \tilde{\mathbf{s}}_{i}^{g}=\left(\tilde{\mathbf{G}}^{\prime}\right)^{-1} \tilde{\mathbf{s}}_{i}$. Plugging this equation into the LM statistic in equation (12.69), with $\tilde{\mathbf{M}}$ chosen as the outer product form, shows that the statistic based on $\tilde{\mathbf{s}}_{i}^{g}$ is identical to that based on $\tilde{\mathbf{s}}_{i}$.

Score statistics based on the estimated Hessian are not generally invariant to reparameterization because they can involve second derivatives of the function $\mathbf{g}(\boldsymbol{\theta})$; see Davidson and MacKinnon (1993, Sect. 13.6) for details. However, when $\mathbf{w}$ partitions as $(\mathbf{x}, \mathbf{y})$, score statistics based on the expected Hessian (conditional on $\mathbf{x}), \mathbf{A}(\mathbf{x}, \boldsymbol{\theta})$, are often invariant. In Chapter 13 we will see that this is always the case for conditional maximum likelihood estimation. Invariance also holds for NLS and WNLS for both the usual and robust LM statistics because any reparameterization comes through the conditional mean. Predicted values and residuals are invariant to reparameterization, and the statistics obtained from regressions (12.72) and (12.75) only involve the residuals and first derivatives of the conditional mean function. As with the outer product LM statistic, the Jacobian in the first derivative cancels out.

\subsection*{12.6.3 Tests Based on the Change in the Objective Function}
When both the restricted and unrestricted models are easy to estimate, a test based on the change in the objective function can greatly simplify the mechanics of obtaining a test statistic: we only need to obtain the value of the objective function with and without the restrictions imposed. However, the computational simplicity comes at a price in terms of robustness. Unlike the Wald and score tests, a test based on the change in the objective function cannot be made robust to general failure of assumption (12.53). Therefore, throughout this subsection we assume that the generalized information matrix equality holds. Because the minimized objective function is invariant with respect to any reparameterization, the test statistic is invariant.

In the context of two-step estimators, we must also assume that $\hat{\gamma}$ has no effect on the asymptotic distribution of the M-estimator. That is, we maintain assumption\\
(12.37) when nuisance parameter estimates appear in the objective function (see Problem 12.8).

We first consider the case where $\sigma_{\mathrm{o}}^{2}=1$, so that $\mathbf{B}_{\mathrm{o}}=\mathbf{A}_{\mathrm{o}}$. Using a second-order Taylor expansion,\\
$\sum_{i=1}^{N} q\left(\mathbf{w}_{i}, \tilde{\boldsymbol{\theta}}\right)-\sum_{i=1}^{N} q\left(\mathbf{w}_{i}, \hat{\boldsymbol{\theta}}\right)=\left(\sum_{i=1}^{N} \mathbf{s}_{i}(\hat{\boldsymbol{\theta}})\right)\left(\hat{\boldsymbol{\theta}}-\boldsymbol{\theta}_{o}\right)+(1 / 2)(\tilde{\boldsymbol{\theta}}-\hat{\boldsymbol{\theta}})^{\prime}\left(\sum_{i=1}^{N} \ddot{\mathbf{H}}_{i}\right)(\tilde{\boldsymbol{\theta}}-\hat{\boldsymbol{\theta}})$,\\
where $\ddot{\mathbf{H}}_{i}$ is the $P \times P$ Hessian evaluate at mean values between $\tilde{\boldsymbol{\theta}}$ and $\hat{\boldsymbol{\theta}}$. Therefore, under $\mathrm{H}_{0}$ (using the first-order condition for $\hat{\boldsymbol{\theta}}$ ), we have\\
$2\left[\sum_{i=1}^{N} q\left(\mathbf{w}_{i}, \tilde{\boldsymbol{\theta}}\right)-\sum_{i=1}^{N} q\left(\mathbf{w}_{i}, \hat{\boldsymbol{\theta}}\right)\right]=[\sqrt{N}(\tilde{\boldsymbol{\theta}}-\hat{\boldsymbol{\theta}})]^{\prime} \mathbf{A}_{0}[\sqrt{N}(\tilde{\boldsymbol{\theta}}-\hat{\boldsymbol{\theta}})]+\mathrm{o}_{p}(1)$,\\
since $N^{-1} \sum_{i=1}^{N} \ddot{\mathbf{H}}_{i}=\mathbf{A}_{\mathrm{o}}+o_{p}$ (1) and $\sqrt{N}(\tilde{\boldsymbol{\theta}}-\hat{\boldsymbol{\theta}})=\mathrm{O}_{p}(1)$. In fact, it follows from equations (12.33) (without $\hat{\gamma}$ ) and (12.66) that $\sqrt{N}(\tilde{\boldsymbol{\theta}}-\hat{\boldsymbol{\theta}})=\mathbf{A}_{\mathrm{o}}^{-1} N^{-1 / 2} \sum_{i=1}^{N} \mathbf{s}_{i}(\tilde{\boldsymbol{\theta}})+ \mathrm{o}_{p}(1)$. Plugging this equation into equation (12.76) shows that


\begin{align*}
Q L R & \equiv 2\left[\sum_{i=1}^{N} q\left(\mathbf{w}_{i}, \tilde{\boldsymbol{\theta}}\right)-\sum_{i=1}^{N} q\left(\mathbf{w}_{i}, \hat{\boldsymbol{\theta}}\right)\right] \\
& =\left(N^{-1 / 2} \sum_{i=1}^{N} \tilde{\mathbf{s}}_{i}\right)^{\prime} \mathbf{A}_{\mathrm{o}}^{-1}\left(N^{-1 / 2} \sum_{i=1}^{N} \tilde{\mathbf{s}}_{i}\right)+\mathrm{o}_{p}(1), \tag{12.77}
\end{align*}


so that $Q L R$ has the same limiting distribution, $\chi_{Q}^{2}$, as the LM statistic under $\mathrm{H}_{0}$. (See equation (12.69), remembering that $\operatorname{plim}(\tilde{\mathbf{M}} / N)=\mathbf{A}_{\mathrm{o}}$.) We call statistic (12.77) the quasi-likelihood ratio (QLR) statistic, which comes from the fact that the leading example of equation (12.77) is the likelihood ratio statistic in the context of maximum likelihood estimation, as we will see in Chapter 13. We could also call equation (12.77) a criterion function statistic, as it is based on the difference in the criterion or objective function with and without the restrictions imposed.

When nuisance parameters are present, the same estimate, say $\hat{\gamma}$, should be used in obtaining the restricted and unrestricted estimates. This is to ensure that $Q L R$ is nonnegative given any sample. Typically, $\hat{\gamma}$ would be based on initial estimation of the unrestricted model.

If $\sigma_{\mathrm{o}}^{2} \neq 1$, we simply divide $Q L R$ by $\hat{\sigma}^{2}$, which is a consistent estimator of $\sigma_{\mathrm{o}}^{2}$ obtained from the unrestricted estimation. For example, consider NLS under Assumptions NLS.1-NLS.3. When equation (12.77) is divided by $\hat{\sigma}^{2}$ in equation (12.57), we obtain $\left(\operatorname{SSR}_{r}-\operatorname{SSR}_{u r}\right) /\left[\operatorname{SSR}_{u r} /(N-P)\right]$, where $\operatorname{SSR}_{r}$ and $\operatorname{SSR}_{u r}$ are the\\
restricted and unrestricted sums of squared residuals. Sometimes an $F$ version of this statistic is used instead, which is obtained by dividing the chi-square version by $Q$ :\\
$F=\frac{\left(\mathrm{SSR}_{r}-\mathrm{SSR}_{u r}\right)}{\mathrm{SSR}_{u r}} \cdot \frac{(N-P)}{Q}$.

This has exactly the same form as the $F$ statistic from classical linear regression analysis. Under the null hypothesis and homoskedasticity, $F$ can be treated as having an approximate $\mathscr{F}_{Q, N-P}$ distribution. (As always, this association is justified because $Q \cdot \mathscr{F}_{Q, N-P} \stackrel{a}{\sim} \chi_{Q}^{2}$ as $N-P \rightarrow \infty$.) Some authors (for example, Gallant, 1987) have found that $F$ has better finite-sample properties than the chi-square version of the statistic.

For weighted NLS, the same statistic works under Assumption WNLS. 3 provided the residuals (both restricted and unrestricted) are weighted by $1 / \sqrt{\hat{h}_{i}}$, where the $\hat{h}_{i}$ are obtained from estimation of the unrestricted model.

\subsection*{12.6.4 Behavior of the Statistics under Alternatives}
To keep the notation and assumptions as simple as possible, and to focus on the computation of valid test statistics under various assumptions, we have only derived the limiting distribution of the classical test statistics under the null hypothesis. It is also important to know how the tests behave under alternative hypotheses in order to choose a test with the highest power.

All the tests we have discussed are consistent against the alternatives they are specifically designed against. While test consistency is desirable, it tells us nothing about the likely finite-sample power that a statistic will have against particular alternatives. A framework that allows us to say more uses the notion of a sequence of local alternatives. Specifying a local alternative is a device that can approximate the finitesample power of test statistics for alternatives "close" to $\mathrm{H}_{0}$. If the null hypothesis is $\mathrm{H}_{0}: \mathbf{c}\left(\boldsymbol{\theta}_{\mathrm{o}}\right)=\mathbf{0}$, then a sequence of local alternatives is\\
$\mathrm{H}_{1}^{N}: \mathbf{c}\left(\boldsymbol{\theta}_{\mathrm{o}, N}\right)=\boldsymbol{\delta}_{\mathrm{o}} / \sqrt{N}$,\\
where $\boldsymbol{\delta}_{\mathrm{o}}$ is a given $Q \times 1$ vector. As $N \rightarrow \infty, \mathrm{H}_{1}^{N}$ approaches $\mathrm{H}_{0}$, since $\boldsymbol{\delta}_{\mathrm{o}} / \sqrt{N} \rightarrow \mathbf{0}$. The division by $\sqrt{N}$ means that the alternatives are local: for given $N$, equation (12.79) is an alternative to $\mathrm{H}_{0}$, but as $N \rightarrow \infty$, the alternative gets closer to $\mathrm{H}_{0}$. Dividing $\boldsymbol{\delta}_{\mathrm{o}}$ by $\sqrt{N}$ ensures that each of the statistics has a well-defined limiting distribution under the alternative that differs from the limiting distribution under $\mathrm{H}_{0}$.

It can be shown that, under equation (12.79), the general forms of the Wald and LM statistics have a limiting noncentral chi-square distribution with $Q$ degrees of\\
freedom under the regularity conditions used to obtain their null limiting distributions. The noncentrality parameter depends on $\mathbf{A}_{\mathrm{o}}, \mathbf{B}_{\mathrm{o}}, \mathbf{C}_{\mathrm{o}}$, and $\boldsymbol{\delta}_{\mathrm{o}}$, and can be estimated by using consistent estimators of $\mathbf{A}_{\mathrm{o}}, \mathbf{B}_{\mathrm{o}}$, and $\mathbf{C}_{\mathrm{o}}$. When we add assumption (12.53), then the special versions of the Wald and LM statistics and the QLR statistics have limiting noncentral chi-square distributions. For various $\boldsymbol{\delta}_{\mathrm{o}}$, we can estimate what is known as the asymptotic local power of the test statistics by computing probabilities from noncentral chi-square distributions.

Consider the Wald statistic where $\mathbf{B}_{\mathrm{O}}=\mathbf{A}_{\mathrm{O}}$. Denote by $\boldsymbol{\theta}_{\mathrm{O}}$ the limit of $\boldsymbol{\theta}_{\mathrm{O}, N}$ as $N \rightarrow \infty$. The usual mean value expansion under $\mathrm{H}_{1}^{N}$ gives\\
$\sqrt{N} \mathbf{c}(\hat{\boldsymbol{\theta}})=\boldsymbol{\delta}_{\mathrm{o}}+\mathbf{C}\left(\boldsymbol{\theta}_{\mathrm{o}}\right) \sqrt{N}\left(\hat{\boldsymbol{\theta}}-\boldsymbol{\theta}_{\mathrm{o}, N}\right)+\mathrm{o}_{p}(1)$\\
and, under standard assumptions, $\sqrt{N}\left(\hat{\boldsymbol{\theta}}-\boldsymbol{\theta}_{\mathrm{o}, N}\right) \stackrel{a}{\sim} \operatorname{Normal}\left(\mathbf{0}, \mathbf{A}_{\mathrm{o}}^{-1}\right)$. Therefore, $\sqrt{N} \mathbf{c}(\hat{\boldsymbol{\theta}}) \stackrel{a}{\sim} \operatorname{Normal}\left(\boldsymbol{\delta}_{\mathrm{o}}, \mathbf{C}_{\mathrm{o}} \mathbf{A}_{\mathrm{o}}^{-1} \mathbf{C}_{\mathrm{o}}^{\prime}\right)$ under the sequence (12.79). This result implies that the Wald statistic has a limiting noncentral chi-square distribution with $Q$ degrees of freedom and noncentrality parameter $\boldsymbol{\delta}_{\mathrm{o}}^{\prime}\left(\mathbf{C}_{\mathrm{o}} \mathbf{A}_{\mathrm{o}}^{-1} \mathbf{C}_{\mathrm{o}}^{\prime}\right)^{-1} \boldsymbol{\delta}_{\mathrm{o}}$. This turns out to be the same noncentrality parameter for the LM and QLR statistics when $\mathbf{B}_{\mathrm{o}}=\mathbf{A}_{\mathrm{o}}$. The details are similar to those under $\mathrm{H}_{0}$; see, for example, Gallant (1987, Sect. 3.6).

The statistic with the largest noncentrality parameter has the largest asymptotic local power. For choosing among the Wald, LM, and QLR statistics, this criterion does not help: they all have the same noncentrality parameters under the local alternatives (12.79) and assumption (12.53). Without assumption (12.53), the robust Wald and LM statistics have the same noncentrality parameter.

The notion of local alternatives is useful when choosing among statistics based on different estimators. Not surprisingly, the more efficient estimator produces tests with the best asymptotic local power under standard assumptions. But we should keep in mind the efficiency versus robustness trade-off, especially when efficient test statistics are computed under tenuous assumptions.

General analyses under local alternatives are available in Gallant (1987), Gallant and White (1988), and White (1994). See Andrews (1989) for innovative suggestions for using local power analysis in applied work.

\subsection*{12.7 Optimization Methods}
In this section we briefly discuss three iterative schemes that can be used to solve the general minimization problem (12.8) or (12.31). In the latter case, the minimization is only over $\boldsymbol{\theta}$, so the presence of $\hat{\gamma}$ changes nothing. If $\hat{\gamma}$ is present, the score and

Hessian with respect to $\boldsymbol{\theta}$ are simply evaluated at $\hat{\gamma}$. These methods are closely related to the asymptotic variance matrix estimators and test statistics we discussed in Sections 12.5 and 12.6.

\subsection*{12.7.1 Newton-Raphson Method}
Iterative methods are defined by an algorithm for going from one iteration to the next. Let $\boldsymbol{\theta}^{\{g\}}$ be the $P \times 1$ vector on the $g$ th iteration, and let $\boldsymbol{\theta}^{\{g+1\}}$ be the value on the next iteration. To motivate how we get from $\boldsymbol{\theta}^{\{g\}}$ to $\boldsymbol{\theta}^{\{g+1\}}$, use a mean value expansion (row by row) to write\\
$\sum_{i=1}^{N} \mathbf{s}_{i}\left(\boldsymbol{\theta}^{\{g+1\}}\right)=\sum_{i=1}^{N} \mathbf{s}_{i}\left(\boldsymbol{\theta}^{\{g\}}\right)+\left[\sum_{i=1}^{N} \mathbf{H}_{i}\left(\boldsymbol{\theta}^{\{g\}}\right)\right]\left(\boldsymbol{\theta}^{\{g+1\}}-\boldsymbol{\theta}^{\{g\}}\right)+\mathbf{r}^{\{g\}}$,\\
where $\mathbf{s}_{i}(\boldsymbol{\theta})$ is the $P \times 1$ score with respect to $\boldsymbol{\theta}$, evaluated at observation $i, \mathbf{H}_{i}(\boldsymbol{\theta})$ is the $P \times P$ Hessian, and $\mathbf{r}^{\{g\}}$ is a $P \times 1$ vector of remainder terms. We are trying to find the solution $\hat{\boldsymbol{\theta}}$ to equation (12.14). If $\boldsymbol{\theta}^{\{g+1\}}=\hat{\boldsymbol{\theta}}$, then the left-hand side of equation (12.80) is zero. Setting the left-hand side to zero, ignoring $\mathbf{r}^{\{g\}}$, and assuming that the Hessian evaluated at $\boldsymbol{\theta}^{\{g\}}$ is nonsingular, we can write\\
$\boldsymbol{\theta}^{\{g+1\}}=\boldsymbol{\theta}^{\{g\}}-\left[\sum_{i=1}^{N} \mathbf{H}_{i}\left(\boldsymbol{\theta}^{\{g\}}\right)\right]^{-1}\left[\sum_{i=1}^{N} \mathbf{s}_{i}\left(\boldsymbol{\theta}^{\{g\}}\right)\right]$.

Equation (12.81) provides an iterative method for finding $\hat{\boldsymbol{\theta}}$. To begin the iterations we must choose a vector of starting values; call this vector $\boldsymbol{\theta}^{\{0\}}$. Good starting values are often difficult to come by, and sometimes we must experiment with several choices before the problem converges. Ideally, the iterations wind up at the same place regardless of the starting values, but this outcome is not guaranteed. Given the starting values, we plug $\boldsymbol{\theta}^{\{0\}}$ into the right-hand side of equation (12.81) to get $\boldsymbol{\theta}^{\{1\}}$. Then, we plug $\boldsymbol{\theta}^{\{1\}}$ into equation (12.81) to get $\boldsymbol{\theta}^{\{2\}}$, and so on.

If the iterations are proceeding toward the minimum, the increments $\boldsymbol{\theta}^{\{g+1\}}-\boldsymbol{\theta}^{\{g\}}$ will eventually become very small: as we near the solution, $\sum_{i=1}^{N} \mathbf{s}_{i}\left(\boldsymbol{\theta}^{\{g\}}\right)$ gets close to zero. Some use as a stopping rule the requirement that the largest absolute change $\left|\boldsymbol{\theta}_{j}^{\{g+1\}}-\boldsymbol{\theta}_{j}^{\{g\}}\right|$, for $j=1,2, \ldots, P$, be smaller than some small constant; others prefer to look at the largest percentage change in the parameter values.

Another popular stopping rule is based on the quadratic form


\begin{equation*}
\left[\sum_{i=1}^{N} \mathbf{s}_{i}\left(\boldsymbol{\theta}^{\{g\}}\right)\right]^{\prime}\left[\sum_{i=1}^{N} \mathbf{H}_{i}\left(\boldsymbol{\theta}^{\{g\}}\right)\right]^{-1}\left[\sum_{i=1}^{N} \mathbf{s}_{i}\left(\boldsymbol{\theta}^{\{g\}}\right)\right], \tag{12.82}
\end{equation*}


where the iterations stop when expression (12.82) is less than some suitably small number, say . 0001 .

The iterative scheme just outlined is usually called the Newton-Raphson method. It is known to work in a variety of circumstances. Our motivation here has been heuristic, and we will not investigate situations under which the Newton-Raphson method does not work well. (See, for example, Quandt, 1983, for some theoretical results.) The Newton-Raphson method has some drawbacks. First, it requires computing the second derivatives of the objective function at every iteration. These calculations are not very taxing if closed forms for the second partials are available, but in many cases they are not. A second problem is that, as we saw for the case of nonlinear least squares, the Hessian evaluated at a particular value of $\boldsymbol{\theta}$ may not be positive definite. If the inverted Hessian in expression (12.81) is not positive definite, the procedure may head in the wrong direction.

We should always check that progress is being made from one iteration to the next by computing the difference in the values of the objective function from one iteration to the next:


\begin{equation*}
\sum_{i=1}^{N} q_{i}\left(\boldsymbol{\theta}^{\{g+1\}}\right)-\sum_{i=1}^{N} q_{i}\left(\boldsymbol{\theta}^{\{g\}}\right) . \tag{12.83}
\end{equation*}


Because we are minimizing the objective function, we should not take the step from $g$ to $g+1$ unless expression (12.83) is negative. (If we are maximizing the function, the iterations in equation (12.81) can still be used because the expansion in equation (12.80) is still appropriate, but then we want expression (12.83) to be positive.)

A slight modification of the Newton-Raphson method is sometimes useful to speed up convergence: multiply the Hessian term in expression (12.81) by a positive number, say $r$, known as the step size. Sometimes the step size $r=1$ produces too large a change in the parameters. If the objective function does not decrease using $r=1$, then try, say, $r=\frac{1}{2}$. Again, check the value of the objective function. If it has now decreased, go on to the next iteration (where $r=1$ is usually used at the beginning of each iteration); if the objective function still has not decreased, replace $r$ with, say, $\frac{1}{4}$. Continue halving $r$ until the objective function decreases. If you have not succeeded in decreasing the objective function after several choices of $r$, new starting values might be needed. Or, a different optimization method might be needed.

\subsection*{12.7.2 Berndt, Hall, Hall, and Hausman Algorithm}
In the context of maximum likelihood estimation, Berndt, Hall, Hall, and Hausman (1974) (hereafter, BHHH ) proposed using the outer product of the score in place of\\
the Hessian. This method can be applied in the general M-estimation case (even though the information matrix equality (12.53) that motivates the method need not hold). The BHHH iteration for a minimization problem is


\begin{equation*}
\boldsymbol{\theta}^{\{g+1\}}=\boldsymbol{\theta}^{\{g\}}-r\left[\sum_{i=1}^{N} \mathbf{s}_{i}\left(\boldsymbol{\theta}^{\{g\}}\right) \mathbf{s}_{i}\left(\boldsymbol{\theta}^{\{g\}}\right)^{\prime}\right]^{-1}\left[\sum_{i=1}^{N} \mathbf{s}_{i}\left(\boldsymbol{\theta}^{\{g\}}\right)\right], \tag{12.84}
\end{equation*}


where $r$ is the step size. (If we want to maximize $\sum_{i=1}^{N} q\left(\mathbf{w}_{i}, \boldsymbol{\theta}\right)$, the minus sign in equation (12.84) should be replaced with a plus sign.) The term multiplying $r$, sometimes called the direction for the next iteration, can be obtained as the $P \times 1$ OLS coefficients from the regression


\begin{equation*}
1 \text { on } \mathbf{s}_{i}\left(\boldsymbol{\theta}^{\{g\}}\right)^{\prime}, \quad i=1,2, \ldots, N . \tag{12.85}
\end{equation*}


The BHHH procedure is easy to implement because it requires computation of the score only; second derivatives are not needed. Further, because the sum of the outer product of the scores is always at least p.s.d., it does not suffer from the potential nonpositive definiteness of the Hessian.

A convenient stopping rule for the BHHH method is obtained as in expression (12.82), but with the sum of the outer products of the score replacing the sum of the Hessians. This is identical to $N$ times the uncentered $R$-squared from regression (12.85). Interestingly, this is the same regression used to obtain the outer product of the score form of the LM statistic when $\mathbf{B}_{\mathrm{o}}=\mathbf{A}_{\mathrm{o}}$, a feature that suggests a natural method for estimating a complicated model after a simpler version of the model has been estimated. Set the starting value, $\boldsymbol{\theta}^{\{0\}}$, equal to the vector of restricted estimates, $\tilde{\boldsymbol{\theta}}$. Then $N R_{0}^{2}$ from the regression used to obtain the first iteration can be used to test the restricted model against the more general model to be estimated; if the restrictions are not rejected, we could just stop the iterations. Of course, as we discussed in Section 12.6.2, the outer-product form of the LM statistic is often ill-behaved even with fairly large sample sizes.

\subsection*{12.7.3 Generalized Gauss-Newton Method}
The final iteration scheme we cover is closely related to the estimator of the expected value of the Hessian in expression (12.44). Let $\mathbf{A}\left(\mathbf{x}, \boldsymbol{\theta}_{\mathrm{o}}\right)$ be the expected value of $\mathbf{H}\left(\mathbf{w}, \boldsymbol{\theta}_{\mathrm{o}}\right)$ conditional on $\mathbf{x}$, where $\mathbf{w}$ is partitioned into $\mathbf{y}$ and $\mathbf{x}$. Then the generalized Gauss-Newton method uses the updating equation


\begin{equation*}
\boldsymbol{\theta}^{\{g+1\}}=\boldsymbol{\theta}^{\{g\}}-r\left[\sum_{i=1}^{N} \mathbf{A}_{i}\left(\boldsymbol{\theta}^{\{g\}}\right)\right]^{-1}\left[\sum_{i=1}^{N} \mathbf{s}_{i}\left(\boldsymbol{\theta}^{\{g\}}\right)\right], \tag{12.86}
\end{equation*}


where $\boldsymbol{\theta}^{\{g\}}$ replaces $\boldsymbol{\theta}_{\mathrm{o}}$ in $\mathbf{A}\left(\mathbf{x}_{i}, \boldsymbol{\theta}_{\mathrm{o}}\right)$. (As before, $\mathbf{A}_{i}$ and $\mathbf{s}_{i}$ might also depend on $\hat{\gamma}$.) This scheme works well when $\mathbf{A}\left(\mathbf{x}, \boldsymbol{\theta}_{\mathrm{o}}\right)$ can be obtained in closed form.

In the special case of nonlinear least squares, we obtain what is traditionally called the Gauss-Newton method (for example, Quandt, 1983). Because $\mathbf{s}_{i}(\boldsymbol{\theta})= -\nabla_{\theta} m_{i}(\boldsymbol{\theta})^{\prime}\left[y_{i}-m_{i}(\boldsymbol{\theta})\right]$, the iteration step is

$$
\boldsymbol{\theta}^{\{g+1\}}=\boldsymbol{\theta}^{\{g\}}+r\left(\sum_{i=1}^{N} \nabla_{\theta} m_{i}^{\{g\} \prime} \nabla_{\theta} m_{i}^{\{g\}}\right)^{-1}\left(\sum_{i=1}^{N} \nabla_{\theta} m_{i}^{\{g\} \prime} u_{i}^{\{g\}}\right)
$$

The term multiplying the step size $r$ is obtained as the OLS coefficients of the regression of the residuals on the gradient, both evaluated at $\boldsymbol{\theta}^{\{g\}}$. The stopping rule can be based on $N$ times the uncentered $R$-squared from this regression. Note how closely the Gauss-Newton method of optimization is related to the regression used to obtain the nonrobust LM statistic (see regression (12.72)).

\subsection*{12.7.4 Concentrating Parameters out of the Objective Function}
In some cases, it is computationally convenient to concentrate one set of parameters out of the objective function. Partition $\boldsymbol{\theta}$ into the vectors $\boldsymbol{\beta}$ and $\boldsymbol{\gamma}$. Then the first-order conditions that define $\hat{\boldsymbol{\theta}}$ are


\begin{equation*}
\sum_{i=1}^{N} \nabla_{\beta} q\left(\mathbf{w}_{i}, \boldsymbol{\beta}, \boldsymbol{\gamma}\right)=\mathbf{0}, \quad \sum_{i=1}^{N} \nabla_{\gamma} q\left(\mathbf{w}_{i}, \boldsymbol{\beta}, \boldsymbol{\gamma}\right)=\mathbf{0} . \tag{12.87}
\end{equation*}


Rather than solving these for $\hat{\boldsymbol{\beta}}$ and $\hat{\gamma}$, suppose that the second set of equations can be solved for $\boldsymbol{\gamma}$ as a function of $\mathbf{W} \equiv\left(\mathbf{w}_{1}, \mathbf{w}_{2}, \ldots, \mathbf{w}_{N}\right)$ and $\boldsymbol{\beta}$ for any outcomes $\mathbf{W}$ and any $\boldsymbol{\beta}$ in the parameter space: $\boldsymbol{\gamma}=\mathbf{g}(\mathbf{W}, \boldsymbol{\beta})$. Then, by construction,


\begin{equation*}
\sum_{i=1}^{N} \nabla_{\gamma} q\left[\mathbf{w}_{i}, \boldsymbol{\beta}, \mathbf{g}(\mathbf{W}, \boldsymbol{\beta})\right]=\mathbf{0} \tag{12.88}
\end{equation*}


When we plug $\boldsymbol{g}(\mathbf{W}, \boldsymbol{\beta})$ into the original objective function, we obtain the concentrated objective function,


\begin{equation*}
Q^{c}(\mathbf{W}, \boldsymbol{\beta})=\sum_{i=1}^{N} q\left[\mathbf{w}_{i}, \boldsymbol{\beta}, \mathbf{g}(\mathbf{W}, \boldsymbol{\beta})\right] \tag{12.89}
\end{equation*}


Under standard differentiability assumptions, the minimizer of equation (12.89) is identical to the $\hat{\boldsymbol{\beta}}$ that solves equations (12.87) (along with $\hat{\boldsymbol{\gamma}}$ ), as can be seen by differentiating equation (12.89) with respect to $\boldsymbol{\beta}$ using the chain rule, setting the result to zero, and using equation (12.88); then $\hat{\gamma}$ can be obtained as $\mathbf{g}(\mathbf{W}, \hat{\boldsymbol{\beta}})$.

As a device for studying asymptotic properties, the concentrated objective function is of limited value because $\mathbf{g}(\mathbf{W}, \boldsymbol{\beta})$ generally depends on all of $\mathbf{W}$, in which case the objective function cannot be written as the sum of independent, identically distributed summands. One setting where equation (12.89) is a sum of i.i.d. functions occurs when we concentrate out individual-specific effects from certain nonlinear panel data models. In addition, the concentrated objective function can be useful for establishing the equivalence of seemingly different estimation approaches.

\subsection*{12.8 Simulation and Resampling Methods}
So far we have focused on the asymptotic properties of M-estimators, as these provide a unified framework for inference. But there are a few good reasons to go beyond asymptotic results, at least in some cases. First, the asymptotic approximations need not be very good, especially with small sample sizes, highly nonlinear models, or unusual features of the population distribution of $\mathbf{w}_{i}$. Simulation methods, while always special, can help determine how well the asymptotic approximations work. Resampling methods can allow us to improve on the asymptotic distribution approximations.

Even if we feel comfortable with asymptotic approximations to the distribution of $\hat{\boldsymbol{\theta}}$, we may not be as confident in the approximations for estimating a nonlinear function of the parameters, say $\boldsymbol{\gamma}_{\mathrm{o}}=\mathbf{g}\left(\boldsymbol{\theta}_{\mathrm{o}}\right)$. Under the assumptions in Section 3.5.2, we can use the delta method to approximate the variance of $\hat{\gamma}=\mathbf{g}(\hat{\boldsymbol{\theta}})$. Depending on the nature of $\mathbf{g}(\cdot)$, applying the delta method might be difficult, and it might not result in a very good approximation. Resampling methods can simplify the calculation of standard errors, confidence intervals, and $p$-values for test statistics, and we can get a good idea of the amount of finite-sample bias in the estimation method. In addition, under certain assumptions and for certain statistics, resampling methods can provide quantifiable improvements to the usual asymptotics.

\subsection*{12.8.1 Monte Carlo Simulation}
In a Monte Carlo simulation, we attempt to estimate the mean and varianceassuming that these exist-and possibly other features of the distribution of the Mestimator, $\hat{\boldsymbol{\theta}}$. The idea is usually to determine how much bias $\hat{\boldsymbol{\theta}}$ has for estimating $\boldsymbol{\theta}_{\mathrm{o}}$, or to determine the efficiency of $\hat{\boldsymbol{\theta}}$ compared with other estimators of $\boldsymbol{\theta}_{\mathrm{o}}$. In addition, we often want to know how well the asymptotic standard errors approximate the standard deviations of the $\hat{\theta}_{j}$.

To conduct a simulation, we must choose a population distribution for $\mathbf{w}$, which depends on the finite-dimensional vector $\boldsymbol{\theta}_{\mathrm{o}}$. We must set the values of $\boldsymbol{\theta}_{\mathrm{o}}$, and decide\\
on a sample size, $N$. We then draw a random sample of size $N$ from this distribution and use the sample to obtain an estimate of $\boldsymbol{\theta}_{\mathrm{o}}$. We draw a new random sample and compute another estimate of $\boldsymbol{\theta}_{\mathrm{o}}$. We repeat the process for several iterations, say $M$. Let $\hat{\boldsymbol{\theta}}^{(m)}$ be the estimate of $\boldsymbol{\theta}_{\mathrm{o}}$ based on the $m$ th iteration. Given $\left\{\hat{\boldsymbol{\theta}}^{(m)}: m=\right. 1,2, \ldots, M\}$, we can compute the sample average and sample variance to estimate $\mathrm{E}(\hat{\boldsymbol{\theta}})$ and $\operatorname{Var}(\hat{\boldsymbol{\theta}})$, respectively. We might also form $t$ statistics or other test statistics to see how well the asymptotic distributions approximate the finite-sample distributions. We can also see how well asymptotic confidence intervals cover the population parameter relative to the nominal confidence level.

A good Monte Carlo study varies the value of $\boldsymbol{\theta}_{\mathrm{o}}$, the sample size, and even the general form of the distribution of $\mathbf{w}$. In fact, it is generally a good idea to check how an estimation method fares when the assumptions on which it is based partially or completely fail. Obtaining a thorough study can be very challenging, especially for a complicated, nonlinear model. First, to get good estimates of the distribution of $\hat{\boldsymbol{\theta}}$, we would like $M$ to be large (perhaps several thousand). But for each Monte Carlo iteration, we must obtain $\hat{\boldsymbol{\theta}}^{(m)}$, and this step can be computationally expensive because it often requires the iterative methods we discussed in Section 12.7. Repeating the simulations for many different sample sizes $N$, values of $\boldsymbol{\theta}_{\mathrm{o}}$, and distributional shapes can be very time-consuming.

In most economic applications, $\mathbf{w}_{i}$ is partitioned as $\left(\mathbf{x}_{i}, \mathbf{y}_{i}\right)$. While we can draw the full vector $\mathbf{w}_{i}$ randomly in the Monte Carlo iterations, sometimes the $\mathbf{x}_{i}$ are fixed at the beginning of the iterations, and then $\mathbf{y}_{i}$ is drawn from the conditional distribution given $\mathbf{x}_{i}$. This method simplifies the simulations because we do not need to vary the distribution of $\mathbf{x}_{i}$ along with the distribution of interest, the distribution of $\mathbf{y}_{i}$ given $\mathbf{x}_{i}$. If we fix the $\mathbf{x}_{i}$ at the beginning of the simulations, the distributional features of $\hat{\boldsymbol{\theta}}$ that we estimate from the Monte Carlo simulations are conditional on $\left\{\mathbf{x}_{1}, \mathbf{x}_{2}, \ldots, \mathbf{x}_{N}\right\}$. This conditional approach is especially common in linear and nonlinear regression contexts, as well as conditional maximum likelihood.

It is important not to rely too much on Monte Carlo simulations. Many estimation methods, including OLS, IV, and panel data estimators, have asymptotic properties that do not depend on underlying distributions. In the nonlinear regression model, the NLS estimator is $\sqrt{N}$-asymptotically normal, and the usual asymptotic variance matrix (12.58) is valid under Assumptions NLS.1-NLS.3. However, in a typical Monte Carlo simulation, the implied error, $u$, is assumed to be independent of $\mathbf{x}$, and the distribution of $u$ must be specified. The Monte Carlo results then pertain to this distribution, and it can be misleading to extrapolate to different settings. In addition, we can never try more than just a small part of the parameter space. Because we never know the population value $\boldsymbol{\theta}_{\mathrm{o}}$, we can never be sure how well our Monte Carlo\\
study describes the underlying population. Hendry (1984) discusses how response surface analysis can be used to reduce the specificity of Monte Carlo studies. See also Davidson and MacKinnon (1993, Chap. 21).

\subsection*{12.8.2 Bootstrapping}
A Monte Carlo simulation, although it is informative about how well the asymptotic approximations can be expected to work in specific situations, does not generally help us refine our inference given a particular sample. (Because we do not know $\boldsymbol{\theta}_{\mathrm{o}}$, we cannot know whether our Monte Carlo findings apply to the population we are studying. Nevertheless, researchers sometimes use the results of a Monte Carlo simulation to obtain rules of thumb for adjusting standard errors or for adjusting critical values for test statistics.) The method of bootstrapping, which is a popular resampling method, can be used as an alternative to asymptotic approximations for obtaining standard errors, confidence intervals, and $p$-values for test statistics.

Though there are several variants of the bootstrap, we begin with one that can be applied to general M-estimation. The goal is to approximate the distribution of $\hat{\boldsymbol{\theta}}$ without relying on the usual first-order asymptotic theory. Let $\left\{\mathbf{w}_{1}, \mathbf{w}_{2}, \ldots, \mathbf{w}_{N}\right\}$ denote the outcome of the random sample used to obtain the estimate. The nonparametric bootstrap is essentially a Monte Carlo simulation where the observed sample is treated as the population. In other words, at each bootstrap iteration, $b$, a random sample of size $N$ is drawn from $\left\{\mathbf{w}_{1}, \mathbf{w}_{2}, \ldots, \mathbf{w}_{N}\right\}$. (That is, we sample with replacement.) In practice, we use a random number generator to obtain $N$ integers from the set $\{1,2, \ldots, N\}$; in the vast majority of iterations some integers will be repeated at least once. These integers index the elements that we draw from $\left\{\mathbf{w}_{1}, \mathbf{w}_{2}, \ldots, \mathbf{w}_{N}\right\}$; call these $\left\{\mathbf{w}_{1}^{(b)}, \mathbf{w}_{2}^{(b)}, \ldots, \mathbf{w}_{N}^{(b)}\right\}$. Next, we use this bootstrap sample to obtain the M-estimate $\hat{\boldsymbol{\theta}}^{(b)}$ by solving\\
$\min _{\boldsymbol{\theta} \in \boldsymbol{\Theta}} \sum_{i=1}^{N} q\left(\mathbf{w}_{i}^{(b)}, \boldsymbol{\theta}\right)$.\\
We iterate the process $B$ times, obtaining $\hat{\boldsymbol{\theta}}^{(b)}, b=1, \ldots, B$. These estimates can now be used as in a Monte Carlo simulation. Computing the average of the $\hat{\boldsymbol{\theta}}^{(b)}$, say $\overline{\hat{\boldsymbol{\theta}}}$, allows us to estimate the bias in $\hat{\boldsymbol{\theta}}$, called the bootstrap bias estimate. The sample variance, $(B-1)^{-1} \sum_{b=1}^{B}\left[\hat{\boldsymbol{\theta}}^{(b)}-\overline{\hat{\boldsymbol{\theta}}}\right] \cdot\left[\hat{\boldsymbol{\theta}}^{(b)}-\overline{\hat{\boldsymbol{\theta}}}\right]^{\prime}$, called the bootstrap variance estimate, can be used to obtain standard errors for the $\hat{\theta}_{j}$-the estimates from the original sample. For a scalar estimate $\hat{\gamma}$, we obtain its bootstrap standard error as $\operatorname{se}_{B}(\hat{\gamma})= \left[(B-1)^{-1} \sum_{b=1}^{B}\left(\hat{\gamma}^{(b)}-\overline{\hat{\gamma}}\right)\right]^{1 / 2}$. Naturally, we can apply this formula to smooth functions of the original parameter estimates $\hat{\boldsymbol{\theta}}$, say $\hat{\gamma}=g(\hat{\boldsymbol{\theta}})$ for a continuously differ-\\
entiable function $g: \mathbb{R}^{P} \rightarrow \mathbb{R}$. We can then use $\operatorname{se}_{B}(\hat{\gamma})$ to construct asymptotic hypotheses tests and confidence intervals for $\gamma_{\mathrm{o}}$.

Especially for computing average partial effects-a topic that will arise repeatedly in Part IV-we often need to estimate a parameter that can be written as $\gamma_{\mathrm{o}}= \mathrm{E}\left[g\left(\mathbf{w}_{i}, \boldsymbol{\theta}_{\mathrm{o}}\right)\right]$. A natural, consistent estimator is $\hat{\gamma}=N^{-1} \sum_{i=1}^{N} g\left(\mathbf{w}_{i}, \hat{\boldsymbol{\theta}}\right)$. To estimate its asymptotic variance, we must account for the randomness in $\mathbf{w}_{i}$ as well as $\hat{\boldsymbol{\theta}}$. As before, we draw bootstrap samples, and, for bootstrap sample $b$, the estimate of $\gamma_{\mathrm{o}}$ is

$$
\hat{\gamma}^{(b)}=N^{-1} \sum_{i=1}^{N} g\left(\mathbf{w}_{i}^{(b)}, \hat{\boldsymbol{\theta}}^{(b)}\right)
$$

Once we have the collection of bootstrap estimates $\left\{\hat{\gamma}^{(b)}: b=1,2, \ldots, B\right\}$, we can compute the bootstrap bias and, more important, the bootstrap standard error using the previous formulas.

Using the bootstrap standard error to construct statistics and confidence intervals is often much easier than the analytical calculations needed to obtain an asymptotic standard error based on first-order asymptotics. But using the bootstrap standard error to construct test statistics cannot be shown to improve on the approximation provided by the usual asymptotic theory. As it turns out, in many cases the bootstrap does improve the approximation of the distribution of test statistics. In other words, the bootstrap can provide an asymptotic refinement compared with the usual asymptotic theory, but one must use some care in computing the bootstrap test statistics.

To show that the bootstrap approximation of a distribution converges more quickly than the usual rates associated with first-order asymptotics, the notion of an asymptotically pivotal statistic is critical. An asymptotically pivotal statistic is one whose limiting distribution does not depend on unknown parameters. Asymptotic $t$ statistics, Wald statistics, score statistics, and quasi-LR statistics are all asymptotically pivotal when they converge to the standard normal distribution, in the case of a $t$ statistic, and to the chi-square distribution, in the case of the other statistics. One must sometimes use care, though, to ensure a statistic is asymptotically pivotal. For example, for a $t$ statistic to be asymptotically pivotal in the context of nonlinear regression with heteroskedasticity, we must use a heteroskedasticity-robust statistic. The Wald and score statistics must often use robust asymptotic variance estimators to deliver an asymptotic chi-square distribution. The quasi-LR statistic is guaranteed to be asymptotically pivotal only when the generalized information matrix equality holds.

To explain how to bootstrap the critical values for a $t$ statistic, consider testing $\mathrm{H}_{0}: \theta_{\mathrm{o}}=c$ for some known value $c$. The $t$ statistic, $t=(\hat{\theta}-c) / \operatorname{se}(\hat{\theta})$, is asymptotically\\
pivotal if $\operatorname{se}(\hat{\theta})$ is appropriately chosen. To obtain a refinement using the bootstrap, we must obtain the empirical distribution of the statistic\\
$t^{(b)}=\left(\hat{\theta}^{(b)}-\hat{\theta}\right) / \operatorname{se}\left(\hat{\theta}^{(b)}\right)$,\\
where $\hat{\theta}$ is the estimate from the original sample, $\hat{\theta}^{(b)}$ is the estimate for bootstrap sample $b$, and $\operatorname{se}\left(\hat{\theta}^{(b)}\right)$ is the standard error estimated from the same bootstrap sample. (So, for example, se $\left(\hat{\theta}^{(b)}\right)$ could be a heteroskedasticity-robust standard error for NLS.) Notice how the $t$ statistic for each bootstrap replication is centered at the original estimate, $\hat{\theta}$, not the hypothesized value. As discussed by Horowitz (2001), centering at the estimate is required to ensure asymptotic refinements of the testing procedure.

The way we obtain bootstrap critical values for a test depends on the nature of the alternative. For a one-sided alternative, say $\mathrm{H}_{0}: \theta_{\mathrm{o}}>c$, we order the statistics $\left\{t^{(b)}: b=1,2, \ldots, B\right\}$, from smallest to largest, and we pick the value representing the desired quantile of the list of ordered values. For example, to obtain a 5 percent test against a greater than one-sided alternative, we choose the critical value as the 95th percentile in the ordered list of $t^{(b)}$. For a two-sided alternative, we must choose between a nonsymmetric test and a symmetric test. For the former, a test with size $\alpha$ chooses critical values as the lower and upper $\alpha / 2$ quantiles of the ordered bootstrap test statistics, and we reject $\mathrm{H}_{0}$ if $t>c v_{u}$ or $t<c v_{l}$. For the latter, we first order the absolute values of the statistics, $\left|t^{(b)}\right|$, and then choose the upper $\alpha$ quantile as the critical value for a test of size $\alpha$. Naturally, we compare $|t|$ with the critical value. This approach to choosing critical values from bootstrapping is called the percentile-t method.

We can use the percentile- $t$ method to compute a bootstrap $\boldsymbol{p}$-value. For example, against a greater than one-sided alternative, we simply find the fraction of bootstrap $t$ statistics $t^{(b)}$ that exceed $t$. A symmetric $p$-value for a two-sided alternative does the same for $\left|t^{(b)}\right|$ and $|t|$.

Testing multiple hypotheses is similar. Suppose that for a $Q$-vector $\boldsymbol{\phi}_{\mathrm{o}}$, we want to test $\mathrm{H}_{\mathrm{O}}: \boldsymbol{\phi}_{\mathrm{O}}=\mathbf{r}$, where $\mathbf{r}$ is a vector of known constants. The Wald statistic computed using the original sample is $W=(\hat{\boldsymbol{\phi}}-\mathbf{r})^{\prime} \hat{\mathbf{V}}^{-1}(\hat{\boldsymbol{\phi}}-\mathbf{r})$. We compute a series of Wald statistics from bootstrap samples as

$$
W^{(b)}=\left(\hat{\boldsymbol{\phi}}^{(b)}-\hat{\boldsymbol{\phi}}\right)^{\prime}\left(\hat{\mathbf{V}}^{(b)}\right)^{-1}\left(\hat{\boldsymbol{\phi}}^{(b)}-\hat{\boldsymbol{\phi}}\right), \quad b=1, \ldots, B
$$

where we must take care so that the calculation of $\hat{\mathbf{V}}$ (and $\hat{\mathbf{V}}^{(b)}$ ) delivers an asymptotic chi-square statistic. The bootstrap $p$-value is the fraction of $W^{(b)}$ that exceed $W$.

The parametric bootstrap is even more similar to a standard Monte Carlo simulation because we assume that the distribution of $\mathbf{w}$ is known up to the parameters $\boldsymbol{\theta}_{\mathrm{o}}$.

Let $f(\cdot, \boldsymbol{\theta})$ denote the parametric density. Then, on each bootstrap iteration, we draw a random sample of size $N$ from $f(\cdot, \hat{\boldsymbol{\theta}})$; this gives $\left\{\mathbf{w}_{1}^{(b)}, \mathbf{w}_{2}^{(b)}, \ldots, \mathbf{w}_{N}^{(b)}\right\}$, and the rest of the calculations are the same as in the nonparametric bootstrap. (With the parametric bootstrap, when $f(\cdot, \boldsymbol{\theta})$ is a continuous density, only rarely would we find repeated values among the $\mathbf{w}_{i}^{(b)}$.)

When $\mathbf{w}_{i}$ is partitioned into $\left(\mathbf{x}_{i}, \mathbf{y}_{i}\right)$, where the $\mathbf{x}_{i}$ are conditioning variables, other resampling schemes are sometimes preferred. In Chapter 13, we study the method of conditional maximum likelihood, where we assume that a model of a conditional density, $f(\mathbf{y} \mid \mathbf{x} ; \boldsymbol{\theta})$, is correctly specified. In that case, we can apply a combination of the nonparametric and parametric bootstrap. Because the distribution of $\mathbf{x}_{i}$ is unspecified, we randomly draw $N$ indexes from $\{1,2, \ldots, N\}$ to obtain a nonparametric bootstrap sample for the conditioning variables, $\left\{\mathbf{x}_{i}^{(b)}: i=1, \ldots, N\right\}$. Then, given the estimate $\hat{\boldsymbol{\theta}}$, we obtain $\mathbf{y}_{i}^{(b)}$ by drawing from the density $f\left(\cdot \mid \mathbf{x}_{i}^{(b)} ; \hat{\boldsymbol{\theta}}\right)$. We then use the bootstrap samples $\left(\mathbf{x}_{i}^{(b)}, \mathbf{y}_{i}^{(b)}\right)$ as before. Compared with the fully nonparametric bootstrap where we resample the entire vector $\mathbf{w}_{i}=\left(\mathbf{x}_{i}, \mathbf{y}_{i}\right)$ from the original data, the method that draws from $f\left(\cdot \mid \mathbf{x}_{i}^{(b)} ; \hat{\boldsymbol{\theta}}\right)$ is not as widely applicable and is computationally more expensive.

In cases where we have even more structure, other alternatives are available. For example, in a nonlinear regression model with $y_{i}=m\left(\mathbf{x}_{i}, \boldsymbol{\theta}_{\mathrm{o}}\right)+u_{i}$, where the error $u_{i}$ is independent of $\mathbf{x}_{i}$, we first compute the NLS estimate $\hat{\boldsymbol{\theta}}$ and the NLS residuals, $\hat{u}_{i}=y_{i}-m\left(\mathbf{x}_{i}, \hat{\boldsymbol{\theta}}\right), i=1,2, \ldots, N$. Then, using the procedure described for the nonparametric bootstrap, a bootstrap sample of residuals, $\left\{\hat{u}_{i}^{(b)}: i=1,2, \ldots, N\right\}$, is obtained, and we compute $y_{i}^{(b)}=m\left(\mathbf{x}_{i}, \hat{\boldsymbol{\theta}}\right)+\hat{u}_{i}^{(b)}$. Using the generated data $\left\{\left(\mathbf{x}_{i}, y_{i}^{(b)}\right): i=1,2, \ldots, N\right\}$, we compute the NLS estimate, $\hat{\boldsymbol{\theta}}^{(b)}$. This procedure is called the nonparametric residual bootstrap. (We resample the residuals and use these to generate a sample on the dependent variable, but we do not resample the conditioning variables, $\mathbf{x}_{i}$.) If the model is nonlinear in $\boldsymbol{\theta}$, this method can be computationally demanding because we want $B$ to be several hundred, if not several thousand. Nonetheless, such procedures are becoming more and more feasible as computational speed increases. When $u_{i}$ has zero conditional mean $\left[\mathrm{E}\left(u_{i} \mid \mathbf{x}_{i}\right)=0\right]$ but is heteroskedastic $\left[\operatorname{Var}\left(u_{i} \mid \mathbf{x}_{i}\right)\right.$ depends on $\left.\mathbf{x}_{i}\right]$, alternative sampling methods, in particular the wild bootstrap, can be used to obtain heteroskedastic-consistent standard errors. See, for example, Horowitz (2001).

Bootstrapping is easily applied to the kinds of panel data structures that we treat in this text because our sampling assumption for panel data is random sampling in the cross section dimension. For panel data, we simply let $\mathbf{w}_{i}=\left(\mathbf{w}_{i 1}, \ldots, \mathbf{w}_{i T}\right)$ denote the outcomes across all $T$ time periods for cross section observation $i$. When we obtain a bootstrap sample, all time periods for a particular unit constitute a single\\
observation. In other words, as with pure cross section applications, we randomly select $N$ integers (with replacement) from $\{1,2, \ldots, N\}$ and obtain the bootstrap sample $\mathbf{w}_{i}^{(b)}=\left(\mathbf{w}_{i 1}^{(b)}, \ldots, \mathbf{w}_{i T}^{(b)}\right)$. For emphasis: we do not resample separate time periods within a unit; we only resample units. As with the first-order asymptotic theory based on random samples, this kind of bootstrapping for panel data is realistic for large cross sections and relatively small time periods.

\subsection*{12.9 Multivariate Nonlinear Regression Methods}
As with linear regression methods, nonlinear regression methods can be extended to systems of equations. For example, suppose that $y_{1}$ and $y_{2}$ are fractional variablessay, shares of pension investments put into stocks and bonds, respectively, with a third "other" category-and we wish to account for the bounded nature of these responses in our model. We might specify logistic regression functions of the form\\
$\mathrm{E}\left(y_{g} \mid \mathbf{x}\right)=\exp \left(\mathbf{x} \boldsymbol{\theta}_{g}\right) /\left[1+\exp \left(\mathbf{x} \boldsymbol{\theta}_{g}\right)\right], \quad g=1,2$,\\
where $\mathbf{x}$ is a row vector of common explanatory variables. We can estimate $\boldsymbol{\theta}_{1}$ and $\boldsymbol{\theta}_{2}$ individually, using NLS or WNLS, but we can also estimate them jointly. Not surprisingly, under certain assumptions, a multivariate procedure that accounts for correlation in the unobservables across equations is more efficient than single equation estimation, as we discuss in Section 12.9.2.

We can also apply nonlinear regression to panel data structures. For example, for a nonnegative response variable $y_{t}$, a dynamic model is\\
$\mathrm{E}\left(y_{t} \mid \mathbf{z}_{t}, y_{t-1}\right)=\exp \left(\mathbf{z}_{t} \boldsymbol{\beta}+\alpha y_{t-1}\right), \quad t=1, \ldots, T$.\\
One way to estimate the parameters is by pooled NLS. As we show in the next two subsections, we can cast both system and panel data problems in one framework.

\subsection*{12.9.1 Multivariate Nonlinear Least Squares}
A general setup for multivariate nonlinear least squares (MNLS) is to let $y_{g}$ be the response variable for equation $g$, which could be a time period. The corresponding explanatory variables are $\mathbf{x}_{g}$. We assume the model for $\mathrm{E}\left(y_{g} \mid \mathbf{x}_{g}\right)$ is correctly specified:


\begin{equation*}
\mathrm{E}\left(y_{g} \mid \mathbf{x}_{g}\right)=m_{g}\left(\mathbf{x}_{g}, \boldsymbol{\theta}_{o g}\right), \quad g=1, \ldots, G, \tag{12.90}
\end{equation*}


where the parameters in different equations can be distinct or there can be restrictions across $g$. Given $N$ randomly sampled observations, it should be no surprise that the $\boldsymbol{\theta}_{o g}$ can be consistently estimated by solving

$$
\min _{\boldsymbol{\theta}} \sum_{i=1}^{N} \sum_{g=1}^{G}\left[y_{i g}-m_{g}\left(\mathbf{x}_{i g}, \boldsymbol{\theta}_{g}\right)\right]^{2}
$$

or, in vector form,


\begin{equation*}
\min _{\boldsymbol{\theta}} \sum_{i=1}^{N}\left[\mathbf{y}_{i}-\mathbf{m}\left(\mathbf{x}_{i}, \boldsymbol{\theta}\right)\right]^{\prime}\left[\mathbf{y}_{i}-\mathbf{m}\left(\mathbf{x}_{i}, \boldsymbol{\theta}\right)\right], \tag{12.91}
\end{equation*}


where $\boldsymbol{\theta}$ denotes the vector of all parameters, $\mathbf{y}_{i}$ is the $G \times 1$ vector of responses for observation $i$, and $\mathbf{m}\left(\mathbf{x}_{i}, \boldsymbol{\theta}\right)$ is the $G \times 1$ vector of conditional mean functions. We call the solution to (12.91) the multivariate nonlinear least squares estimator. As with univariate NLS, identification requires that $\boldsymbol{\theta}_{\mathrm{o}}$ is the only solution to the corresponding population problem, and we would assume mean functions twice continuously differentiable in the parameters. Generally, the assumptions are as weak as in the univariate case.

If there are $G$ separate parameter vectors without restrictions across $g$, then the solutions to (12.91) are the same as NLS on each equation. Sometimes, for example with cost share equations derived from the theory of the firm, the parameters will be restricted across equations. Then the MNLS estimator can be used to estimate the parameters with the restrictions imposed.

In the panel data case, where the conditional mean model is written for a common set of parameters,


\begin{equation*}
\mathrm{E}\left(y_{i t} \mid \mathbf{x}_{i t}\right)=m\left(\mathbf{x}_{i t}, \boldsymbol{\theta}_{\mathrm{o}}\right), \tag{12.92}
\end{equation*}


the MNLS estimator is the pooled nonlinear least squares (PNLS) estimator. Problem 12.6 asks you to study this estimator in more detail. Consistency and $\sqrt{N}$-asymptotic normality of the PNLS estimator follows under general assumptions on the smoothness of the mean function, but generally one should use a fully robust variance matrix estimator to account for possible heteroskedasticity and serial correlation. If the conditional mean model is dynamically complete (see Problem 12.6), then one need not worry about serial correlation in the errors when estimating the variance matrix of the PNLS estimator. But heteroskedasticity could still be an issue. If one adds the homoskedasticity assumption $\operatorname{Var}\left(y_{i t} \mid \mathbf{x}_{i t}\right)=\sigma_{\mathrm{o}}^{2}$ to the appropriate no serial correlation assumption, the usual statistics from the pooled NLS analysis are asymptotically valid.

The PNLS estimator is attractive when one wishes only to impose (12.92), without making the stronger assumption that the covariates are strictly exogenous. Even\\
under strict exogeneity, PNLS is usually needed as a first step in obtaining an asymptotically more efficient estimator. We study that possibility next.

\subsection*{12.9.2 Weighted Multivariate Nonlinear Least Squares}
Under certain assumptions, we can use generalized least squares (GLS) methods to more efficiently estimate the parameters appearing in a set of conditional mean functions. Here we focus on the case where the explanatory variables are strictly exogenous. This often applies to systems of equations and sometimes applies to panel data methods. Specifically, we assume\\
$\mathrm{E}\left(\mathbf{y}_{i} \mid \mathbf{x}_{i}\right)=\mathbf{m}\left(\mathbf{x}_{i}, \boldsymbol{\theta}_{\mathrm{o}}\right), \quad$ some $\boldsymbol{\theta}_{\mathrm{o}} \in \boldsymbol{\Theta} \subset \mathbb{R}^{P}$,\\
where again, $\mathbf{y}_{i}$ is a $G \times 1$ vector on the dependent variable and $\mathbf{m}\left(\mathbf{x}_{i}, \boldsymbol{\theta}\right)$ is a $G \times 1$ vector of conditional mean functions. It is important to note that the entire vector of covariates, $\mathbf{x}_{i}$, is conditioned on in (12.93). This is the sense in which the explanatory variables are strictly exogenous: if some elements of $\mathbf{x}_{i}$ are omitted from one of the functions $m_{g}\left(\mathbf{x}_{i}, \boldsymbol{\theta}_{\mathrm{o}}\right)$, then they should have no partial effect on $\mathrm{E}\left(y_{i g} \mid \mathbf{x}_{i}\right)$. As we discussed in Chapter 7, this is often intended in seemingly unrelated regressions (SUR)type applications; in fact, often the same set of explanatory variables appears in each equation. But assumption (12.93) is violated for panel data models with lagged dependent variables and, as we saw in Chapter 7 for linear models, perhaps in models without lagged dependent variables. Nevertheless, we have seen linear models where strict exogeneity holds after including the time average of the covariates, and we will see this again in later chapters-including Chapters 15, 16, 17, and 18 for nonlinear models.

The most general estimator we consider here is the weighted multivariate nonlinear least squares (WMNLS) estimator. We can easily motivate WMNLS. Let W $\left(\mathbf{x}_{i}, \gamma\right)$ be a model for the $G \times G$ conditional variance matrix $\operatorname{Var}\left(\mathbf{y}_{i} \mid \mathbf{x}_{i}\right)$. If this model is correctly specified, then, generally, we can obtain a $\sqrt{N}$-consistent estimator of the "true" parameters in the variance matrix, $\gamma_{\mathrm{o}}$. In fact, we would probably use the residuals from an initial MNLS estimation (more on this shortly). Given $\hat{\gamma}$, and assuming that $\mathbf{W}\left(\mathbf{x}_{i}, \hat{\gamma}\right)$ is nonsingular for all $i$ (often true by construction), we can estimate $\boldsymbol{\theta}_{\mathrm{O}}$ by solving


\begin{equation*}
\min _{\boldsymbol{\theta}} \sum_{i=1}^{N}\left[\mathbf{y}_{i}-\mathbf{m}\left(\mathbf{x}_{i}, \boldsymbol{\theta}\right)\right]^{\prime}\left[\mathbf{W}\left(\mathbf{x}_{i}, \hat{\gamma}\right)\right]^{-1}\left[\mathbf{y}_{i}-\mathbf{m}\left(\mathbf{x}_{i}, \boldsymbol{\theta}\right)\right] . \tag{12.94}
\end{equation*}


The solution to (12.94) is the WMNLS estimator.

As we will see in later chapters, the WMNLS estimator can be attractive for a broad class of nonlinear models, particularly when we cannot, or do not wish to, specify a complete distribution $\mathrm{D}\left(\mathbf{y}_{i} \mid \mathbf{x}_{i}\right)$. Of importance, under (12.93), the WMNLS estimator is generally consistent for $\boldsymbol{\theta}_{\mathrm{O}}$ even if the inverse of the weighting matrix is misspecified for $\boldsymbol{\theta}_{\mathrm{o}}$. The intuition for the robustness of the WMNLS estimator is straightforward from the general M -estimation theory. Let $\gamma^{*}$ be the probability limit of $\hat{\gamma}$, whether or not the variance is misspecified. Then the WMNLS estimator is consistent for $\boldsymbol{\theta}_{\mathrm{O}}$ if $\boldsymbol{\theta}_{\mathrm{O}}$ uniquely solves the population problem

$$
\min _{\boldsymbol{\theta}} \mathrm{E}\left\{\left[\mathbf{y}_{i}-\mathbf{m}\left(\mathbf{x}_{i}, \boldsymbol{\theta}\right)\right]^{\prime}\left[\mathbf{W}\left(\mathbf{x}_{i}, \gamma^{*}\right)\right]^{-1}\left[\mathbf{y}_{i}-\mathbf{m}\left(\mathbf{x}_{i}, \boldsymbol{\theta}\right)\right]\right\} .
$$

But straightforward algebra shows that

$$
\begin{aligned}
{\left[\mathbf{y}_{i}-\right.} & \left.\mathbf{m}\left(\mathbf{x}_{i}, \boldsymbol{\theta}\right)\right]^{\prime}\left[\mathbf{W}\left(\mathbf{x}_{i}, \gamma^{*}\right)\right]^{-1}\left[\mathbf{y}_{i}-\mathbf{m}\left(\mathbf{x}_{i}, \boldsymbol{\theta}\right)\right] \\
= & {\left[\mathbf{y}_{i}-\mathbf{m}\left(\mathbf{x}_{i}, \boldsymbol{\theta}_{\mathrm{o}}\right)\right]^{\prime}\left[\mathbf{W}\left(\mathbf{x}_{i}, \gamma^{*}\right)\right]^{-1}\left[\mathbf{y}_{i}-\mathbf{m}\left(\mathbf{x}_{i}, \boldsymbol{\theta}_{\mathrm{o}}\right)\right] } \\
& -2\left[\mathbf{y}_{i}-\mathbf{m}\left(\mathbf{x}_{i}, \boldsymbol{\theta}_{\mathrm{o}}\right)\right]^{\prime}\left[\mathbf{W}\left(\mathbf{x}_{i}, \gamma^{*}\right)\right]^{-1}\left[\mathbf{m}\left(\mathbf{x}_{i}, \boldsymbol{\theta}_{\mathrm{o}}\right)-\mathbf{m}\left(\mathbf{x}_{i}, \boldsymbol{\theta}\right)\right] \\
& +\left[\mathbf{m}\left(\mathbf{x}_{i}, \boldsymbol{\theta}_{\mathrm{o}}\right)-\mathbf{m}\left(\mathbf{x}_{i}, \boldsymbol{\theta}\right)\right]^{\prime}\left[\mathbf{W}\left(\mathbf{x}_{i}, \gamma^{*}\right)\right]^{-1}\left[\mathbf{m}\left(\mathbf{x}_{i}, \boldsymbol{\theta}_{\mathrm{o}}\right)-\mathbf{m}\left(\mathbf{x}_{i}, \boldsymbol{\theta}\right)\right]
\end{aligned}
$$

By the law of interated expectations, the term in the middle has zero conditional mean by (12.93), and so


\begin{align*}
\mathrm{E}\left\{\left[\mathbf{y}_{i}-\right.\right. & \left.\left.\mathbf{m}\left(\mathbf{x}_{i}, \boldsymbol{\theta}\right)\right]^{\prime}\left[\mathbf{W}\left(\mathbf{x}_{i}, \gamma^{*}\right)\right]^{-1}\left[\mathbf{y}_{i}-\mathbf{m}\left(\mathbf{x}_{i}, \boldsymbol{\theta}\right)\right]\right\} \\
= & \mathrm{E}\left\{\left[\mathbf{y}_{i}-\mathbf{m}\left(\mathbf{x}_{i}, \boldsymbol{\theta}_{\mathrm{o}}\right)\right]^{\prime}\left[\mathbf{W}\left(\mathbf{x}_{i}, \gamma^{*}\right)\right]^{-1}\left[\mathbf{y}_{i}-\mathbf{m}\left(\mathbf{x}_{i}, \boldsymbol{\theta}_{\mathrm{o}}\right)\right]\right\} \\
& +\mathrm{E}\left\{\left[\mathbf{m}\left(\mathbf{x}_{i}, \boldsymbol{\theta}_{\mathrm{o}}\right)-\mathbf{m}\left(\mathbf{x}_{i}, \boldsymbol{\theta}\right)\right]^{\prime}\left[\mathbf{W}\left(\mathbf{x}_{i}, \gamma^{*}\right)\right]^{-1}\left[\mathbf{m}\left(\mathbf{x}_{i}, \boldsymbol{\theta}_{\mathrm{o}}\right)-\mathbf{m}\left(\mathbf{x}_{i}, \boldsymbol{\theta}\right)\right]\right\} \tag{12.95}
\end{align*}


The first term on the right-hand side in (12.95) does not depend on $\boldsymbol{\theta}$ and the second term is zero when $\boldsymbol{\theta}=\boldsymbol{\theta}_{\mathrm{o}}$. As always, identification requires that the latter term be zero only when $\boldsymbol{\theta}=\boldsymbol{\theta}_{\mathrm{o}}$.

Not surprisingly, the WMNLS estimator has desirable asymptotic efficiency properties if (12.94) holds, along with\\
$\operatorname{Var}\left(\mathbf{y}_{i} \mid \mathbf{x}_{i}\right)=\mathbf{W}\left(\mathbf{x}_{i}, \gamma_{\mathrm{o}}\right) \quad$ for some $\gamma_{\mathrm{o}}$.

Under these assumptions, it is easy to show that the conditional information matrix equality (12.96) holds. Matrix algebra can be used to show directly that the asymptotic variance of Avar $\sqrt{N}\left(\hat{\boldsymbol{\theta}}-\boldsymbol{\theta}_{\mathrm{o}}\right)$ under correct second moment specification is\\
smaller than that of any other WMNLS estimator using weights that are a function of $\mathbf{x}_{i}$ (including, of course, constant weights). In particular, the WMNLS estimator is more efficient than MNLS. But remember that MNLS does not require the strict exogeneity assumption for consistency. In Chapter 14, we will develop a general efficiency framework that allows us to conclude that WMNLS is the asymptotically efficient estimator in a broad class of instrumental variables estimators.

Even if we admit the possibility that our model $\mathbf{W}\left(\mathbf{x}_{i}, \gamma\right)$ is misspecified for the conditional variance, there are still good reasons to apply WMNLS, at least under (12.93). As with the linear case, it will often be the case that a misspecified model of the variance matrix that nevertheless captures key features of the conditional second moments might lead to a more efficient estimator of $\boldsymbol{\theta}_{\mathrm{o}}$ than an estimator that ignores variances and covariances: the MNLS estimator. This is the key insight in the generalized estimating equation (GEE) literature in statistics-see, for example, Liang and Zeger (1986)-which is typically applied to panel data sets (and cluster samples) but whose insights also apply to SUR-like systems of equations. Borrowing from GEE nomenclature, we refer to $\mathbf{W}(\mathbf{x}, \gamma)$ as a working variance matrix, which is allowed, and in many cases is known, to be misspecified. We discuss some ways of choosing this matrix below. The GEE approach is closely related to quasi-maximum likelihood methods, and we cover these, along with GEE for panel data, in Chapter 13.

Given $\hat{\gamma}$ from a first-step estimation, the first-order conditions for the WMNLS estimator, $\hat{\boldsymbol{\theta}}$, are


\begin{equation*}
\sum_{i=1}^{N} \nabla_{\boldsymbol{\theta}} \mathbf{m}\left(\mathbf{x}_{i}, \hat{\boldsymbol{\theta}}\right)^{\prime}\left[\mathbf{W}\left(\mathbf{x}_{i}, \hat{\gamma}\right)\right]^{-1}\left[\mathbf{y}_{i}-\mathbf{m}\left(\mathbf{x}_{i}, \hat{\boldsymbol{\theta}}\right)\right]=\mathbf{0} \tag{12.97}
\end{equation*}


(The GEE approach works off a very similar set of moment conditions-namely, the first occurrence of $\hat{\boldsymbol{\theta}}$ is replaced with an initial, $\sqrt{N}$-consistent estimator of $\boldsymbol{\theta}_{\mathrm{o}}$-and the GEE estimator is $\sqrt{N}$-equivalent to the WMNLS estimator. Here we focus on equation (12.97).)

With a possibly misspecified conditional variance matrix, the asymptotic variance of $\hat{\boldsymbol{\theta}}$ should be estimated using a Huber-White sandwich form,


\begin{align*}
&\left(\sum_{i=1}^{N} \nabla_{\boldsymbol{\theta}} \mathbf{m}_{i}(\hat{\boldsymbol{\theta}})^{\prime}\left[\mathbf{W}_{i}(\hat{\gamma})\right]^{-1} \nabla_{\boldsymbol{\theta}} \mathbf{m}_{i}(\hat{\boldsymbol{\theta}})\right)^{-1}\left(\sum_{i=1}^{N} \nabla_{\boldsymbol{\theta}} \mathbf{m}_{i}(\hat{\boldsymbol{\theta}})^{\prime}\left[\mathbf{W}_{i}(\hat{\gamma})\right]^{-1} \hat{\mathbf{u}}_{i} \hat{\mathbf{u}}_{i}^{\prime}\left[\mathbf{W}_{i}(\hat{\gamma})\right]^{-1} \nabla_{\boldsymbol{\theta}} \mathbf{m}_{i}(\hat{\boldsymbol{\theta}})\right) \\
& \cdot\left(\sum_{i=1}^{N} \nabla_{\boldsymbol{\theta}} \mathbf{m}_{i}(\hat{\boldsymbol{\theta}})^{\prime}\left[\mathbf{W}_{i}(\hat{\gamma})\right]^{-1} \nabla_{\boldsymbol{\theta}} \mathbf{m}_{i}(\hat{\boldsymbol{\theta}})\right)^{-1}, \tag{12.98}
\end{align*}


where the notation should be clear. (As usual, saying that (12.98) is valid for $\widehat{\operatorname{Avar}(\hat{\boldsymbol{\theta}})}$ means that $\mathrm{Avar} \sqrt{N}\left(\hat{\boldsymbol{\theta}}-\boldsymbol{\theta}_{\mathrm{o}}\right)$ is consistently estimated by dividing (12.98) by $N$.) Problem 12.11 asks you to derive this formula, along with the simplified formulawith the final two terms in (12.98) dropped-that is valid under $\operatorname{Var}\left(\mathbf{y}_{i} \mid \mathbf{x}_{i}\right)= \mathbf{W}\left(\mathbf{x}_{i}, \gamma_{\mathrm{o}}\right)$ for some $\gamma_{\mathrm{o}}$.

When $\mathbf{W}\left(\mathbf{x}_{i}, \gamma\right)$ is chosen to be a diagonal matrix, the WMNLS estimator is typically robust to violations of the strict exogeneity assumption. For example, in a panel data setting where $\mathrm{E}\left(y_{i t} \mid \mathbf{x}_{i t}\right)=m_{t}\left(\mathbf{x}_{i t}, \boldsymbol{\theta}_{\mathrm{o}}\right)$, we can choose $\mathbf{W}\left(\mathbf{x}_{i}, \gamma\right)$ to be diagonal where the $t$ th diagonal depends only on $\mathbf{x}_{i t}$. The resulting estimator is a pooled weighted nonlinear least squares (PWNLS) estimator, which is often useful for dynamic panel data models, or other panel data models without strict exogeneity; see Problem 12.13. In Chapters 13 and 18 we will cover quasi-maximum likelihood estimators that are asymptotically equivalent to PWNLS but can be obtained more easily via one-step estimation.

Sometimes we might choose $\mathbf{W}\left(\mathbf{x}_{i}, \gamma\right)=\boldsymbol{\Omega}$, that is, use a matrix where the variances and covariances do not depend on $\mathbf{x}_{i}$. A consistent estimator of the unconditional variance matrix of $\mathbf{u}_{i} \equiv \mathbf{y}_{i}-\mathbf{m}\left(\mathbf{x}_{i}, \boldsymbol{\theta}_{\mathrm{o}}\right)$ is


\begin{equation*}
\hat{\boldsymbol{\Omega}}=N^{-1} \sum_{i=1}^{N} \check{\mathbf{u}}_{i} \check{\mathbf{u}}_{i}^{\prime} \tag{12.99}
\end{equation*}


where the $\check{\mathbf{u}}_{i}$ are the vectors of MNLS residuals. When we use $\hat{\boldsymbol{\Omega}}$ in (12.99), we have what is sometimes called the nonlinear SUR estimator. Equation (12.99) places no restrictions on the unconditional variances and covariances. In panel data cases we might, for example, restrict $\hat{\boldsymbol{\Omega}}$ to have a random effects or an $\operatorname{AR}(1)$ structure. The asymptotic analysis of the estimator with constant $\mathbf{W}\left(\mathbf{x}_{i}, \gamma\right)$ follows from the general WMNLS framework. In particular, the nonlinear SUR estimator is consistent and $\sqrt{N}$-asymptotically normal under (12.93) whether or not $\operatorname{Var}\left(\mathbf{y}_{i} \mid \mathbf{x}_{i}\right)$ is constant. Of course, if $\operatorname{Var}\left(\mathbf{y}_{i} \mid \mathbf{x}_{i}\right)$ depends on $\mathbf{x}_{i}$, there could be an alternative WMNLS estimator that is more efficient than the nonlinear SUR estimator, but that hinges on our ability to find $\mathbf{W}\left(\mathbf{x}_{i}, \gamma\right)$ that provides a "better" approximation to $\operatorname{Var}\left(\mathbf{y}_{i} \mid \mathbf{x}_{i}\right)$ than a constant matrix.

Models for conditional variances $\operatorname{Var}\left(y_{i g} \mid \mathbf{x}_{i}\right)$ are relatively straightforward because we have many univariate distributions to draw on (see Section 13.11.3). Directly specifying models for conditional covariances, when they actually depend on $\mathbf{x}_{i}$, is more difficult. A fruitful approach is to specify conditional variances for each $g$ but then to nominally assume constant conditional correlations. Then, if $\rho_{g h}$ is (nominally) $\operatorname{Corr}\left(y_{i g}, y_{i h} \mid \mathbf{x}_{i}\right)$, we have $\operatorname{Cov}\left(y_{i g}, y_{i h} \mid \mathbf{x}_{i}\right)=$\\
$\rho_{g h}\left[\operatorname{Var}\left(y_{i g} \mid \mathbf{x}_{i}\right) \operatorname{Var}\left(y_{i h} \mid \mathbf{x}_{i}\right)\right]^{1 / 2}$. Let $\mathbf{V}(\mathbf{x}, \boldsymbol{\omega})$ be the $G \times G$ diagonal matrix with the proposed variances down its diagonal, and let $\mathbf{R}(\boldsymbol{\rho})$ denote the $G \times G$ matrix of proposed constant correlations, which depends on the $J$-vector of parameters $\boldsymbol{\rho}$. In the GEE literature, $\mathbf{R}(\boldsymbol{\rho})$ is called a working correlation matrix because there is no presumption that it truly contains the conditional correlations. In fact, there is no presumption that the conditional correlations are even constant.

Given $\mathbf{V}(\mathbf{x}, \boldsymbol{\omega})$ and $\mathbf{R}(\boldsymbol{\rho})$, we can write the working variance matrix as\\
$\mathbf{W}\left(\mathbf{x}_{i}, \gamma\right)=\mathbf{V}\left(\mathbf{x}_{i}, \boldsymbol{\omega}\right)^{1 / 2} \mathbf{R}(\boldsymbol{\rho}) \mathbf{V}\left(\mathbf{x}_{i}, \boldsymbol{\omega}\right)^{1 / 2}$,\\
where $\mathbf{V}\left(\mathbf{x}_{i}, \boldsymbol{\omega}\right)^{1 / 2}$ is the matrix square root. Note that $\gamma$ contains $\boldsymbol{\omega}$ and $\boldsymbol{\rho}$.\\
To implement WMNLS (GEE) under (12.100), we need to estimate $\boldsymbol{\omega}$ and $\boldsymbol{\rho}$. This typically proceeds by MNLS to first obtain residuals $\check{u}_{i g}$, then by using $\check{u}_{i g}^{2}$ and the variance models, say $v_{g}\left(\mathbf{x}_{i}, \boldsymbol{\omega}\right)$, to estimate the variances, say $\check{v}_{i g} \equiv v_{i g}\left(\mathbf{x}_{i}, \hat{\boldsymbol{\omega}}\right)$. (In many cases, the $v_{g}(\mathbf{x}, \boldsymbol{\omega})$ functions depend on the mean parameters $\boldsymbol{\theta}$ and a single additional parameter; we will see this explicitly in Section 13.11.3 and Chapter 18.) We can then use the standardized residuals $\check{u}_{i g} / \sqrt{\check{v}_{i g}}$ for all $i$ and $g$ to estimate the parameters in the working correlation matrix $\mathbf{R}(\boldsymbol{\rho})$. The details depend on how $\mathbf{R}(\boldsymbol{\rho})$ is specified.

The most general specification (once we restrict ourselves to a matrix that does not depend on $\mathbf{x}$ ) is an unstructured working correlation matrix, where the elements of $\mathbf{R}(\boldsymbol{\rho})$ are unrestricted (except, of course, for requiring the matrix to be a valid correlation matrix):\\
$\mathbf{R}(\boldsymbol{\rho})=\left(\begin{array}{ccccc}1 & \rho_{12} & \rho_{13} & \cdots & \rho_{1 G} \\ \rho_{12} & 1 & \rho_{23} & \cdots & \rho_{2 G} \\ \rho_{13} & \rho_{23} & \ddots & \ddots & \vdots \\ \vdots & \vdots & \vdots & 1 & \rho_{G-1, G} \\ \rho_{1 G} & \rho_{2 G} & \cdots & \rho_{G-1, G} & 1\end{array}\right)$.

Then, we can estimate each $\rho_{g h}$ as\\
$\hat{\rho}_{g h}=$ Sample Correlation $\left(\check{u}_{i g} / \sqrt{\check{v}_{i g}}, \check{u}_{i h} / \sqrt{\check{v}_{i h}}\right)$.

Under general conditions, the $\hat{\rho}_{g h}$ converge to, say, $\rho_{g h}^{*}$ the population correlation between $u_{i g} / \sqrt{v_{g}\left(\mathbf{x}_{i}, \boldsymbol{\omega}^{*}\right)}$ and $u_{i h} / \sqrt{v_{h}\left(\mathbf{x}_{i}, \boldsymbol{\omega}^{*}\right)}$, where $\boldsymbol{\omega}^{*}=\operatorname{plim}(\hat{\boldsymbol{\omega}})$.

A common form of $\mathbf{R}(\boldsymbol{\rho})$ in panel data settings is an exchangeable working correlation matrix, which introduces a single correlation parameter, $\rho$, for all pairs. Then, $\hat{\rho}$ is obtained by averaging the $\hat{\rho}_{g h}$ (but where $g$ and $h$ denote different time periods) across all nonredundant pairs with $g \neq h$. An exchangeable correlation matrix allows\\
for a random effects-type correlation structure but with conditional variances changing over time. We will have more to say on this setup in Chapter 13.

Once $\hat{\boldsymbol{\omega}}$ and $\hat{\boldsymbol{\rho}}$ have been obtained, the working variance matrix estimates, $\mathbf{W}_{i}(\hat{\gamma})=\mathbf{V}\left(\mathbf{x}_{i}, \hat{\boldsymbol{\omega}}\right)^{1 / 2} \mathbf{R}(\hat{\boldsymbol{\rho}}) \mathbf{V}\left(\mathbf{x}_{i}, \hat{\boldsymbol{\omega}}\right)^{1 / 2}$, are easy to obtain, and they can be used in WMNLS. In special cases it may be reasonable to assume $\mathbf{W}(\mathbf{x}, \gamma)$ is correctly specified for $\operatorname{Var}\left(\mathbf{y}_{i} \mid \mathbf{x}_{i}\right)$, but most of the time one should use equation (12.98) as the variance matrix estimator.

Finally, a final word of caution. The variance matrix estimator (12.98) assumes that the conditional mean is correctly specified, as stated explicitly in equation (12.93). Thus, this estimator is fully robust only if the conditional mean is correctly specified; it is only semirobust if we entertain misspecification of $\mathrm{E}\left(y_{i} \mid \mathbf{x}_{i}\right)$. Unfortunately, because WMNLS is a two-step estimator, a variance matrix estimator that allows conditional mean misspecification is much more complicated than (12.98) because it depends on the sampling variability of $\hat{\gamma}$. In particular, it is not enough just to replace the outer ends of the sandwich with the unconditional Hessian evaluated at the estimates. (This is why, in the application of WMNLS to GEE, the standard errors from (12.98) are often explicitly labeled as semirobust.)

\subsection*{12.10 Quantile Estimation}
The introduction to this chapter included a brief discussion of the problem of estimating a conditional median function, and we discussed how the consistency of LAD follows from Theorem 12.2. Estimation of conditional medians and, more generally, conditional quantiles, is increasingly popular in empirical research. It has been long recognized that while estimating conditional mean functions is valuable, the partial effect of an explanatory variable can have very different effects across different segments of a population. Quantile estimation allows us to study such effects. Koenker and Bassett (1978) developed the theory of quantile regression, and Buchinsky (1998), Peracchi (2001), and Koenker (2005) provide recent treatments. Chamberlain (1994) and Buchinsky (1994) are influential applications of quantile regression to estimate changes in the wage distribution in the United States.

\subsection*{12.10.1 Quantiles, the Estimation Problem, and Consistency}
Let $y_{i}$ denote a random draw from a population. Then, for $0<\tau<1, q(\tau)$ is a $\tau$ th quantile of the distribution of $y_{i}$ if $\mathrm{P}\left(y_{i} \leq q(\tau)\right) \geq \tau$ and $\mathrm{P}\left(y_{i} \geq q(\tau)\right) \geq 1-\tau$. A special case is the median when $\tau=1 / 2$. For notational convenience, we will write the $\tau$ th quantile of $y_{i}$ as Quant $_{\tau}\left(y_{i}\right)$.

Typically, we are interested in modeling quantiles conditional on a set of covariates $\mathbf{x}_{i}$. In most applications, the assumption is that the quantiles are linear in parameters (and we will show them linear in $\mathbf{x}_{i}$, even though we can, as usual, choose $\mathbf{x}_{i}$ to include nonlinear functions of underlying explanatory variables). Under linearity, we have\\
$\operatorname{Quant}_{\tau}\left(y_{i} \mid \mathbf{x}_{i}\right)=\alpha_{\mathrm{o}}(\tau)+\mathbf{x}_{i} \boldsymbol{\beta}_{\mathrm{o}}(\tau)$,\\
where, for reasons to be seen shortly, we explicitly introduce an intercept and explicitly show the intercept and slopes depending on $\tau$.

To estimate the parameters in a conditional quantile function, it is very helpful to know whether a population quantile solves a population extremum problem. We know that the conditional mean minimizes the expected squared error and, in the introduction, we asserted that the conditional median (when $\tau=.5$ ) minimizes the expected absolute error. Generally, if $q_{\mathrm{o}}(\tau)$ is the $\tau$ th quantile of $y_{i}$, then $q_{\mathrm{o}}(\tau)$ solves


\begin{equation*}
\min _{q \in \mathbb{R}} \mathrm{E}\left\{\left(\tau 1\left[y_{i}-q \geq 0\right]+(1-\tau) 1\left[y_{i}-q<0\right]\right)\left|y_{i}-q\right|\right\}, \tag{12.104}
\end{equation*}


where $1[\cdot]$ is the indicator function equal to one if the statement in brackets is true and zero otherwise. The function\\
$c_{\tau}(u)=(\tau 1[u \geq 0]+(1-\tau) 1[u<0])|u|=(\tau-1[u<0]) u$\\
is called the asymmetric absolute loss function, the $\boldsymbol{\tau}$-absolute loss function, or the check function (because its graph resembles a check mark) (see, for example, Manski, 1988, Sect. 4.2.4). The slope of $c_{\tau}(u)$ is $\tau$ when $u>0$ and $-(1-\tau)$ when $u<0$; the slope is undefined at $u=0$. When $\tau=.5$, the check function is simply the absolute value divided by two, and so it is symmetric about zero. If $\tau>.5$, the slope of the loss for $u>0$ is greater than the absolute value of the slope for $u<0$; the opposite holds if $\tau<.5$.

It follows immediately that a conditional quantile minimizes the asymmetric absolute loss function conditional on $\mathbf{x}_{i}$. (Of course, when $\tau=.5$, we are back to the fact that the median minimizes the absolute error.) Therefore, we can immediately apply the analogy principle to obtain consistent estimators of the parameters in equation (12.103):\\
$\min _{\alpha \in \mathbb{R}, \boldsymbol{\beta} \in \mathbb{R}^{K}} \sum_{i=1}^{N} c_{\tau}\left(y_{i}-\alpha-\mathbf{x}_{i} \boldsymbol{\beta}\right)$.

Under the assumption that $\boldsymbol{\theta}_{\mathrm{o}}(\tau)=\left(\alpha_{\mathrm{o}}(\tau), \boldsymbol{\beta}_{\mathrm{o}}(\tau)^{\prime}\right)^{\prime}$ is the unique minimizer of $\mathrm{E}\left[c_{\tau}\left(y_{i}-\alpha-\mathbf{x}_{i} \boldsymbol{\beta}\right)\right]$-it is guaranteed to be a solution-the quantile regression estimator is consistent under very weak regularity conditions. Note that $c_{\tau}\left(y_{i}-\alpha-\mathbf{x}_{i} \boldsymbol{\beta}\right)$ is continuous in the parameters because the check function is continuous. However, the check function is not differentiable at zero.

Before we discuss estimation across different quantiles, it is useful to spend some time on the leading case of the conditional median. In many applications, the LAD estimator is applied along with OLS, often to supposedly demonstrate the sensitivity of OLS to influential observations. It is no secret that OLS, because it minimizes the sum of squared residuals, can be sensitive to the inclusion of extreme observations. In the context of specific models of data contamination, one can make precise the notion that OLS is "nonrobust" to influential observations, or "outliers." By contrast, the LAD estimator (and quantile estimators more generally) are "robust" to influential observations. We need not develop a formal framework for defining robustness to outlying data-as in, for example, Huber (1981)-to understand the main point: OLS is sensitive to changes in extreme data points because the mean is sensitive to changes in extreme values; LAD is insensitive to changes in extreme data points because the median is insensitive to changes in extreme values. This point is easy to illustrate by selecting three positive integers, computing the mean and median, multiplying the largest value by 10 , and computing the mean and median again. The mean can increase dramatically, while the median will not change.

The insensitivity of the median to changes in extreme values is desirable, but we should not overlook an important point: sometimes, probably more often than not, we are interested in partial effects on the conditional mean. If that is the case, then we must recognize that LAD does not generally consistently estimate parameters in a correctly specified conditional mean. Only least squares does (assuming we rule out error distributions with very thick tails). Therefore, one must be very careful in attributing differences between LAD and OLS to outliers; there are other reasons the estimates may differ significantly. If we define robustness to mean consistently estimating the parameters of the conditional mean, LAD is not a robust estimator of conditional mean parameters because consistency holds only under additional restrictions on the conditional distribution. (Other so-called robust estimatorswhere "robust" means insensitivity to outlying observations-are not robust for estimating the conditional mean in that they also rely on symmetry for consistency. See Huber (1981) for a general treatment and also Peracchi (2001).)

To study the assumptions under which LAD and OLS estimate the same parameters, it is helpful to write a model for a random draw $i$ as\\
$y_{i}=\alpha_{\mathrm{o}}+\mathbf{x}_{i} \boldsymbol{\beta}_{\mathrm{o}}+u_{i}$.

If we assume\\
$\mathrm{D}\left(u_{i} \mid \mathbf{x}_{i}\right)$ is symmetric about zero,\\
then $\mathrm{E}\left(u_{i} \mid \mathbf{x}_{i}\right)=\operatorname{Med}\left(u_{i} \mid \mathbf{x}_{i}\right)=0$, which means that the deterministic part of (12.106), $\alpha_{\mathrm{o}}+\mathbf{x}_{i} \boldsymbol{\beta}_{\mathrm{o}}$, is both the conditional mean and conditional median of $y_{i}\left(\right.$ and $\mathrm{D}\left(y_{i} \mid \mathbf{x}_{i}\right)$ is symmetric about $\alpha_{\mathrm{o}}+\mathbf{x}_{i} \boldsymbol{\beta}_{\mathrm{o}}$ ). In discussions of the sensitivity of OLS to outliers, and the superiority of LAD under such circumstances, conditional symmetry is often maintained, if only implicitly. If we maintain conditional symmetry, then the analysis is fairly clean, as the trade-offs between LAD and OLS are readily obtained. Under conditional symmetry, LAD can have a smaller asymptotic variance-see the next section for derivation-than OLS for fat-tailed distributions. But, as is well known, OLS will be more efficient for estimating the mean (median) parameters under certain thin-tailed distributions, such as the normal.

Are there other assumptions under which OLS and LAD are both consistent for the parameters in (12.106)? Yes. Suppose that, instead of (12.107), we assume independence and take as the normalization $\mathrm{E}\left(u_{i}\right)=0$ :\\
$\mathrm{D}\left(u_{i} \mid \mathbf{x}_{i}\right)=\mathrm{D}\left(u_{i}\right) \quad$ and $\quad \mathrm{E}\left(u_{i}\right)=0$.

Under (12.108), OLS consistently estimates $\alpha_{\mathrm{o}}$ and $\boldsymbol{\beta}_{\mathrm{o}}$. Notice that $u_{i}$ need not have a symmetric distribution. When it does not, $\operatorname{Med}\left(u_{i}\right) \equiv \eta_{\mathrm{o}} \neq 0$. Nevertheless, by independence,\\
$\operatorname{Med}\left(y_{i} \mid \mathbf{x}_{i}\right)=\alpha_{\mathrm{o}}+\mathbf{x}_{i} \boldsymbol{\beta}_{\mathrm{o}}+\operatorname{Med}\left(u_{i}\right)=\left(\alpha_{\mathrm{o}}+\eta_{\mathrm{o}}\right)+\mathbf{x}_{i} \boldsymbol{\beta}_{\mathrm{o}}$.\\
This equation immediately implies that the LAD slope estimators are consistent for $\boldsymbol{\beta}_{\mathrm{o}}$. Therefore, OLS and LAD should provide similar estimates of the slope parameters. But they will estimate different intercepts.

Unfortunately, in many applications of LAD, the independence assumption is clearly violated. Heteroskedasticity in $\operatorname{Var}\left(u_{i} \mid \mathbf{x}_{i}\right)$ is common when $y_{i}$ is a variable such as wealth, income, or pension contributions. In addition, conditional wealth and income distributions tend to be skewed. Therefore, when LAD methods are applied alongside OLS, there are often reasons to think a priori that OLS and LAD will not produce similar slope estimates. (In fact, it is unlikely that the conditional mean and conditional median are both linear in $\mathbf{x}_{i}$.) But important differences in the OLS and LAD estimates need have nothing to do with the presence of outliers.

Sometimes one can use a transformation to ensure conditional symmetry or the independence assumption in (12.108). When $y_{i}>0$, the most common transforma-\\
tion is the natural log. Often, the linear model $\log \left(y_{i}\right)=\alpha_{\mathrm{o}}+\mathbf{x}_{i} \boldsymbol{\beta}_{\mathrm{o}}+u_{i}$ is more likely to satisfy symmetry or independence. Suppose that symmetry about zero holds in the linear model for $\log \left(y_{i}\right)$. Then, because the median passes through monotonic functions (unlike the expectation), $\operatorname{Med}\left(y_{i} \mid \mathbf{x}_{i}\right)=\exp \left(\operatorname{Med}\left[\log \left(y_{i}\right) \mid \mathbf{x}_{i}\right]\right)=\exp \left(\alpha_{\mathrm{o}}+\mathbf{x}_{i} \boldsymbol{\beta}_{\mathrm{o}}\right)$, and so we can easily recover the partial effects on the median of $y_{i}$ itself. By contrast, we cannot generally find $\mathrm{E}\left(y_{i} \mid \mathbf{x}_{i}\right)=\exp \left(\alpha_{\mathrm{o}}+\mathbf{x}_{i} \boldsymbol{\beta}_{\mathrm{o}}\right) \mathrm{E}\left[\exp \left(u_{i}\right) \mid \mathbf{x}_{i}\right]$. If instead we assume (12.108), then $\operatorname{Med}\left(y_{i} \mid \mathbf{x}_{i}\right)$ and $\mathrm{E}\left(y_{i} \mid \mathbf{x}_{i}\right)$ are both exponential functions of $\mathbf{x}_{i} \boldsymbol{\beta}_{\mathrm{o}}$, but with different "intercepts" inside the exponential function.

Incidentally, the fact that the median passes through monotonic functions is very handy for applying LAD to a variety of problems, including some in Chapters 15 and 17. But the expectation operator has useful properties that the median does not: linearity and the law of iterated expectations. For example, suppose we begin with a random coefficient model $y_{i}=a_{i}+\mathbf{x}_{i} \mathbf{b}_{i}$, where $a_{i}$ is the heterogeneous intercept and $\mathbf{b}_{i}$ is a $1 \times K$ vector of heterogeneous slopes ("random coefficients"). If we assume that $\left(a_{i}, \mathbf{b}_{i}\right)$ is independent of $\mathbf{x}_{i}$, then\\
$\mathrm{E}\left(y_{i} \mid \mathbf{x}_{i}\right)=\mathrm{E}\left(a_{i} \mid \mathbf{x}_{i}\right)+\mathbf{x}_{i} \mathrm{E}\left(\mathbf{b}_{i} \mid \mathbf{x}_{i}\right) \equiv \alpha_{\mathrm{o}}+\mathbf{x}_{i} \boldsymbol{\beta}_{\mathrm{o}}$,\\
where $\alpha_{\mathrm{o}}=\mathrm{E}\left(a_{i}\right)$ and $\boldsymbol{\beta}_{\mathrm{o}}=\mathrm{E}\left(\mathbf{b}_{i}\right)$. Because OLS consistently estimates the parameters of a conditional mean linear in those parameters, OLS consistently estimates the population averaged effects, or average partial effects, $\boldsymbol{\beta}_{\mathrm{o}}$ (see also Section 4.4.4). Generally, even under independence, there is no way to derive $\operatorname{Med}\left(y_{i} \mid \mathbf{x}_{i}\right)$, and it cannot be shown that LAD estimation of a linear model estimates average or median partial effects. Angrist, Chernozhukov, and Fernandez-Val (2006) provide a treatment of LAD (and quantile regression) under misspecification and characterize the probability limit of the LAD estimator.

We now turn to general quantile estimation. In many applications, one estimates linear conditional quantile functions for various quantiles. Having a set of estimated linear quantiles allows one to see how various explanatory variables differentially affect different parts of the distribution $\mathrm{D}\left(y_{i} \mid \mathbf{x}_{i}\right)$. Of course, nothing guarantees that all or even several of the conditional quantiles are actually linear. With additive errors independent of the regressors, we have equation (12.103) but with common slopes, $\boldsymbol{\beta}_{\mathrm{o}}(\tau)=\boldsymbol{\beta}_{\mathrm{o}}$ for all $\tau$. Then, we can estimate the common slopes using any quantile with $0<\tau<1$. In most applications of quantile regression, the whole point is to see how the effects of the covariates change with the quantile, and so different linear models are estimated for different quantiles. Practically speaking, the conditions under which the quantile functions do not cross are quite restrictive, and the estimated quantiles often do not show uniformly increasing or decreasing slopes as $\tau$ ranges between zero and one. Koenker (2005) provides further discussion.

\subsection*{12.10.2 Asymptotic Inference}
As mentioned in Section 12.3, LAD estimation falls outside Theorem 12.3 because the objective function is not twice continuously differentiable with positive definite expected Hessian (at $\boldsymbol{\theta}_{\mathrm{o}}$ ). For general quantile regression, the nondifferentiability in the check function occurs only at zero. To simplify notation, for a given $\tau$ we now write\\
$y_{i}=\mathbf{x}_{i} \boldsymbol{\theta}_{\mathrm{o}}+u_{i}, \quad \operatorname{Quant}_{\tau}\left(u_{i} \mid \mathbf{x}_{i}\right)=0$,\\
where the first element of $\mathbf{x}_{i}$ is unity. Provided there is sufficient variation in the distribution of $\mathbf{x}_{i}$, so that the probability that $y_{i}-\mathbf{x}_{i} \hat{\boldsymbol{\theta}}$ is zero is sufficiently small, the nonsmoothness of $c_{\tau}(u)$ at $u=0$ does not cause a serious problem. The real complication is that the second derivative of the check function is zero everywhere it is defined, that is, for all $u \neq 0$. Interestingly, although the usual mean value expansion of the score can no longer be applied, it is possible to modify the argument and obtain an influence function representation for the quantile regression estimator. If we write the objective function as\\
$q\left(\mathbf{w}_{i}, \boldsymbol{\theta}\right)=\tau 1\left[y_{i}-\mathbf{x}_{i} \boldsymbol{\theta} \geq 0\right]\left(y_{i}-\mathbf{x}_{i} \boldsymbol{\theta}\right)-(1-\tau) 1\left[y_{i}-\mathbf{x}_{i} \boldsymbol{\theta}<0\right]\left(y_{i}-\mathbf{x}_{i} \boldsymbol{\theta}\right)$,\\
then we can define a score function as\\
$\mathbf{s}_{i}(\boldsymbol{\theta})=-\mathbf{x}_{i}^{\prime}\left\{\tau 1\left[y_{i}-\mathbf{x}_{i} \boldsymbol{\theta} \geq 0\right]-(1-\tau) 1\left[y_{i}-\mathbf{x}_{i} \boldsymbol{\theta}<0\right]\right\}$.

We refer to $\mathbf{s}_{i}(\boldsymbol{\theta})$ as a score function because, for $\boldsymbol{\theta}$ such that $y_{i}-\mathbf{x}_{i} \boldsymbol{\theta} \neq 0$, $\mathbf{s}_{i}(\boldsymbol{\theta})$ is in fact the transpose of the gradient of the objective function, just as before. The hope is that we do not have to worry about what happens when the objective function is nondifferentiable. If $u_{i}$ has a continuous distribution at zero, then $\mathrm{P}\left(y_{i}-\mathbf{x}_{i} \boldsymbol{\theta}_{\mathrm{o}}=0\right)=0$. In other words, at $\boldsymbol{\theta}=\boldsymbol{\theta}_{\mathrm{o}}$, we can ignore the possibility of observing data where the objective function is nondifferentiable at the true value. Because $\hat{\boldsymbol{\theta}}$ is consistent for $\boldsymbol{\theta}_{\mathrm{o}}$, as the sample size grows there is less and less chance that $q\left(\mathbf{w}_{i}, \boldsymbol{\theta}\right)$ is nondifferentiable at $\hat{\boldsymbol{\theta}}$, and so we can use a first-order condition to obtain $\hat{\boldsymbol{\theta}}$.

We can show directly that the score satisfies $\mathrm{E}\left[\mathbf{s}_{i}\left(\boldsymbol{\theta}_{\mathrm{o}}\right) \mid \mathbf{x}_{i}\right]=0$ because $\mathrm{E}\left(1\left[y_{i}-\mathbf{x}_{i} \boldsymbol{\theta}_{\mathrm{o}} \geq 0\right] \mid \mathbf{x}_{i}\right)=\mathrm{P}\left(y_{i} \geq \mathbf{x}_{i} \boldsymbol{\theta}_{\mathrm{o}} \mid \mathbf{x}_{i}\right)=(1-\tau)$ by definition of a quantile. Further, $\mathrm{E}\left(1\left[y_{i}-\mathbf{x}_{i} \boldsymbol{\theta}_{\mathrm{o}}<0\right] \mid \mathbf{x}_{i}\right)=\mathrm{P}\left(y_{i}<\mathbf{x}_{i} \boldsymbol{\theta}_{\mathrm{o}} \mid \mathbf{x}_{i}\right)=\tau$ (when we assume $u_{i}$ has a continuous distribution at zero), and so\\
$\mathrm{E}\left[\mathbf{s}_{i}\left(\boldsymbol{\theta}_{\mathrm{o}}\right) \mid \mathbf{x}_{i}\right]=-\mathbf{x}_{i}^{\prime}(\tau(1-\tau)-(1-\tau) \tau)=\mathbf{0}$.\\
The asymptotic theory requires that we consider solutions that satisfy $N^{-1 / 2} \sum_{i=1}^{N} \mathbf{s}_{i}(\hat{\boldsymbol{\theta}})=o_{p}(1)$. (Technically, there is a chance that the solution to the\\
minimization problem does not solve the first-order condition, but it diminishes.) The key to obtaining an influence function representation is to first compute the expected value of the score, then obtain the Jacobian. If we first obtain the Jacobian of the score, then it is identically zero. Therefore, let $\mathbf{a}(\boldsymbol{\theta}) \equiv \mathrm{E}\left[\mathbf{s}_{i}(\boldsymbol{\theta})\right]$ and assume that the $P \times P$ Jacobian of $\mathbf{a}(\cdot)$, say $\mathbf{A}(\cdot)$, exists. Further, assume that $\mathbf{A}\left(\boldsymbol{\theta}_{\mathrm{o}}\right)$ is nonsingular; in fact, it will almost always be positive definite. By first taking the expectation of $\mathbf{s}_{i}(\boldsymbol{\theta})$, we smooth it out, and under weak conditions the expected value of the score is well behaved. The starting point is to find the expectation conditional on $\mathbf{x}_{i}$ :


\begin{align*}
\mathrm{E}\left[\mathbf{s}_{i}(\boldsymbol{\theta}) \mid \mathbf{x}_{i}\right] & =-\mathbf{x}_{i}^{\prime}\left\{\tau P\left[u_{i} \geq \mathbf{x}_{i}\left(\boldsymbol{\theta}-\boldsymbol{\theta}_{\mathrm{o}}\right) \mid \mathbf{x}_{i}\right]-(1-\tau) P\left[u_{i}<\mathbf{x}_{i}\left(\boldsymbol{\theta}-\boldsymbol{\theta}_{\mathrm{o}}\right) \mid \mathbf{x}_{i}\right]\right. \\
& =-\mathbf{x}_{i}^{\prime}\left\{\tau\left[1-F_{u}\left(\mathbf{x}_{i}\left(\boldsymbol{\theta}-\boldsymbol{\theta}_{\mathrm{o}}\right) \mid \mathbf{x}_{i}\right)\right]-(1-\tau) F_{u}\left(\mathbf{x}_{i}\left(\boldsymbol{\theta}-\boldsymbol{\theta}_{\mathrm{o}}\right) \mid \mathbf{x}_{i}\right)\right\} \\
& =-\mathbf{x}_{i}^{\prime}\left[\tau-F_{u}\left(\mathbf{x}_{i}\left(\boldsymbol{\theta}-\boldsymbol{\theta}_{\mathrm{o}}\right) \mid \mathbf{x}_{i}\right)\right] \tag{12.111}
\end{align*}


where we assume that the conditional cumulative distribution function $F_{u}(\cdot \mid \mathbf{x})$ is continuous at zero. In fact, we assume that $F_{u}(\cdot \mid \mathbf{x})$ is continuously differentiable and denote the density $f_{u}(\cdot \mid \mathbf{x})$. Of course, $\mathrm{E}\left[\mathbf{s}_{i}(\boldsymbol{\theta})\right]$ is just the expected value of (12.111) across the distribution of $\mathbf{x}_{i}$. Assuming that we can interchange the expectation and the Jacobian-which holds under general conditions, as described in Bartle (1966, Chap. 4)-we can first compute the Jacobian of $\mathrm{E}\left[\mathbf{s}_{i}(\boldsymbol{\theta}) \mid \mathbf{x}_{i}\right]$ and then obtain its expected value to obtain $\mathbf{A}(\boldsymbol{\theta})$. But from (12.111), $\nabla_{\theta} \mathrm{E}\left[\mathbf{s}_{i}(\boldsymbol{\theta}) \mid \mathbf{x}_{i}\right]= f_{u}\left(\mathbf{x}_{i}\left(\boldsymbol{\theta}-\boldsymbol{\theta}_{\mathrm{o}}\right) \mid \mathbf{x}_{i}\right) \mathbf{x}_{i}^{\prime} \mathbf{x}_{i}$, and evaluating this Jacobian at $\boldsymbol{\theta}_{\mathrm{o}}$ and taking the expectation (which we assume can be interchanged with the Jacobian) gives\\
$\mathbf{A}_{\mathrm{o}} \equiv \mathbf{A}\left(\boldsymbol{\theta}_{\mathrm{o}}\right)=\mathrm{E}\left[f_{u}\left(0 \mid \mathbf{x}_{i}\right) \mathbf{x}_{i}^{\prime} \mathbf{x}_{i}\right]$.

Using the methods of Huber (1967) and Newey and McFadden (1994, Sect. 7), one can derive the representation


\begin{equation*}
\sqrt{N}\left(\hat{\boldsymbol{\theta}}-\boldsymbol{\theta}_{\mathrm{o}}\right)=-\mathbf{A}_{\mathrm{o}}^{-1} N^{-1 / 2} \sum_{i=1}^{N} \mathbf{s}_{i}\left(\boldsymbol{\theta}_{\mathrm{o}}\right)+o_{p}(1) \tag{12.113}
\end{equation*}


Because $\mathbf{s}_{i}\left(\boldsymbol{\theta}_{\mathrm{O}}\right)$ has zero mean, it follows (under mild conditions) that\\
$\sqrt{N}\left(\hat{\boldsymbol{\theta}}-\boldsymbol{\theta}_{\mathrm{o}}\right) \xrightarrow{d} \operatorname{Normal}\left(\mathbf{0}, \mathbf{A}_{\mathrm{o}}^{-1} \mathbf{B}_{\mathrm{o}} \mathbf{A}_{\mathrm{o}}^{-1}\right)$,\\
where\\
$\mathbf{B}_{\mathrm{o}} \equiv \mathrm{E}\left[\mathbf{s}_{i}\left(\boldsymbol{\theta}_{\mathrm{o}}\right) \mathbf{s}_{i}\left(\boldsymbol{\theta}_{\mathrm{o}}\right)^{\prime}\right]=\tau(1-\tau) \mathrm{E}\left(\mathbf{x}_{i}^{\prime} \mathbf{x}_{i}\right)$.\\
(Problem 12.14 asks you to derive the expression for $\mathbf{B}_{\mathrm{o}}$.)

The matrix $\mathbf{B}_{\mathrm{o}}$ is simple to estimate (recall that we have chosen $\tau$ ):\\
$\hat{\mathbf{B}}=\tau(1-\tau)\left(N^{-1} \sum_{i=1}^{N} \mathbf{x}_{i}^{\prime} \mathbf{x}_{i}\right)$.

We can also use the average outer product of the scores, $\mathbf{s}_{i}(\hat{\boldsymbol{\theta}})$. This estimator is generally consistent when $\mathrm{P}\left(y_{i}-\mathbf{x}_{i} \boldsymbol{\theta}_{\mathrm{o}}=0\right)=0$ by Lemma 12.1 (allowing the function to be discontinuous at $\boldsymbol{\theta}_{\mathrm{o}}$ with probability zero).

The matrix $\mathbf{A}_{\mathrm{o}}$ is much more difficult to estimate because, at first glance, it seems to require estimation of $f_{u}\left(0 \mid \mathbf{x}_{i}\right)$, the conditional density of $u_{i}$ at zero. One approach is to use a nonparametric density estimator, but these estimators can be imprecise, especially when the dimension of $\mathbf{x}_{i}$ is large. It turns out that we can estimate $\mathbf{A}_{\mathrm{o}}$ more simply, using an approach due to Powell (1991). To sketch the approach, we can use the assumption that $f_{u}\left(\cdot \mid \mathbf{x}_{i}\right)$ is differentiable at zero, along with the definition of a derivative, to approximate $f_{u}\left(0 \mid \mathbf{x}_{i}\right)$ as

$$
\begin{aligned}
{\left[F_{u}\left(h_{N} \mid \mathbf{x}_{i}\right)-F_{u}\left(-h_{N} \mid \mathbf{x}_{i}\right)\right] / 2 h_{N} } & =\mathrm{P}\left(-h_{N} \leq u_{i} \leq h_{N} \mid \mathbf{x}_{i}\right) / 2 h_{N} \\
& =\mathrm{P}\left(\left|u_{i}\right| \leq h_{N} \mid \mathbf{x}_{i}\right) / 2 h_{N},
\end{aligned}
$$

where $\left\{h_{N}\right\}$ is a sequence of positive numbers with $h_{N} \rightarrow 0$. (We will see in a moment why we subscript $h_{N}$ with the sample size, $N$.) It is convenient to write the conditional probability as a conditional expectation using indicator functions, $\mathrm{P}\left(\left|u_{i}\right| \leq h_{N} \mid \mathbf{x}_{i}\right)= \mathrm{E}\left(1\left[\left|u_{i}\right| \leq h_{N}\right] \mid \mathbf{x}_{i}\right)$. Therefore, an approximation of $\mathrm{E}\left[f_{u}\left(0 \mid \mathbf{x}_{i}\right) \mathbf{x}_{i}^{\prime} \mathbf{x}_{i}\right]$ for "small" $h_{N}$ is


\begin{equation*}
\left(2 h_{N}\right)^{-1} \mathrm{E}\left\{\left(1\left[\left|u_{i}\right| \leq h_{N}\right] \mid \mathbf{x}_{i}\right) \mathbf{x}_{i}^{\prime} \mathbf{x}_{i}\right\}=\left(2 h_{N}\right)^{-1} \mathrm{E}\left(1\left[\left|u_{i}\right| \leq h_{N}\right] \mathbf{x}_{i}^{\prime} \mathbf{x}_{i}\right), \tag{12.117}
\end{equation*}


where the equality holds by iterated expectations. Switching from the conditional to the unconditional expectation of $1\left[\left|u_{i}\right| \leq h_{N}\right]$ is an important simplification, and is reminiscent of the argument used to obtain the heteroskedasticity-robust variance matrix estimator. In that case, the matrix to be estimated, in the middle of the sandwich, is $\mathrm{E}\left[\mathrm{E}\left(u_{i}^{2} \mid \mathbf{x}_{i}\right) \mathbf{x}_{i}^{\prime} \mathbf{x}_{i}\right]=\mathrm{E}\left(u_{i}^{2} \mathbf{x}_{i}^{\prime} \mathbf{x}_{i}\right)$. The latter expression is much easier to estimate directly because it circumvents the need to estimate the conditional expectation $\mathrm{E}\left(u_{i}^{2} \mid \mathbf{x}_{i}\right)$.

Proceeding with the quantile regression case, we estimate (12.117) using the sample analogue, as usual:


\begin{equation*}
\hat{\mathbf{A}}=\left(2 h_{N}\right)^{-1} N^{-1} \sum_{i=1}^{N} 1\left[\left|\hat{u}_{i}\right| \leq h_{N}\right] \mathbf{x}_{i}^{\prime} \mathbf{x}_{i}=\left(2 N h_{N}\right)^{-1} \sum_{i=1}^{N} 1\left[\left|\hat{u}_{i}\right| \leq h_{N}\right] \mathbf{x}_{i}^{\prime} \mathbf{x}_{i}, \tag{12.118}
\end{equation*}


where $\hat{u}_{i}=y_{i}-\mathbf{x}_{i} \hat{\boldsymbol{\theta}}$ are the residuals from the quantile regression. Unfortunately, unlike in previous cases, we cannot so easily assert that $\hat{\mathbf{A}}$ is consistent for $\mathbf{A}_{\mathrm{o}}$ : we have said nothing about how quickly $h_{N}$ decreases to zero with the sample size, and the indicator function is not continuous. Powell (1991) (see also Koenker [2005]) shows that, under weak conditions, $h_{N} \rightarrow 0$ such that $\sqrt{N} h_{N} \rightarrow \infty$ are sufficient for consistency. The second condition controls how quickly $h_{N}$ shrinks to zero. For example, $h_{N}=a N^{-1 / 3}$ for any $a>0$ satisfies these conditions. The practical problem is choosing $a$ (or choosing $h_{N}$ more generally). Koenker (2005) contains specific recommendations, and also discusses related estimators. In equation (12.118), observation $i$ does not contribute if $\left|\hat{u}_{i}\right|>h_{N}$. Other methods allow each observation to enter the sum but with a weight that declines as $\left|\hat{u}_{i}\right|$ increases. In practice, one uses the most convenient estimate (that may be programmed into an existing econometrics package that performs quantile regression). The nonparametric bootstrap can be applied to quantile regression, but if the data set is large, the computation using several hundred bootstrap samples can be costly.

If we assume that $u_{i}$ is independent of $\mathbf{x}_{i}$ then $f_{u}\left(0 \mid \mathbf{x}_{i}\right)=f_{u}(0)$ and the asymptotic variance in equation (12.114) simplifies to


\begin{equation*}
\frac{\tau(1-\tau)}{\left[f_{u}(0)\right]^{2}}\left[\mathrm{E}\left(\mathbf{x}_{i}^{\prime} \mathbf{x}_{i}\right)\right]^{-1} \tag{12.119}
\end{equation*}


its estimator has the general form


\begin{equation*}
\frac{\tau(1-\tau)}{\left[\hat{f}_{u}(0)\right]^{2}}\left(N^{-1} \sum_{i=1}^{N} \mathbf{x}_{i}^{\prime} \mathbf{x}_{i}\right)^{-1} \tag{12.120}
\end{equation*}


and a simple, consistent estimate of $f_{u}(0)$ is


\begin{equation*}
\hat{f}_{u}(0)=\left(2 N h_{N}\right)^{-1} \sum_{i=1}^{N} 1\left[\left|\hat{u}_{i}\right| \leq h_{N}\right] \tag{12.121}
\end{equation*}


where $h_{N}$ satisfies the same conditions as before. The estimator in (12.121) is easily seen to be a simple histogram estimator of the density of $u_{i}$ at $u=0$, where the bin width is $2 h_{N}$ (and we must use the residuals in place of $u_{i}$ ). Alternatively, one can use a different kernel density estimator applied to $\left\{\hat{u}_{i}\right\}$; see, for example, Cameron and Trivedi (2005, Sect. 9.3).

Example 12.1 (Quantile Regression for Financial Wealth): We use the data set on single individuals $($ fsize $=1)$ in the data file 401KSUBS.RAW (from Abadie (2003))

\begin{table}[h]
\begin{center}
\captionsetup{labelformat=empty}
\caption{Table 12.1\\
Mean and Quantile Regression for Net Total Financial Wealth}
\begin{tabular}{|l|l|l|l|l|l|l|}
\hline
\multicolumn{7}{|l|}{Dependent} \\
\hline
Variable & \multicolumn{6}{|l|}{nettfa} \\
\hline
 & (1) & (2) & (3) & (4) & (5) & (6) \\
\hline
Explanatory &  &  &  & Median &  & . 90 \\
\hline
Variable & Mean (OLS) & . 10 Quantile & . 25 Quantile & (LAD) & . 75 Quantile & Quantile \\
\hline
\multirow[t]{2}{*}{inc} & . 783 & -. 0179 & . 0713 & . 324 & . 798 & 1.291 \\
\hline
 & (.104) & (.0177) & (.0072) & (.012) & (.025) & (.048) \\
\hline
\multirow[t]{2}{*}{age} & -1.568 & -. 0663 & . 0336 & -. 244 & -1.386 & -3.579 \\
\hline
 & (1.076) & (.2307) & (.0955) & (.146) & (.287) & (.501) \\
\hline
\multirow[t]{3}{*}{age ${ }^{2}$} & . 0284 & . 0024 & . 0004 & . 0048 & . 0242 & . 0605 \\
\hline
 & (.0138) & (.0027) & (.0011) &  & (.0017) & (.0034) \\
\hline
 & (.0059) &  &  &  &  &  \\
\hline
$e 401 k$ & 6.837 & . 949 & 1.281 & 2.598 & 4.460 & 6.001 \\
\hline
 & (2.173) & (.617) & (.263) & (.404) & (.801) & (1.437) \\
\hline
$N$ & 2,017 & 2,017 & 2,017 & 2,017 & 2,017 & 2,017 \\
\hline
\end{tabular}
\end{center}
\end{table}

The OLS standard errors (in parentheses) are robust to heteroskedasticity.\\
The quantile regression estimates, along with standard errors (in parentheses), were obtained using Stata 9.0. The variance matrix is of the form in equation (12.120)-that is, it assumes independence between the error and regressors.\\
to estimate conditional quantiles for net total financial wealth (nettfa). The explanatory variables are income, age, and a binary variable indicating eligibility to participate in a $401(\mathrm{k})$ pension plan through one's employer. The estimates, including OLS estimates of a linear model for the conditional mean, are given in Table 12.1. Financial wealth and income are both measured in thousands of dollars.

There are no surprises in the OLS estimates. The mean relationship between financial wealth and income is strong and very statistically significant. Further, eligibility in a $401(\mathrm{k})$ plan, holding income and age fixed, is estimated to increase expected financial wealth by about $\$ 6,800$. The heteroskedasticity-robust $t$ statistic is over three. The coefficients on age and age ${ }^{2}$ may appear puzzling at first, but they imply an increasing effect on nettfa starting at age $=27.6$. This makes sense because the youngest person in the sample is 25 . (In fact, the fit of the model is hardly changed-the $R$-squared goes from .1273 to .1272 -if we replace age and age $^{2}$ with the single explanatory variable $(\text { age }-25)^{2}$, and the restriction certainly cannot be rejected.)

The picture of the effects of income and $401(\mathrm{k})$ eligibility is very different when we look across the wealth distribution. How do we interpret the coefficient on inc for the quantile regressions? Consider the median regression result. Holding age and $e 401 k$ fixed, the coefficient on inc implies that if we compare two groups of people whose\\
income differs by $\$ 1,000$, the median financial wealth is estimated to be $\$ 324$ higher for the group with higher income. The effect on the median is less than half of the effect on the mean (\$783). The effect of income at the low end of the nettfa distribution, the .10 quantile, is nonexistent. Income has a large effect at the upper end of the financial wealth distribution. For example, increasing income by $\$ 1,000$ dollars increases the .90 quantile of nettfa by $\$ 1,291$, or more than $\$ 1,000$. The coefficient on inc is statistically greater than one.

The effects of $401(\mathrm{k})$ eligibility on nettfa also increase as we move up the wealth distribution, even conditional on income. Wealthy people can afford to contribute the maximums allowable by law, and so the option of contributing to tax-deferred savings plans, such as $401(\mathrm{k})$ plans, leads to a larger effect as we move up the distribution. (The coefficient on $e 401 k$ increases to about 9.7 at the .95 quantile.)

\subsection*{12.10.3 Quantile Regression for Panel Data}
Quantile regression methods can be applied to panel data, too. For a given quantile $0<\tau<1$, suppose we specify\\
$\operatorname{Quant}_{\tau}\left(y_{i t} \mid \mathbf{x}_{i t}\right)=\mathbf{x}_{i t} \boldsymbol{\theta}_{\mathrm{o}}, \quad t=1, \ldots, T$,\\
where $\mathbf{x}_{i t}$ probably allows for a full set of time period intercepts. Of course, we can write $y_{i t}=\mathbf{x}_{i t} \boldsymbol{\theta}_{\mathrm{o}}+u_{i t}$ where Quant $_{\tau}\left(u_{i t} \mid \mathbf{x}_{i t}\right)=0$. The natural estimator of $\boldsymbol{\theta}_{\mathrm{o}}$ is the pooled quantile regression estimator, $\hat{\boldsymbol{\theta}}$, which solves\\
$\min _{\theta \in \Theta} \sum_{i=1}^{N} \sum_{t=1}^{T} c_{\tau}\left(y_{i t}-\mathbf{x}_{i t} \boldsymbol{\theta}\right)$,\\
where $c_{\tau}(\cdot)$ is the check function. Now, when we define a score $\mathbf{s}_{i}(\boldsymbol{\theta})=\sum_{t=1}^{T} \mathbf{s}_{i t}(\boldsymbol{\theta})$, we generally have to account for serial correlation in $\mathbf{s}_{i t}\left(\boldsymbol{\theta}_{\mathrm{o}}\right)$ (although see Problem 12.15 for the case of dynamically complete quantiles). One technical issue in using the outer product of the score to estimate $\mathbf{B}_{\mathrm{o}}$, that is,\\
$\hat{\mathbf{B}}=N^{-1} \sum_{i=1}^{N} \mathbf{s}_{i}(\hat{\boldsymbol{\theta}}) \mathbf{s}_{i}(\hat{\boldsymbol{\theta}})^{\prime}=N^{-1} \sum_{i=1}^{N} \sum_{t=1}^{T} \sum_{r=1}^{T} \mathbf{s}_{i t}(\hat{\boldsymbol{\theta}}) \mathbf{s}_{i r}(\hat{\boldsymbol{\theta}})^{\prime}$,\\
is that the score function is discontinuous for $\boldsymbol{\theta}$ such that $y_{i t}=\mathbf{x}_{i t} \boldsymbol{\theta}$, for some $t \in\{1, \ldots, T\}$. Neverthless, if we assume $\mathrm{P}\left(u_{i 1} \neq 0, \ldots, u_{i T} \neq 0\right)=1$, then $\mathbf{s}_{i}(\boldsymbol{\theta})$ is continuous at $\boldsymbol{\theta}_{\mathrm{o}}$ with probability one, and this is enough to apply a slightly generalized version of Lemma 12.1. Because $\hat{\mathbf{B}}$ includes the terms $\mathbf{s}_{i t}(\hat{\boldsymbol{\theta}}) \mathbf{s}_{i r}(\hat{\boldsymbol{\theta}})^{\prime}$ for $t \neq r$, $\hat{\mathbf{B}}$ accounts for any kind of neglected dynamics in Quant $_{\tau}\left(y_{i t} \mid \mathbf{x}_{i t}\right)$. The terms\\
$\mathbf{s}_{i t}(\hat{\boldsymbol{\theta}}) \mathbf{s}_{i t}(\hat{\boldsymbol{\theta}})^{\prime}$ could be simplified along the lines of (12.116), but it is unecessary and seems pointless when the terms for $t \neq r$ are included.

Estimation of $\mathbf{A}_{\mathrm{o}}$ is similar to the cross section case. A fully robust estimator, one that does not assume independence between $u_{i t}$ and $\mathbf{x}_{i t}$ and allows the distribution of $u_{i t}$ to change across $t$, is an extension of equation (12.118):\\
$\hat{\mathbf{A}}=\left(2 N h_{N}\right)^{-1} \sum_{i=1}^{N} \sum_{t=1}^{T} 1\left[\left|\hat{u}_{i t}\right| \leq h_{N}\right] \mathbf{x}_{i t}^{\prime} \mathbf{x}_{i t}$,\\
or, we can replace the indicator function with a smoothed version. Rather than using $\hat{\mathbf{A}}^{-1} \hat{\mathbf{B}} \hat{\mathbf{A}}^{-1} / N$ as the estimate of $\operatorname{Avar}(\hat{\boldsymbol{\theta}})$, the bootstrap can be applied using the method for panel data described in Section 12.8.2.

Allowing explicitly for unobserved effects in quantile regression is trickier. For a given quantile $0<\tau<1$, a natural specification that incorporates strict exogeneity conditional on $c_{i}$ is\\
$\operatorname{Quant}_{\tau}\left(y_{i t} \mid \mathbf{x}_{i}, c_{i}\right)=\operatorname{Quant}_{\tau}\left(y_{i t} \mid \mathbf{x}_{i t}, c_{i}\right)=\mathbf{x}_{i t} \boldsymbol{\theta}_{\mathrm{o}}+c_{i}, \quad t=1, \ldots, T$,\\
which is reminiscent of the way we specified the conditional mean in Chapter 10. Equivalently, we can write\\
$y_{i t}=\mathbf{x}_{i t} \boldsymbol{\theta}_{\mathrm{o}}+c_{i}+u_{i t}, \quad$ Quant $_{\tau}\left(u_{i t} \mid \mathbf{x}_{i}, c_{i}\right)=0, \quad t=1, \ldots, T$.\\
Unfortunately, unlike in the case of estimating effects on the conditional mean, we cannot proceed without further assumptions. A "fixed effects" approach, where we allow $\mathrm{D}\left(c_{i} \mid \mathbf{x}_{i}\right)$ to be unrestricted, is attractive. Generally, there are no simple transformations to eliminate $c_{i}$ and estimate $\boldsymbol{\theta}_{\mathrm{o}}$. If we treat the $c_{i}$ as parameters to estimate along with $\boldsymbol{\theta}_{\mathrm{o}}$, the resulting estimator generally suffers from an incidental parameters problem, a topic that comes up in Chapter 13 and at several places in Part IV. Briefly, if we try to estimate $c_{i}$ for each $i$, then, with large $N$ and small $T$, the poor quality of the estimates of $c_{i}$ causes the accompanying estimate of $\boldsymbol{\theta}_{\mathrm{o}}$ to be badly behaved. Recall that this was not the case when we used the FE estimator for a conditional mean: treating the $c_{i}$ as parameters led us to the within estimator. Koenker (2004) derives asymptotic properties of this estimation procedure when $T$ grows along with $N$, but he adds the assumptions that the regressors are fixed and $\left\{u_{i t}: t=1, \ldots, T\right\}$ is serially independent.

An alternative approach is suggested by Abrevaya and Dahl (2008) for $T=2$. Motivated by Chamberlain's approach to linear unobserved effects models (see Section 11.1.2), Abrevaya and Dahl estimate separate quantile regressions Quant $_{\tau}\left(y_{i t} \mid \mathbf{x}_{i 1}, \mathbf{x}_{i 2}\right)$ (with intercepts, of course) for $t=1,2$. They define the partial\\
effects in a way that mimics a representation of the partial effects in Chamberlain's correlated random effects (CRE) approach.

For quantile regression, CRE approaches are generically hampered because finding quantiles of sums of random variables is difficult. For example, suppose we impose the Mundlak representation $c_{i}=\psi_{\mathrm{o}}+\overline{\mathbf{x}}_{i} \boldsymbol{\xi}_{\mathrm{o}}+a_{i}$. Then we can write $y_{i t}= \psi_{\mathrm{o}}+\mathbf{x}_{i t} \boldsymbol{\theta}_{\mathrm{o}}+\overline{\mathbf{x}}_{i} \boldsymbol{\xi}_{\mathrm{o}}+a_{i}+u_{i t} \equiv y_{i t}=\psi_{\mathrm{o}}+\mathbf{x}_{i t} \boldsymbol{\theta}_{\mathrm{o}}+\overline{\mathbf{x}}_{i} \boldsymbol{\xi}_{\mathrm{o}}+v_{i t}$, where $v_{i t}$ is the composite error. Now, if we assume $v_{i t}$ is independent of $\mathbf{x}_{i}$, then we can estimate $\boldsymbol{\theta}_{\mathrm{o}}$ and $\boldsymbol{\xi}_{\mathrm{o}}$ using pooled quantile regression of $y_{i t}$ on $1, \mathbf{x}_{i t}$, and $\overline{\mathbf{x}}_{i}$. (The intercept does not estimate a quantity of particular interest.) But independence is very strong, and, if we truly believe it, then we probably believe all quantile functions are parallel. Of course, we can always just assert that the effect of interest is the set of coefficients on $\mathbf{x}_{i t}$ in the pooled quantile estimation, and we allow these, along with the intercept and coefficients on $\overline{\mathbf{x}}_{i}$, to change across quantiles. The asymptotic variance matrix estimator discussed for pooled quantile regression applies directly once we define the explanatory variables at time $t$ to be $\left(1, \mathbf{x}_{i t}, \overline{\mathbf{x}}_{i}\right)$.

We have more flexibility if we are interested in the median, and a few simple approaches suggest themselves. Write the model $\operatorname{Med}\left(y_{i t} \mid \mathbf{x}_{i}, c_{i}\right)=\operatorname{Med}\left(y_{i t} \mid \mathbf{x}_{i t}, c_{i}\right)= \mathbf{x}_{i t} \boldsymbol{\theta}_{\mathrm{o}}+c_{i}$ in error form as\\
$y_{i t}=\mathbf{x}_{i t} \boldsymbol{\theta}_{\mathrm{o}}+c_{i}+u_{i t}, \quad \operatorname{Med}\left(u_{i t} \mid \mathbf{x}_{i}, c_{i}\right)=0, \quad t=1, \ldots, T$,\\
and consider the multivariate conditional distribution $\mathbf{D}\left(\mathbf{u}_{i} \mid \mathbf{x}_{i}\right)$. If this distribution is symmetric about zero in the sense that $\mathrm{D}\left(\mathbf{u}_{i} \mid \mathbf{x}_{i}\right)=\mathrm{D}\left(-\mathbf{u}_{i} \mid \mathbf{x}_{i}\right)$-which is sometimes called centrally symmetric-then the distribution of $\mathbf{g}^{\prime} \mathbf{u}_{i}$ given $\mathbf{x}_{i}$ is symmetric about zero for any linear combination $\mathbf{g}$ (see, for example, Serfling (2006) for discussion). In particular, the time-demeaned errors $\ddot{u}_{i t}$ have (univariate) conditional distributions symmetric about zero, which means we can consistently estimate $\boldsymbol{\theta}_{\mathrm{o}}$ by applying pooled least absolute deviations to the time-demeaned equation $\ddot{y}_{i t}=\ddot{\mathbf{x}}_{i t} \boldsymbol{\theta}_{\mathrm{o}}+\ddot{u}_{i t}$, being sure to obtain fully robust standard errors by using equations (12.124) and (12.125) on the time-demeaned data.

Alternatively, under the centrally symmetric assumption, the difference in the errors, $\Delta u_{i t}=u_{i t}-u_{i, t-1}$, have symmetric distributions about zero, so one can apply pooled LAD to $\Delta y_{i t}=\Delta \mathbf{x}_{i t} \boldsymbol{\theta}_{\mathrm{o}}+\Delta u_{i t}, t=2, \ldots, T$. From Honoré (1992) applied to the uncensored case, LAD on the first differences is consistent when $\left\{u_{i t}: t=1, \ldots, T\right\}$ is an i.i.d. sequence conditional on $\left(\mathbf{x}_{i}, c_{i}\right)$, even if the common distribution is not symmetric, and this may afford robustness for LAD on the first differences rather than on the time-demeaned data. (Interestingly, it follows from the discussion in Honoré (1992, Appendix 1) that when $T=2$, applying LAD on the first differences is equivalent to estimating the $c_{i}$ along with $\boldsymbol{\theta}_{\mathrm{o}}$. So, in this case, there is no incidental\\
parameters problem in estimating the $c_{i}$ as long as $u_{i 2}-u_{i 1}$ has a symmetric distribution.) Although not an especially weak assumption, central symmetry of $\mathrm{D}\left(\mathbf{u}_{i} \mid \mathbf{x}_{i}\right)$ allows for serial dependence and heteroskedasticity in the $u_{i t}$ (both of which can depend on $\mathbf{x}_{i}$ or on $t$ ). As always, we should be cautious in comparing the pooled OLS and pooled LAD estimates of $\boldsymbol{\theta}_{\mathrm{o}}$ on the demeaned or differenced data because they are only expected to be similar under the conditional symmetry assumption.

If we impose the Mundlak device, we can get by with conditional symmetry of a sequence of bivariate distributions. Write $y_{i t}=\psi_{\mathrm{o}}+\mathbf{x}_{i t} \boldsymbol{\theta}_{\mathrm{o}}+\overline{\mathbf{x}}_{i} \boldsymbol{\xi}_{\mathrm{o}}+a_{i}+u_{i t}$, where $\operatorname{Med}\left(u_{i t} \mid \mathbf{x}_{i}, a_{i}\right)=0$. If $\mathrm{D}\left(a_{i}, u_{i t} \mid \mathbf{x}_{i}\right)$ has a symmetric distribution around zero, then $\mathrm{D}\left(a_{i}+u_{i t} \mid \mathbf{x}_{i}\right)$ is symmetric about zero, and, if this holds for each $t$, pooled LAD of $y_{i t}$ on $1, \mathbf{x}_{i t}$, and $\overline{\mathbf{x}}_{i}$ consistently estimates $\left(\psi_{\mathrm{o}}, \boldsymbol{\theta}_{\mathrm{o}}, \boldsymbol{\xi}_{\mathrm{o}}\right)$. (Therefore, we can estimate the partial effects on $\operatorname{Med}\left(y_{i t} \mid \mathbf{x}_{i t}, c_{i}\right)$ and also test if $c_{i}$ is correlated with $\overline{\mathbf{x}}_{i}$.) The assumptions used for this approach are not as weak as we would like, but, as in using pooled LAD on the time-demeaned data, adding $\overline{\mathbf{x}}_{i}$ to pooled LAD gives a way to compare with the usual FE estimate of $\boldsymbol{\theta}_{\mathrm{o}}$. (Remember, if we use pooled OLS with $\overline{\mathbf{x}}_{i}$ included, we obtain the FE estimate.) Fully robust inference can be obtained by computing $\hat{\mathbf{B}}$ and $\hat{\mathbf{A}}$ in (12.124) and (12.125), respectively.

\section*{Problems}
12.1. a. Use equation (12.4) to show that $\boldsymbol{\theta}_{\mathrm{o}}$ minimizes $\mathrm{E}\left\{[y-m(\mathbf{x}, \boldsymbol{\theta})]^{2} \mid \mathbf{x}\right\}$ over $\Theta$ for any $\mathbf{x}$.\\
b. Explain why the result in part a is stronger than stating that $\boldsymbol{\theta}_{\mathrm{o}}$ solves problem (12.3).\\
12.2. Consider the model $\mathrm{E}(y \mid \mathbf{x})=m\left(\mathbf{x}, \boldsymbol{\theta}_{\mathrm{o}}\right), \operatorname{Var}(y \mid \mathbf{x})=\exp \left(\alpha_{\mathrm{o}}+\mathbf{x} \boldsymbol{\gamma}_{\mathrm{o}}\right)$, where $\mathbf{x}$ is $1 \times K$. The vector $\boldsymbol{\theta}_{\mathrm{o}}$ is $P \times 1$ and $\gamma_{\mathrm{o}}$ is $K \times 1$.\\
a. Define $u \equiv y-\mathrm{E}(y \mid \mathbf{x})$. Show that $\mathrm{E}\left(u^{2} \mid \mathbf{x}\right)=\exp \left(\alpha_{\mathrm{o}}+\mathbf{x} \boldsymbol{\gamma}_{\mathrm{o}}\right)$.\\
b. Let $\hat{u}_{i}$ denote the residuals from estimating the conditional mean by NLS. Argue that $\alpha_{\mathrm{o}}$ and $\gamma_{\mathrm{o}}$ can be consistently estimated by a nonlinear regression where $\hat{u}_{i}^{2}$ is the dependent variable and the regression function is $\exp \left(\alpha_{\mathrm{o}}+\mathbf{x} \boldsymbol{\gamma}_{\mathrm{o}}\right)$. (Hint: Use the results on two-step estimation.)\\
c. Using part b , propose a (feasible) weighted least squares procedure for estimating $\boldsymbol{\theta}_{\mathrm{o}}$.\\
d. If the error $u$ is divided by $[\operatorname{Var}(u \mid \mathbf{x})]^{1 / 2}$, we obtain $v \equiv \exp \left[-\left(\alpha_{\mathrm{o}}+\mathbf{x} \gamma_{\mathrm{o}}\right) / 2\right] u$. Argue that if $v$ is independent of $\mathbf{x}$, then $\boldsymbol{\gamma}_{\mathrm{o}}$ is consistently estimated from the regression $\log \left(\hat{u}_{i}^{2}\right)$ on $1, \mathbf{x}_{i}, i=1,2, \ldots, N$. (The intercept from this regression will\\
not consistently estimate $\alpha_{\mathrm{o}}$, but this fact does not matter, since $\exp \left(\alpha_{\mathrm{o}}+\mathbf{x} \boldsymbol{\gamma}_{\mathrm{o}}\right)= \sigma_{\mathrm{o}}^{2} \exp \left(\mathbf{x} \boldsymbol{\gamma}_{\mathrm{o}}\right)$, and $\sigma_{\mathrm{o}}^{2}$ can be estimated from the WNLS regression.)\\
e. What would you do after running WNLS if you suspect the variance function is misspecified?\\
12.3. Consider the exponential regression function $m(\mathbf{x}, \boldsymbol{\theta})=\exp (\mathbf{x} \boldsymbol{\theta})$, where $\mathbf{x}$ is $1 \times K$.\\
a. Suppose you have estimated a special case of the model, $\hat{\mathrm{E}}(y \mid \mathbf{z})=\exp \left[\hat{\theta}_{1}+\right. \left.\hat{\theta}_{2} \log \left(z_{1}\right)+\hat{\theta}_{3} z_{2}\right]$, where $z_{1}$ and $z_{2}$ are the conditioning variables. Show that $\hat{\theta}_{2}$ is approximately the elasticity of $\hat{\mathrm{E}}(y \mid \mathbf{z})$ with respect to $z_{1}$.\\
b. In the same estimated model from part a, how would you approximate the percentage change in $\hat{\mathrm{E}}(y \mid \mathbf{z})$ given $\Delta z_{2}=1$ ?\\
c. Now suppose a square of $z_{2}$ is added: $\hat{\mathrm{E}}(y \mid \mathbf{z})=\exp \left[\hat{\theta}_{1}+\hat{\theta}_{2} \log \left(z_{1}\right)+\hat{\theta}_{3} z_{2}+\right. \left.\hat{\theta}_{4} z_{2}^{2}\right]$, where $\hat{\theta}_{3}>0$ and $\hat{\theta}_{4}<0$. How would you compute the value of $z_{2}$ where the partial effect of $z_{2}$ on $\hat{\mathrm{E}}(y \mid \mathbf{z})$ becomes negative?\\
d. Now write the general model as $\exp (\mathbf{x} \boldsymbol{\theta})=\exp \left(\mathbf{x}_{1} \boldsymbol{\theta}_{1}+\mathbf{x}_{2} \boldsymbol{\theta}_{2}\right)$, where $\mathbf{x}_{1}$ is $1 \times K_{1}$ (and probably contains unity as an element) and $\mathbf{x}_{2}$ is $1 \times K_{2}$. Derive the usual (nonrobust) and heteroskedasticity-robust LM tests of $\mathrm{H}_{0}: \boldsymbol{\theta}_{\mathrm{o} 2}=\mathbf{0}$, where $\boldsymbol{\theta}_{\mathrm{o}}$ indexes $\mathrm{E}(y \mid \mathbf{x})$.\\
12.4. a. Show that the score for WNLS is $\mathbf{s}_{i}(\boldsymbol{\theta} ; \boldsymbol{\gamma})=-\nabla_{\theta} m\left(\mathbf{x}_{i}, \boldsymbol{\theta}\right)^{\prime} u_{i}(\boldsymbol{\theta}) / h\left(\mathbf{x}_{i}, \boldsymbol{\gamma}\right)$.\\
b. Show that, under Assumption WNLS.1, $\mathrm{E}\left[\mathbf{s}_{i}\left(\boldsymbol{\theta}_{\mathrm{o}} ; \boldsymbol{\gamma}\right) \mid \mathbf{x}_{i}\right]=\mathbf{0}$ for any value of $\boldsymbol{\gamma}$.\\
c. Show that, under Assumption WNLS.1, $\mathrm{E}\left[\nabla_{\gamma} \mathbf{s}_{i}\left(\boldsymbol{\theta}_{\mathrm{o}} ; \boldsymbol{\gamma}\right)\right]=\mathbf{0}$ for any value of $\boldsymbol{\gamma}$.\\
d. How would you estimate $\operatorname{Avar}(\hat{\boldsymbol{\theta}})$ without Assumption WNLS.3?\\
e. Verify that equation (12.59) is valid under Assumption WNLS.3.\\
12.5. a. For the regression model\\
$m(\mathbf{x}, \boldsymbol{\theta})=G\left[\mathbf{x} \boldsymbol{\beta}+\delta_{1}(\mathbf{x} \boldsymbol{\beta})^{2}+\delta_{2}(\mathbf{x} \boldsymbol{\beta})^{3}\right]$,\\
where $G(\cdot)$ is a known, twice continuously differentiable function with derivative $g(\cdot)$, derive the standard LM test of $\mathrm{H}_{0}: \delta_{\mathrm{o} 2}=0, \delta_{\mathrm{o} 3}=0$ using NLS. Show that, when $G(\cdot)$ is the identity function, the test reduces to RESET from Section 6.2.3.\\
b. Explain how to implement the variable addition version of the test.\\
12.6. Consider a panel data model for a random draw $i$ from the population:\\
$y_{i t}=m\left(\mathbf{x}_{i t}, \boldsymbol{\theta}_{\mathrm{o}}\right)+u_{i t}, \quad \mathrm{E}\left(u_{i t} \mid \mathbf{x}_{i t}\right)=0, \quad t=1, \ldots, T$.\\
a. If you apply pooled nonlinear least squares to estimate $\boldsymbol{\theta}_{\mathrm{O}}$, how would you estimate its asymptotic variance without further assumptions?\\
b. Suppose that the model is dynamically complete in the conditional mean, so that $\mathrm{E}\left(u_{i t} \mid \mathbf{x}_{i t}, u_{i, t-1}, \mathbf{x}_{i, t-1}, \ldots\right)=0$ for all $t$. In addition, $\mathrm{E}\left(u_{i t}^{2} \mid \mathbf{x}_{i t}\right)=\sigma_{\mathrm{o}}^{2}$. Show that the usual statistics from a pooled NLS regression are valid. (Hint: The objective function for each $i$ is $q_{i}(\boldsymbol{\theta})=\sum_{t=1}^{T}\left[y_{i t}-m\left(\mathbf{x}_{i t}, \boldsymbol{\theta}\right)\right]^{2} / 2$ and the score is $\mathbf{s}_{i}(\boldsymbol{\theta})= -\sum_{t=1}^{T} \nabla_{\theta} m\left(\mathbf{x}_{i t}, \boldsymbol{\theta}\right)^{\prime} u_{i t}(\boldsymbol{\theta})$. Now show that $\mathbf{B}_{\mathrm{o}}=\sigma_{\mathrm{o}}^{2} \mathbf{A}_{\mathrm{o}}$ and that $\sigma_{o}^{2}$ is consistently estimated by $(N T-P)^{-1} \sum_{i=1}^{N} \sum_{t=1}^{T} \hat{u}_{i t}^{2}$.)\\
c. For the mean model $m\left(\mathbf{x}_{i t}, \boldsymbol{\beta}, \boldsymbol{\delta}\right)$, consider testing $\mathrm{H}_{0}: \boldsymbol{\delta}_{\mathrm{o}}=\overline{\boldsymbol{\delta}}$. Show that under dynamic completeness and homoskedasticity (under $\mathrm{H}_{0}$ ), a valid version of the LM statistic is obtained as $N T R_{u}^{2}$ where $R_{u}^{2}$ is the uncentered $R$-squared from the pooled OLS regresssion\\
$\tilde{u}_{i t}$ on $\nabla_{\beta} \tilde{m}_{i t}, \nabla_{\delta} \tilde{m}_{i t}, \quad t=1, \ldots, T ; i=1, \ldots, N$, where $\tilde{u}_{i t}=y_{i t}-m\left(\mathbf{x}_{i t}, \tilde{\boldsymbol{\beta}}, \overline{\boldsymbol{\delta}}\right)$ and the gradients are evaluated at $(\tilde{\boldsymbol{\beta}}, \overline{\boldsymbol{\delta}})$.\\
12.7. Consider a nonlinear analogue of the SUR system from Chapter 7:\\
$\mathrm{E}\left(y_{i g} \mid \mathbf{x}_{i}\right)=\mathrm{E}\left(y_{i g} \mid \mathbf{x}_{i g}\right)=m_{g}\left(\mathbf{x}_{i g}, \boldsymbol{\theta}_{\mathrm{o} g}\right), \quad g=1, \ldots, G$.\\
Thus, each $\boldsymbol{\theta}_{\text {o } g}$ can be estimated by NLS using only equation $g$; call these $\check{\boldsymbol{\theta}}_{g}$. Suppose also that $\operatorname{Var}\left(\mathbf{y}_{i} \mid \mathbf{x}_{i}\right)=\boldsymbol{\Omega}_{\mathrm{o}}$, where $\boldsymbol{\Omega}_{\mathrm{o}}$ is $G \times G$ and positive definite.\\
a. Explain how to consistently estimate $\boldsymbol{\Omega}_{\mathrm{o}}$ (as usual, with $G$ fixed and $N \rightarrow \infty$ ). Call this estimator $\hat{\boldsymbol{\Omega}}$.\\
b. Let $\hat{\boldsymbol{\theta}}$ be the nonlinear SUR estimator that solves\\
$\min _{\boldsymbol{\theta}} \sum_{i=1}^{N}\left[\mathbf{y}_{i}-\mathbf{m}\left(\mathbf{x}_{i}, \boldsymbol{\theta}\right)\right]^{\prime} \hat{\boldsymbol{\Omega}}^{-1}\left[\mathbf{y}_{i}-\mathbf{m}\left(\mathbf{x}_{i}, \boldsymbol{\theta}\right)\right] / 2$,\\
where $\mathbf{m}\left(\mathbf{x}_{i}, \boldsymbol{\theta}\right)$ is the $G \times 1$ vector of conditional mean functions and $\mathbf{y}_{i}$ is $G \times 1$. Show that\\
$\operatorname{Avar} \sqrt{N}\left(\hat{\boldsymbol{\theta}}-\boldsymbol{\theta}_{\mathrm{o}}\right)=\left\{\mathrm{E}\left[\nabla_{\theta} \mathbf{m}\left(\mathbf{x}_{i}, \boldsymbol{\theta}_{\mathrm{o}}\right)^{\prime} \boldsymbol{\Omega}_{\mathrm{o}}^{-1} \nabla_{\theta} \mathbf{m}\left(\mathbf{x}_{i}, \boldsymbol{\theta}_{\mathrm{o}}\right)\right]\right\}^{-1}$.\\
(Hint: Under standard regularity conditions, $N^{-1 / 2} \sum_{i=1}^{N} \nabla_{\theta} \mathbf{m}\left(\mathbf{x}_{i}, \boldsymbol{\theta}_{\mathrm{o}}\right)^{\prime} \hat{\boldsymbol{\Omega}}^{-1}\left[\mathbf{y}_{i}-\right. \left.\mathbf{m}\left(\mathbf{x}_{i}, \boldsymbol{\theta}_{\mathrm{o}}\right)\right]=N^{-1 / 2} \sum_{i=1}^{N} \nabla_{\theta} \mathbf{m}\left(\mathbf{x}_{i}, \boldsymbol{\theta}_{\mathrm{o}}\right)^{\prime} \boldsymbol{\Omega}_{\mathrm{o}}^{-1}\left[\mathbf{y}_{i}-\mathbf{m}\left(\mathbf{x}_{i}, \boldsymbol{\theta}_{\mathrm{o}}\right)\right]+\mathrm{o}_{p}(1)$.)\\
c. How would you estimate $\operatorname{Avar}(\hat{\boldsymbol{\theta}})$ ?\\
d. If $\boldsymbol{\Omega}_{\mathrm{o}}$ is diagonal and if the assumptions stated previously hold, show that NLS equation by equation is just as asymptotically efficient as the nonlinear SUR estimator.\\
e. Is there a nonlinear analogue of Theorem 7.7 for linear systems in the sense that nonlinear SUR and NLS equation by equation are asymptotically equivalent when the same explanatory variables appear in each equation? (Hint: When would $\nabla_{\theta} \mathbf{m}\left(\mathbf{x}_{i}, \boldsymbol{\theta}_{\mathrm{o}}\right)$ have the form needed to apply the hint in Problem 7.5? You might try $\mathrm{E}\left(y_{g} \mid \mathbf{x}\right)=\exp \left(\mathbf{x} \boldsymbol{\theta}_{\mathrm{o} g}\right)$ for all $g$ as an example.)\\
12.8. Consider the M -estimator with estimated nuisance parameter $\hat{\gamma}$, where $\sqrt{N}\left(\hat{\gamma}-\gamma_{\mathrm{o}}\right)=\mathrm{O}_{p}(1)$. If assumption (12.37) holds under the null hypothesis, show that the QLR statistic still has a limiting chi-square distribution, assuming also that $\mathbf{A}_{\mathrm{o}}=\mathbf{B}_{\mathrm{o}}$. (Hint: Start from equation (12.76) but where $\sqrt{N}(\tilde{\boldsymbol{\theta}}-\hat{\boldsymbol{\theta}})=\mathbf{A}_{\mathrm{o}}^{-1} N^{-1 / 2}$. $\sum_{i=1}^{N} \mathbf{s}_{i}(\tilde{\boldsymbol{\theta}} ; \hat{\boldsymbol{\gamma}})+\mathrm{o}_{p}(1)$. Now use a mean value expansion of the score about $\left(\tilde{\boldsymbol{\theta}}, \boldsymbol{\gamma}_{\mathrm{o}}\right)$ to show that $\sqrt{N}(\tilde{\boldsymbol{\theta}}-\hat{\boldsymbol{\theta}})=\mathbf{A}_{\mathrm{o}}^{-1} N^{-1 / 2} \sum_{i=1}^{N} \mathbf{s}_{i}\left(\tilde{\boldsymbol{\theta}} ; \boldsymbol{\gamma}_{\mathrm{o}}\right)+\mathrm{o}_{p}(1)$.)\\
12.9. For scalar $y$, suppose that $y=m\left(\mathbf{x}, \boldsymbol{\beta}_{\mathrm{o}}\right)+u$, where $\mathbf{x}$ is a $1 \times K$ vector.\\
a. If $\mathrm{E}(u \mid \mathbf{x})=0$, what can you say about $\operatorname{Med}(y \mid \mathbf{x})$ ?\\
b. Suppose that $u$ and $\mathbf{x}$ are independent. Show that $\mathrm{E}(y \mid \mathbf{x})-\operatorname{Med}(y \mid \mathbf{x})$ does not depend on $\mathbf{x}$.\\
c. What does part b imply about $\partial \mathrm{E}(y \mid \mathbf{x}) / \partial x_{j}$ and $\partial \operatorname{Med}(y \mid \mathbf{x}) / \partial x_{j}$ ?\\
12.10. For each $i$, let $y_{i}$ be a nonnegative integer with a conditional binomial distribution with upper bound $n_{i}$ (a positive integer) and probability of success $p\left(\mathbf{x}_{i}, \boldsymbol{\beta}_{\mathrm{o}}\right)$, where $0<p(\mathbf{x}, \boldsymbol{\beta})<1$ for all $\mathbf{x}$ and $\boldsymbol{\beta}$. (A leading case is the logistic function.) Therefore, $\mathrm{E}\left(y_{i} \mid \mathbf{x}_{i}, n_{i}\right)=n_{i} p\left(\mathbf{x}_{i}, \boldsymbol{\beta}_{\mathrm{o}}\right)$ and $\operatorname{Var}\left(y_{i} \mid \mathbf{x}_{i}, n_{i}\right)=n_{i} p\left(\mathbf{x}_{i}, \boldsymbol{\beta}_{\mathrm{o}}\right)\left[1-p\left(\mathbf{x}_{i}, \boldsymbol{\beta}_{\mathrm{o}}\right)\right]$. Explain in detail how to obtain the weighted nonlinear least squares estimator of $\boldsymbol{\beta}_{\mathrm{o}}$.\\
12.11. a. Derive equation (12.98) for the WMNLS estimator. You can use equation (12.97) to show that estimation of $\boldsymbol{\gamma}^{*}$ can be ignored in obtaining the limiting distribution of $\sqrt{N}\left(\hat{\boldsymbol{\theta}}-\boldsymbol{\theta}_{\mathrm{O}}\right)$.\\
b. If assumption (12.96) holds (and $\hat{\gamma}$ is $\sqrt{N}$-consistent for $\gamma_{\mathrm{o}}$ ), what is the asymptotic variance of $\sqrt{N}\left(\hat{\boldsymbol{\theta}}-\boldsymbol{\theta}_{\mathrm{o}}\right)$ ?\\
c. Explain how you would compute the QLR statistic for testing hypotheses about $\boldsymbol{\theta}_{\mathrm{o}}$ under assumption (12.96). Naturally, you should assume the mean is correctly specified.\\
12.12. Let $y_{i}$ be a scalar response, and let $\mathbf{x}_{i}$ and $\mathbf{w}_{i}$ be vectors. Suppose that\\
$\mathrm{E}\left(y_{i} \mid \mathbf{x}_{i}, \mathbf{w}_{i}\right)=m\left(\mathbf{x}_{i}, \mathbf{v}\left(\mathbf{w}_{i}, \boldsymbol{\delta}_{\mathrm{o}}\right), \boldsymbol{\theta}_{\mathrm{o}}\right)$,\\
where $\mathbf{v}(\mathbf{w}, \boldsymbol{\delta})$ is a known function of $\mathbf{w}$ and the $J$ parameters $\boldsymbol{\delta}$ and $\boldsymbol{\theta}_{\mathrm{o}}$ are a $P \times 1$ vector. Assume that we have a $\sqrt{N}$-asymptotically normal estimator of $\boldsymbol{\delta}_{\mathrm{o}}$ that\\
satisfies\\
$\sqrt{N}\left(\hat{\boldsymbol{\delta}}-\boldsymbol{\delta}_{\mathrm{o}}\right)=N^{-1 / 2} \sum_{i=1}^{N} \mathbf{r}\left(\mathbf{w}_{i}, \boldsymbol{\delta}_{\mathrm{o}}\right)+o_{p}(1)$.\\
Let $\hat{\boldsymbol{\theta}}$ be the two-step NLS estimator that solves\\
$\min _{\theta \in \Theta} \sum_{i=1}^{N}\left[y_{i}-m\left(\mathbf{x}_{i}, \mathbf{v}\left(\mathbf{w}_{i}, \hat{\boldsymbol{\delta}}\right), \boldsymbol{\theta}\right)\right]^{2} / 2$.\\
a. Show that, under standard regularity conditions, the asymptotic variance of $\sqrt{N}\left(\hat{\boldsymbol{\theta}}-\theta_{\mathrm{o}}\right)$ is at least as large as the asymptotic variance if we knew $\boldsymbol{\delta}_{\mathrm{o}}$.\\
b. Propose a consistent estimator of $\operatorname{Avar} \sqrt{N}\left(\hat{\boldsymbol{\theta}}-\boldsymbol{\theta}_{\mathrm{O}}\right)$.\\
12.13. Let $\left\{\left(\mathbf{x}_{i t}, y_{i t}\right): t=1, \ldots, T\right\}$ be a panel data set, and assume that for some $\boldsymbol{\theta}_{\mathrm{o}} \in \boldsymbol{\Theta} \subset \mathbb{R}^{P}, \mathrm{E}\left(y_{i t} \mid \mathbf{x}_{i t}\right)=m\left(\mathbf{x}_{i t}, \boldsymbol{\theta}_{\mathrm{o}}\right), t=1, \ldots, T$. Further, let $h\left(\mathbf{x}_{i t}, \boldsymbol{\gamma}\right)$ be a model for $\operatorname{Var}\left(y_{i t} \mid \mathbf{x}_{i t}\right)$ for each $t$. Generally, let $\hat{\gamma}$ be a $\sqrt{N}$-consistent estimator for some $\boldsymbol{\gamma}^{*}$, where $h\left(\mathbf{x}_{i t}, \boldsymbol{\gamma}^{*}\right)$ need not equal $\operatorname{Var}\left(y_{i t} \mid \mathbf{x}_{i t}\right)$.\\
a. Let $\hat{\boldsymbol{\theta}}$ denote the pooled weighted nonlinear least squares estimator. Is strict exogeneity of $\left\{\mathbf{x}_{i t}\right\}$ needed for consistency of $\hat{\boldsymbol{\theta}}$ for $\boldsymbol{\theta}_{\mathrm{o}}$ ? Explain.\\
b. Propose a consistent estimator of $\operatorname{Avar} \sqrt{N}\left(\hat{\boldsymbol{\theta}}-\boldsymbol{\theta}_{\mathrm{o}}\right)$ without further assumptions.\\
c. If the mean is dynamically complete, how would you estimate Avar $\sqrt{N}\left(\hat{\boldsymbol{\theta}}-\boldsymbol{\theta}_{\mathrm{o}}\right)$ ?\\
d. If the mean is dynamically complete and $\operatorname{Var}\left(y_{i t} \mid \mathbf{x}_{i t}\right)=\sigma_{\mathrm{o}}^{2} h\left(\mathbf{x}_{i t}, \gamma_{\mathrm{o}}\right)$ and $\hat{\gamma}$ is $\sqrt{N}$ consistent for $\boldsymbol{\gamma}_{\mathrm{o}}$, how would you estimate $\operatorname{Avar} \sqrt{N}\left(\hat{\boldsymbol{\theta}}-\boldsymbol{\theta}_{\mathrm{o}}\right)$ ?\\
12.14. Derive equation (12.115). (Hint: The product of the two indicator functions appearing in $\mathbf{s}_{i}\left(\boldsymbol{\theta}_{\mathrm{O}}\right)$ is identically zero.)\\
12.15. Let $\left\{\left(\mathbf{x}_{i t}, y_{i t}\right): t=1, \ldots, T\right\}$ be a panel data set and assume that, for a given quantile $0<\tau<1$, Quant ${ }_{\tau}\left(y_{i t} \mid \mathbf{x}_{i t}\right)=\mathbf{x}_{i t} \boldsymbol{\theta}_{\mathrm{o}}, t=1, \ldots, T$. Let $\hat{\boldsymbol{\theta}}$ be the pooled quantile regression estimator discussed in Section 12.10.3.\\
a. Write down the approximate first-order condition solved by $\hat{\boldsymbol{\theta}}$. In particular, define a suitable "score" for each $t$ and then for a random draw $i$.\\
b. Show that if the quantile is dynamically complete in the sense that

Quant $_{\tau}\left(y_{i t} \mid \mathbf{x}_{i t}, y_{i, t-1}, \mathbf{x}_{i, t-1}, \ldots, y_{i 1} \mathbf{x}_{i 1}\right)=$ Quant $_{\tau}\left(y_{i t} \mid \mathbf{x}_{i t}\right)=\mathbf{x}_{i t} \boldsymbol{\theta}_{\mathrm{o}}, \quad t=1, \ldots, T$, then $\mathbf{B}_{\mathrm{o}} \equiv \mathrm{E}\left[\mathbf{s}_{i}\left(\boldsymbol{\theta}_{\mathrm{o}}\right) \mathbf{s}_{i}\left(\boldsymbol{\theta}_{\mathrm{o}}\right)^{\prime}\right]=\tau(1-\tau) \sum_{t=1}^{T} \mathrm{E}\left(\mathbf{x}_{i t}^{\prime} \mathbf{x}_{i t}\right)$. How would you estimate $\mathbf{B}_{\mathrm{o}}$ in this case?\\
c. Show that, whether or not the quantile is dynamically complete,\\
$\mathbf{A}_{\mathrm{o}}=\sum_{t=1}^{T} \mathrm{E}\left[f_{u_{t}}\left(0 \mid \mathbf{x}_{i t}\right) \mathbf{x}_{i t}^{\prime} \mathbf{x}_{i t}\right]$,\\
where $f_{u_{t}}\left(\cdot \mid \mathbf{x}_{t}\right)$ is the density of $u_{i t}$ given $\mathbf{x}_{i t}=\mathbf{x}_{t}$.\\
12.16. Consider a linear model with an endogenous explanatory variable, $y_{2}$, along with a reduced form for $y_{2}$ :\\
$y_{1}=\mathbf{z}_{1} \boldsymbol{\delta}_{1}+\alpha_{1} y_{2}+u_{1}$\\
$y_{2}=\mathbf{z} \boldsymbol{\pi}_{2}+v_{2}$,\\
where $\mathbf{z}=\left(\mathbf{z}_{1}, \mathbf{z}_{2}\right)$ (with first element of $\mathbf{z}_{1}$ unity) and $\boldsymbol{\pi}_{2}=\left(\boldsymbol{\pi}_{21}^{\prime}, \boldsymbol{\pi}_{22}^{\prime}\right)^{\prime}$; for notational simplicity, we do not use " $o$ " subscripts on the true parameters.\\
a. Under the assumption $\operatorname{Med}\left(v_{2} \mid \mathbf{z}\right)=0$, how would you estimate $\boldsymbol{\pi}_{2}$ ?\\
b. Suppose that $\operatorname{Med}\left(u_{1} \mid y_{2}, \mathbf{z}\right)=\operatorname{Med}\left(u_{1} \mid v_{2}\right)=\rho_{1} v_{2}$. Propose a two-step consistent estimator of $\boldsymbol{\delta}_{1}$ and $\alpha_{1}$.\\
c. How would you state the null hypothesis that $y_{2}$ is exogenous, and how would you test it?\\
d. Generally, how would you obtain legitimate standard errors for $\hat{\boldsymbol{\delta}}_{1}$ and $\hat{\alpha}_{1}$ ?\\
e. If the bivariate conditional distribution $\mathrm{D}\left(u_{1}, v_{2} \mid \mathbf{z}\right)$ satisfies the centrally symmetric condition $\mathrm{D}\left(u_{1}, v_{2} \mid \mathbf{z}\right)=\mathrm{D}\left(-u_{1},-v_{2} \mid \mathbf{z}\right)$ (and all variables have at least finite second moments), explain why we should expect the procedure from part $b$ and the usual 2SLS estimator to provide similar estimates in large samples.\\
12.17. Let $\hat{\boldsymbol{\theta}}$ be an M-estimator of $\boldsymbol{\theta}_{\mathrm{o}}$ with score $\mathbf{s}_{i}(\boldsymbol{\theta}) \equiv \mathbf{s}\left(\mathbf{w}_{i}, \boldsymbol{\theta}\right)$ and expected Hes$\operatorname{sian} \mathbf{A}_{\mathrm{o}}$ (evaluated at $\left.\boldsymbol{\theta}_{\mathrm{o}}\right)$. Let $\mathbf{g}(\mathbf{w}, \boldsymbol{\theta})$ be an $M \times 1$ function of the random vector $\mathbf{w}$ and the parameter vector, and suppose we wish to estimate $\boldsymbol{\delta}_{0} \equiv \mathrm{E}\left[\mathbf{g}\left(\mathbf{w}_{i}, \boldsymbol{\theta}_{\mathrm{o}}\right)\right]$. The natural estimator is $\hat{\boldsymbol{\delta}} \equiv N^{-1} \sum_{i=1}^{N} \mathbf{g}\left(\mathbf{w}_{i}, \hat{\boldsymbol{\theta}}\right)$. Assume that $\sqrt{N}\left(\hat{\boldsymbol{\theta}}-\boldsymbol{\theta}_{\mathrm{o}}\right) \xrightarrow{d} \operatorname{Normal}\left(\mathbf{0}, \mathbf{A}_{\mathrm{o}}^{-1} \mathbf{B}_{\mathrm{o}} \mathbf{A}_{\mathrm{o}}^{-1}\right)$ where $\mathbf{B}_{\mathrm{o}} \equiv \mathrm{E}\left[\mathbf{s}_{i}\left(\boldsymbol{\theta}_{\mathrm{o}}\right) \mathbf{s}_{i}\left(\boldsymbol{\theta}_{\mathrm{o}}\right)^{\prime}\right]$, as usual.\\
a. Assuming that $\mathbf{g}(\mathbf{w}, \cdot)$ is continuously differentiable on $\operatorname{int}(\boldsymbol{\Theta}), \boldsymbol{\theta}_{\mathrm{o}} \in \operatorname{int}(\boldsymbol{\Theta})$, and other regularity conditions, find $\operatorname{Avar} \sqrt{N}\left(\hat{\boldsymbol{\delta}}-\boldsymbol{\delta}_{\mathrm{o}}\right)$. (Hint: The asymptotic variance depends on $\boldsymbol{\delta}_{\mathrm{o}}$ and $\mathbf{G}_{\mathrm{o}} \equiv \mathrm{E}\left[\nabla_{\boldsymbol{\theta}} \mathbf{g}\left(\mathbf{w}_{i}, \boldsymbol{\theta}_{\mathrm{o}}\right)\right]$.)\\
b. How would you consistently estimate Avar $\sqrt{N}\left(\hat{\boldsymbol{\delta}}-\boldsymbol{\delta}_{\mathrm{o}}\right)$ ?\\
c. Show that if $\boldsymbol{g}(\mathbf{w}, \boldsymbol{\theta})=\mathbf{g}(\mathbf{x}, \boldsymbol{\theta})$, where $\mathbf{x}$ is exogenous in the estimation problem used to obtain $\hat{\boldsymbol{\theta}}$ (so that $\mathrm{E}\left[\mathbf{s}\left(\mathbf{w}_{i}, \boldsymbol{\theta}_{\mathrm{o}}\right) \mid \mathbf{x}_{i}\right]=\mathbf{0}$ ), then $\operatorname{Avar} \sqrt{N}\left(\hat{\boldsymbol{\delta}}-\boldsymbol{\delta}_{\mathrm{o}}\right)= \operatorname{Var}\left[\mathbf{g}\left(\mathbf{x}_{i}, \boldsymbol{\theta}_{\mathrm{o}}\right)\right]+\mathbf{G}_{\mathrm{o}}\left[\operatorname{Avar} \sqrt{N}\left(\hat{\boldsymbol{\theta}}-\boldsymbol{\theta}_{\mathrm{o}}\right)\right] \mathbf{G}_{\mathrm{o}}^{\prime}$.

\section*{13 Maximum Likelihood Methods}

\end{document}
